% PHP Gallery Fullständig produktdokumentation
% Project: PHP Gallery
% Repository: https://github.com/klusik/PHP_gallery
% File: docs/PHP_Gallery_Manual_SV.tex
% Module Type: LaTeX-produktmanual
% Purpose: Förklara administrationsflöden, drift och underhåll av PHP Gallery.
% Responsibilities:
%   - Underhåll användarguiden och den tekniska referensen med återskapbar PDF-typografi.
% Author: Rudolf Klusal
% Bygg från denna mapp med:
%   pdflatex PHP_Gallery_Manual_SV.tex
%   makeindex PHP_Gallery_Manual_SV.idx
%   pdflatex PHP_Gallery_Manual_SV.tex
%   pdflatex PHP_Gallery_Manual_SV.tex
%
% Markdown- och implementationsfilerna i referenserna är denna utgåvas
% redaktionella och tekniska källor. Granska det lokala projektet
% när programversionen ändras.

\documentclass[11pt,a4paper]{article}

\usepackage[utf8]{inputenc}
\usepackage[T1]{fontenc}
\usepackage{lmodern}
\usepackage[swedish]{babel}
\usepackage{geometry}
\usepackage{microtype}
\usepackage{xcolor}
\usepackage{graphicx}
\usepackage{booktabs}
\usepackage{array}
\usepackage{ragged2e}
\usepackage{longtable}
\usepackage{tabularx}
\usepackage{enumitem}
\usepackage{listings}
\usepackage{fancyhdr}
\usepackage{titlesec}
\usepackage[bottom,hang,flushmargin]{footmisc}
\usepackage{makeidx}
\usepackage{hyperref}
\usepackage{bookmark}

\geometry{top=23mm,bottom=25mm,left=24mm,right=24mm,headheight=15pt}
\setlength{\parindent}{0pt}
\setlength{\emergencystretch}{1.5em}
\setlength{\parskip}{0.65em}
\setcounter{tocdepth}{3}
\setcounter{secnumdepth}{3}
\makeatletter
% Håll tvåsiffriga underavsnittsnummer åtskilda från deras rubriker i innehållsförteckningen.
\renewcommand*\l@subsection{\@dottedtocline{2}{1.5em}{4.1em}}
\makeatother
\setlist{itemsep=0.3em,topsep=0.45em}
\renewcommand{\arraystretch}{1.2}

\definecolor{GalleryBlue}{HTML}{245A86}
\definecolor{GalleryLightBlue}{HTML}{EAF2F8}
\definecolor{GalleryDark}{HTML}{263238}
\definecolor{GalleryGray}{HTML}{F3F5F7}
\definecolor{GalleryCode}{HTML}{1F2933}
\definecolor{GalleryGreen}{HTML}{2E6B45}
\definecolor{GalleryRed}{HTML}{8A2F2F}

% Utgåvemetadata finns här så att en utgåveuppdatering kräver en enda versionsändring.
\newcommand{\version}{0.118.4}
\newcommand{\manualdate}{5~oktober~2026}

\hypersetup{
  colorlinks=true,
  linkcolor=GalleryBlue,
  citecolor=GalleryGreen,
  urlcolor=GalleryBlue,
  pdftitle={PHP Gallery Fullständig produktdokumentation},
  pdfauthor={Rudolf Klusal},
  pdfsubject={Referens för installation, administration, arkitektur, säkerhet, underhåll och utveckling},
  pdfkeywords={PHP Gallery, galleri-CMS, administration, arkitektur, dokumentation},
  pdfdisplaydoctitle=true,
  bookmarksopen=true,
  bookmarksnumbered=true
}

\pagestyle{fancy}
\fancyhf{}
\fancyhead[L]{\small PHP Gallery}
\fancyhead[R]{\small Fullständig produktdokumentation}
\fancyfoot[L]{\small Version \version}
\fancyfoot[C]{\thepage}
\fancyfoot[R]{\small \manualdate}
\renewcommand{\headrulewidth}{0.4pt}
\renewcommand{\footrulewidth}{0.2pt}

\titleformat{\section}{\Large\bfseries\color{GalleryBlue}}{\thesection}{0.75em}{}
\titleformat{\subsection}{\large\bfseries\color{GalleryDark}}{\thesubsection}{0.75em}{}
\titleformat{\subsubsection}{\normalsize\bfseries\color{GalleryDark}}{\thesubsubsection}{0.75em}{}
\pretocmd{\section}{\clearpage}{}{}

\lstdefinestyle{gallery}{
  basicstyle=\ttfamily\small\color{GalleryCode},
  backgroundcolor=\color{GalleryGray},
  frame=single,
  rulecolor=\color{GalleryLightBlue},
  breaklines=true,
  breakatwhitespace=false,
  columns=fullflexible,
  keepspaces=true,
  showstringspaces=false,
  tabsize=4,
  xleftmargin=0.8em,
  xrightmargin=0.8em,
  framexleftmargin=0.4em,
  aboveskip=0.8em,
  belowskip=0.8em
}
\lstset{style=gallery}

\newcommand{\product}{\textbf{PHP Gallery}}
\newcommand{\setting}[1]{\nolinkurl{#1}}
\newcommand{\filepath}[1]{\path{#1}}
\newcommand{\ui}[1]{\textbf{#1}}
\newcommand{\codeword}[1]{\nolinkurl{#1}}
\newcommand{\important}[1]{%
  \begin{center}\fcolorbox{GalleryBlue}{GalleryLightBlue}{\parbox{0.91\linewidth}{\textbf{Viktigt.} #1}}\end{center}}
\newcommand{\securitynote}[1]{%
  \begin{center}\fcolorbox{GalleryRed}{GalleryGray}{\parbox{0.91\linewidth}{\textbf{Säkerhetsinformation.} #1}}\end{center}}
\newcommand{\performancenote}[1]{%
  \begin{center}\fcolorbox{GalleryGreen}{GalleryGray}{\parbox{0.91\linewidth}{\textbf{Prestandainformation.} #1}}\end{center}}

\makeindex

\begin{document}

\pagenumbering{gobble}
\begin{titlepage}
  \centering
  \vspace*{24mm}
  {\Huge\bfseries\color{GalleryBlue} PHP Gallery\par}
  \vspace{6mm}
  {\LARGE Fullständig produktdokumentation\par}
  \vspace{14mm}
  {\large Referens för installation, administration, arkitektur, säkerhet, underhåll och utveckling\par}
  \vfill
  \begin{tabular}{rl}
    \textbf{Programversion:} & \version\\
    \textbf{Manualutgåva:} & \manualdate\\
    \textbf{Dokumentformat:} & A4, PDF med index och länkar\\
    \textbf{Projektförfattare:} & Rudolf Klusal\\
    \textbf{Licens:} & MIT-licensen\\
  \end{tabular}
  \vfill
  {\small Vägledning för att skapa, organisera, presentera, skydda, dela och bevara en fotografisamling med PHP Gallery version \version.\par}
\end{titlepage}

\pagenumbering{roman}
\begingroup
\small
\setlength{\parskip}{0pt}
\tableofcontents
\endgroup
\clearpage
\pagenumbering{arabic}

\section{Syfte och läsanvisning}
\index{d00065@dokumentation}
\index{m00157@manual!o00179@omfattning}

\product{} är ett fotogalleri som du driver på din egen server. Fotografierna ligger kvar på lagring som du kontrollerar, medan programmet tillför katalog, navigering, åtkomstregler, miniatyrbilder, kartor, sökning, hämtningar och administrationsverktyg. En enda installation kan användas för en offentlig portfolio, ett familjearkiv, kundleveranser, evenemangsbilder, resesamlingar eller ett stort personligt bibliotek.

Manualen är upplagd efter det arbete som en ägare faktiskt utför. De första avsnitten behandlar installationen och det första galleriet; mitten av manualen handlar om publicering, medier, integritet, presentation och underhåll. Arkitektur och utvecklarmaterial ligger mot slutet så att de inte avbryter den vanliga administrationen.

Börja med avsnitt~\ref{sec:first-hour} vid en ny installation. Planering av samlingen behandlas i avsnitt~\ref{sec:planning}; löpande redigering av gallerier och fotografier finns i avsnitten~\ref{sec:create-edit-gallery} och~\ref{sec:photo-workflow}. Avsnitt~\ref{sec:audiences} är användbart före publicering eftersom det jämför flera vanliga målgrupper. Innehållsförteckningen och indexet hjälper dig att gå direkt till en inställning eller ett arbetsflöde när webbplatsen redan används.

\subsection{Så använder du manualen}
\index{m00157@manual!a00011@användning}

Namn i administrationen, såsom \ui{Tema} och \ui{Systemstatus}, visas med fetstil. Text med fast teckenbredd betecknar filnamn, inställningsnamn, värden eller kommandon. Instruktioner skrivs som numrerade listor när ordningsföljden är viktig; checklistor och referensmaterial använder punkter eller tabeller. Tekniska kommentarer finns bara där de förklarar ett synligt beteende, ett beslut om återställning eller ett underhållskrav.

\subsection{Ytterligare projektreferenser}

Referenser för utvecklare och webbhotell hålls åtskilda från denna ägarmanual. Referensavsnittet anger de dokument i kodförrådet som används för arkitektur, databas, ansvar för källkod, testning och regler för bidrag \cite{readme,architecture,database,codemap,testing,agents}.\footnote{Dessa filer är bättre källor när du ändrar kod, undersöker schemadetaljer eller förbereder en utgåva.}

\subsection{Terminologi}

\begin{description}
  \item[Galleri] En namngiven samling fotografier. Ett galleri kan innehålla fotografier och mindre gallerier.
  \item[Fotografi eller bild] Ett objekt i ett galleri, tillsammans med dess titel, beskrivning, taggar, synlighet och tillgängliga kamerauppgifter.
  \item[Miniatyrbild] En mindre kopia som skapas för snabb visning i gallerirutnät, sökresultat, omslag och förhandsvisningar.
  \item[Besökare] En person som tittar på webbplatsens offentliga del.
  \item[Administratör] En betrodd person som loggar in för att skapa och hantera innehåll och inställningar.
  \item[Valt galleri] Det galleri som för tillfället är öppet på skärmen.
  \item[Gallerigren] Ett galleri tillsammans med alla mindre gallerier som ligger under det.
\end{description}

\section{Produktöversikt}
\label{sec:overview}
\index{p00192@produktöversikt}
\index{d00058@delat webbhotell}

\product{} gör en mappsamling till en fotowebbplats som går att bläddra i. Besökare ser omslag, titlar, beskrivningar, bildrutnät, helskärmsvisning, kartor, taggar, sökning och hämtningar enligt ägarens val. Administratörer får privata verktyg för att lägga till fotografier, ordna samlingar, styra integritet, förbättra beskrivningar och hålla installationen i gott skick.

Vanliga webbhotell med PHP/MySQL är den huvudsakliga miljön för driftsättning. Ägaren behåller originalfotografierna i igenkännbara mappar, medan databasen tillhandahåller katalogen och det programtillstånd som behövs för att bläddra.\footnote{Projektöversikten ger en kort sammanfattning av installationen och funktionerna \cite{readme}.}

\subsection{Grundläggande funktioner}

\begin{itemize}
  \item Skapa gallerier inuti andra gallerier, t.ex. \emph{2026 / Italien / Rom} eller \emph{Kunder / Northwind / Slutligt urval}.
  \item Ge varje galleri en titel, beskrivning, datumperiod, omslag, offentlig adress, integritetsinställning och eget utseende.
  \item Ladda upp fotografier genom webbläsaren eller upptäck filer som har kopierats till gallerimappar.
  \item Lägg till titlar, berättelser, taggar, synlighet, ordning, datum och platsuppgifter till fotografier.
  \item Låt besökare söka, bläddra bland taggar, se fotografier i helskärm, följa kartor, rösta och hämta tillåtna samlingar.
  \item Skydda familjesamlingar, kundsamlingar eller känsliga samlingar med åtkomstkontroller och delningsmetoder.
  \item Hitta sannolika dubbletter och granska dem med direkta länkar till båda platserna.
  \item Håll installationen aktuell med vägledda uppdateringar, statuskontroller, säkerhetskopior och underhållsverktyg.
  \item Automatisera uppladdningar med avgränsad behörighet, utbyt gallerier med kompatibla installationer och förbered hämtning av valda fotografier eller hela gallerier.
  \item Generera valfria AI-metadata och översättningsutkast, och låt varje genererat värde granskas före publicering.
  \item Publicera robotregler och bildwebbplatskartor som sökmotorer kan läsa för det offentliga innehåll som besökare faktiskt kan komma åt.
\end{itemize}

\subsection{Samspelet mellan filsystem och databas}
\index{f00080@filsystembaserad utformning}
\index{d00049@databas}

PHP Gallery lagrar beständiga data både i filsystemet och i databasen. Gallerimapparna innehåller originalfotografierna i en igenkännbar trädstruktur. Databasen lagrar uppgifter som inte hör hemma i själva filen: titlar, beskrivningar, taggar, ordning, åtkomstinställningar, röster, datum, poster för genererade medier och annat programtillstånd.

Ingen av delarna är en fullständig säkerhetskopia på egen hand. Om bara databasen återställs saknar katalogposterna sina originalfiler; om bara galleriträdet återställs går filernas administrations- och presentationstillstånd förlorade. Säkerhetskopiering och återställning bör därför hantera databasen, gallerimapparna och den lokala konfigurationen som en helhet. Avsnitt~\ref{sec:user-data} återkommer till detta mer utförligt.

\section{Krav och miljö}
\label{sec:requirements}
\index{k00140@krav}
\index{p00184@PHP}
\index{m00169@MySQL}
\index{m00158@MariaDB}

Använd detta avsnitt när du väljer webbhotell eller kontrollerar ett befintligt konto. Kravtabellen lämpar sig också för att skickas direkt till ett webbhotell.

\subsection{Minimikrav för körning}

\begin{longtable}{>{\bfseries}>{\hspace{0pt}}p{0.27\linewidth}>{\hspace{0pt}}p{0.65\linewidth}}
\toprule
Komponent & Krav\\
\midrule
\endhead
PHP & Version 8.1 eller senare.\\
Databas & MySQL 5.7 eller senare, eller MariaDB 10.2 eller senare.\\
PDO & Tillägget PDO MySQL för programmets frågor och migreringar.\\
Bildbehandling & GD för vanlig generering av miniatyrbilder. Ytterligare lokala bildverktyg kan stödja specialformat.\\
Arkiv & ZipArchive för paket-, hämtnings- och överföringsflöden som skapar eller läser ZIP-data.\\
Filsystem & Skrivbara platser för programdata, gallerier, cache och genererade medier.\\
Webbserver & URL-omskrivning av Apache-typ rekommenderas för rena URL:er; dirigering med frågesträngar stöder kompatibla miljöer utan omskrivning.\\
\bottomrule
\end{longtable}

Utgående HTTPS-åtkomst stöder programuppdateringar och fjärrhämtade utgåvemetadata. Lokal drift och manuellt installerade utgåvor kan fungera med mer begränsad utgående trafik.\footnote{Nätverkspolicyn bör följa installationsägarens krav på uppdateringar och integritet. Uppdateringskoden validerar paket och bevarar återställningsdata enligt den lokala uppdaterarens implementation.}

\subsection{Behörigheter}
\index{b00022@behörigheter!f00079@filsystem}

Webbprocessen behöver läsåtkomst till programkoden och skrivåtkomst till utsedda föränderliga platser. Håll programkoden skrivskyddad där webbhotellet tillåter det, utom under en kontrollerad uppdatering. Originalfilerna i gallerierna behöver samma säkerhetskopieringsskydd som databasen eftersom hela programtillståndet omfattar båda delarna.

\important{Säkerhetskopiera databasen, gallerifilerna och installationskonfigurationen tillsammans. En säkerhetskopia av enbart databasen bevarar metadata men saknar det ursprungliga medieträdet.}

\section{Installation och grundinställningar}
\label{sec:installation}
\index{i00123@installation}
\index{g00118@grundinställningar}

\subsection{Installation i webbläsaren}

\subsubsection{Installationsordning}

Webbläsarinstallationen samlar in databasinställningar, initierar schemat genom migreringar, skapar den första administratören, förbereder skrivbara platser och skriver den lokala konfigurationen. En ny begäran omdirigeras till installationen när nödvändig konfiguration saknas.\footnote{Installationstillstånd och genererad konfiguration bör skyddas på samma sätt som inloggningsuppgifter. Källkodsarkiv ska fortsatt utesluta den lokala \filepath{config.php}.}

Vanliga steg:

\begin{enumerate}
  \item Skapa en tom MySQL- eller MariaDB-databas i webbhotellets kontrollpanel.
  \item Ladda upp programpaketet till den avsedda webbkatalogen.
  \item Öppna webbplatsen och följ installationsguiden.
  \item Ange anslutningsuppgifter för databasen och inloggningsuppgifter för den första administratören.
  \item Bekräfta skrivbara mappar och kontroller av tillägg.
  \item Slutför installationen och gå in i administrationen.
\end{enumerate}

\subsection{Installation från kommandoraden}

\subsubsection{Lokala kommandon}

Kodförrådet innehåller vanliga PHP-skript för migrering och skapande av administratörer:

\begin{lstlisting}[language=bash]
php scripts/migrate.php
php scripts/create_admin.php <username> <password>
\end{lstlisting}

Kommandoradsinstallationen följer samma modell för beständigt schema som webbläsarinstallationen. Granska \filepath{config.example.php}, skapa den lokala konfigurationen, förbered databasen, kör migreringarna och skapa den första administratören.

\subsection{Första genomgången i administrationen}

Granska följande efter installationen:

\begin{itemize}
  \item Webbplatsnamn, språk, tema, sidbredd och inställningar för offentlig bildvisning.
  \item Behörigheter för galleriroten och resultat från upptäckt.
  \item Miniatyrbildernas format, mått, kvalitet och underhållsstatus.
  \item Inställningar för åtkomst, lösenordsåterställning, e-post, telemetri, uppdateringskanal och säkerhetskopiering.
  \item Kompatibilitet för URL-omskrivning och kanoniska offentliga URL:er.
  \item Diagnostik i Systemstatus och väntande migreringar.
\end{itemize}

\subsection{Rena URL:er och omskrivningskompatibilitet}
\index{u00259@URL-omskrivning}
\index{d00062@dirigering med frågesträngar}
\index{r00201@rena URL:er}

PHP Gallery stöder rena offentliga sökvägar när värdmiljön tillämpar projektets omskrivningsregler. Rena URL:er är ett val för presentation och dirigering, inte ett krav för galleriets datamodell. Rutter med frågesträngar är fortsatt kompatibilitetsvägen och fungerar när omskrivning är avstängd eller inte kan verifieras tillförlitligt.

Inställningen för URL-omskrivning styr vilket format PHP Gallery använder i genererade länkar. När omskrivning är avstängd använder hjälpfunktionerna adresser med frågesträngar för gallerier, fotografier, miniatyrbilder och andra programrutter, i stället för att förutsätta att webbservern översätter en ren sökväg. URL-namn, gallerimappar, databasrader och originalfotografier behöver inte döpas om när inställningen ändras.

Om rena sökvägar inte fungerar på ett nytt webbhotell, stäng först av URL-omskrivning i administrationen och kontrollera att frågesträngsformatet, t.ex. \codeword{?page=gallery}, fungerar. På Apache eller LiteSpeed kontrollerar du sedan att projektets \filepath{.htaccess} respekteras och att omskrivningsstöd finns. På andra webbservrar bör kompatibilitetsläget behållas om inte motsvarande dirigeringsregler har konfigurerats.

\section{Din första timme med PHP Gallery}
\label{sec:first-hour}
\index{f00094@första timmen}
\index{k00135@komma igång}
\index{g00113@galleriägare}
\index{a00003@administrationens översikt}

Håll den första sessionen liten: skapa ett galleri, lägg till några fotografier och öppna resultatet i en webbläsare där du inte är inloggad som administratör. Temaarbete och djupare organisering kan vänta tills du vet att grundflödet fungerar.

\subsection{Öppna administrationen}
\index{a00005@administratörsinloggning}

Välj \ui{Administratörsinloggning} i den offentliga sidans sidhuvud och ange kontot som skapades vid installationen. Översikten är utgångspunkten för den dagliga driften. Den sammanfattar gallerier, fotografier, systemstatus och tillgängliga underhållsåtgärder.\footnote{Översiktens exakta innehåll beror på aktiverade funktioner, genomförda migreringar och utvecklingsinställningar. En välfungerande installation kan ändå visa informationsmeddelanden som uppmanar en administratör att slutföra valfritt arbete.}

Ägna några minuter åt att hitta följande områden:

\begin{itemize}
  \item \ui{Gallerier}, där samlingar skapas, upptäcks, redigeras och ordnas;
  \item \ui{Tema}, där den offentliga webbplatsens utseende och standardinställningar för presentation väljs;
  \item \ui{Taggar}, där återanvändbara ämnen och kategorier underhålls;
  \item \ui{Systemstatus}, där migreringar, miniatyrbilder, lagring och konsekvens granskas;
  \item \ui{Konto}, där inloggningsuppgifter, återställningsadress och e-postleverans hanteras.
\end{itemize}

\subsection{Skapa ett första galleri}
\index{g00096@galleri!s00223@skapa det första galleriet}

Öppna galleriadministrationen och skapa ett galleri med en kort titel, t.ex. \emph{Sommarpromenad}. Granska den föreslagna offentliga sökvägen och mappnamnet och lägg sedan till en mening som beskriver samlingens innehåll.

Gör galleriet offentligt för denna första kontroll och låt avancerad profilering och åsidosättningar av ärvda inställningar behålla standardvärdena. Lägg till tre till tio fotografier. Det räcker för att verifiera ordning, miniatyrbildsgenerering, titlar och ljuslådan utan att göra övningen till en stor import.

\subsection{Förhandsgranska som besökare}
\index{a00009@anonym förhandsgranskning}
\index{b00025@besökarförhandsgranskning}

Använd länken till det offentliga galleriet eller kontrollen för besökarförhandsgranskning. För en realistisk kontroll öppnar du den offentliga adressen i ett privat webbläsarfönster utan administratörssession. Kontrollera att titel, beskrivning, fotografier och navigering visas som avsett. Öppna ett fotografi, gå framåt och bakåt i visaren och återvänd till galleriet.

Administratörsvyer kan innehålla redigeringskontroller och ytterligare resurser. Anonym förhandsgranskning ger den tydligaste bilden av upplevelsen för familj, kunder eller allmänheten.\footnote{Skyddade gallerier och åldersbegränsat innehåll behöver ett separat anonymt test eftersom en aktiv administratörssession redan har bredare åtkomst.}

\subsection{Slutför checklistan för första timmen}

\begin{enumerate}
  \item Ange webbplatsens offentliga namn och språk.
  \item Skapa ett galleri med titel och beskrivning.
  \item Ladda upp en liten grupp fotografier.
  \item Lägg till en meningsfull titel för minst ett fotografi.
  \item Byt ordning på två fotografier och kontrollera den nya ordningen offentligt.
  \item Öppna galleriet som anonym besökare.
  \item Kontrollera galleriet på en skärm i telefonstorlek.
  \item Kontrollera att Systemstatus inte rapporterar något brådskande migrerings- eller behörighetsproblem.
  \item Skapa en samordnad säkerhetskopia när grundinställningarna är klara.
\end{enumerate}

\section{Planera en gallerisamling}
\label{sec:planning}
\index{p00185@planera gallerier}
\index{s00211@samlingens utformning}
\index{g00100@gallerihierarki!p00186@planering}

En genomtänkt struktur hjälper besökarna att förstå samlingen och hjälper ägaren att underhålla den över tid. Börja med hur människor kommer att leta efter fotografier. Mappnamn och databasdetaljer kan följa den mänskliga organiseringen.

\subsection{Välj en struktur som känns igen}

Vanliga strukturer:

\begin{description}
  \item[Efter år och händelse] Ett år på den högsta nivån innehåller bröllop, resor, firanden och projekt. Det passar familjearkiv och historiska arkiv.
  \item[Efter plats] Länder innehåller regioner och städer. Det passar rese-, flyg-, landskaps- och dokumentärfotografi.
  \item[Efter kund eller projekt] Varje kund innehåller uppdrag och leveransgallerier. Det passar professionellt arbete med separata åtkomstregler.
  \item[Efter ämne] Djurliv, arkitektur, porträtt, flygplan och liknande ämnen bildar bestående kategorier. Taggar kan ge ett andra sätt att bläddra mellan dem.
  \item[Efter arbetsflöde] Inkommande, utvalda, levererade och arkiverade gallerier representerar olika steg. Synlighetsinställningar håller arbetsmaterial privat medan färdiga urval blir offentliga.
\end{description}

Använd gallerier för tydliga navigeringsgränser och taggar för samband som går över gränserna. Ett fotografi från Prag kan ligga i \emph{Resor / Tjeckien / Prag} och ha taggar som \emph{arkitektur}, \emph{natt} och \emph{spårvagn}.\footnote{Ett fotografi har en enda galleriplats i den vanliga innehållsmodellen, medan återanvändbara taggar ger flera innehållsmässiga kopplingar. Verktyg för kopiering och flyttning stöder avsiktliga ändringar av den fysiska organiseringen.}

\subsection{Bestäm hierarkins djup}
\index{g00100@gallerihierarki!d00063@djup}

Underordnade gallerier stöder djupa samlingar, men de flesta besökare bläddrar lättare när varje nivå har ett tydligt syfte. Två till fyra nivåer fungerar bra för många webbplatser. Ett djupare arkiv är fortfarande praktiskt när brödsmulor, meningsfulla titlar och konsekventa datum vägleder läsaren.

Skapa ett undergalleri när det behöver egen titel, beskrivning, omslag, åtkomstpolicy, datumperiod eller offentlig adress. Håll fotografierna tillsammans när en ny nivå bara skulle tillföra ett extra klick.

\subsection{Förbered titlar och beskrivningar}
\index{g00112@gallerititel}
\index{g00097@galleribeskrivning}

En gallerititel ska vara begriplig utanför administrationsgränssnittet. En beskrivning kan besvara tre frågor:

\begin{itemize}
  \item Vad visar samlingen?
  \item Var och när skapades den?
  \item Vilket sammanhang gör fotografierna meningsfulla?
\end{itemize}

Beskrivningar kan vara korta för vardagliga album och omfattande för dokumentärt arbete. Låt första meningen ge den viktigaste sammanfattningen eftersom besökare ofta skummar innan de läser hela texten.

\subsection{Planera integriteten före uppladdning}
\index{i00129@integritetsplanering}
\index{g00110@gallerisynlighet}

Välj om en samling ska vara offentlig, opublicerad under administrativ förberedelse eller skyddad för en känd målgrupp. Övergalleriets policy kan påverka undergallerier, så samla gallerier med liknande målgrupper där det är praktiskt.

Exempel:

\begin{itemize}
  \item en offentlig portfolio med noggrant utvalda fotografier;
  \item en lösenordsskyddad familjegren;
  \item ett opublicerat korrekturgalleri för en kund under förberedelsen;
  \item ett offentligt evenemangsgalleri där vissa enskilda fotografier markerats som privata;
  \item en åldersbegränsad gren för innehåll som kräver besökarens bekräftelse.
\end{itemize}
\section{Skapa och redigera gallerier}
\label{sec:create-edit-gallery}
\index{g00096@galleri!s00224@skapande}
\index{g00096@galleri!r00200@redigering}
\index{g00108@galleriredigerare}

\subsection{Skapa ett galleri från administrationen}

Använd kontrollen \ui{+} i ett offentligt galleri eller \ui{Skapa galleri här} i dess redigerare för att skapa ett undergalleri. Administrationspanelen till höger frågar bara efter den nya titeln; den känner redan till övergalleriet som kontrollen tillhör. Det nya galleriet börjar som \ui{Opublicerat}. Efter skapandet öppnar samma panel den fullständiga redigeraren för beskrivning, datum, språk, taggar, åtkomst, visning, medier och fotografier. Den offentliga sidan lämnas inte under detta utökade arbetsflöde.

Om du börjar på den fristående sidan för att skapa gallerier i administrationen anger du en titel och kan även ange SimBrief-uppgifter, beskrivning, källspråk och taggar innan du skickar formuläret. Det stängda avsnittet \ui{Fler alternativ} innehåller mappnamn, inledande synlighet, datumperiod, övergalleri, röstning, filnamnsvisning och märket för antal bilder. En kort sammanfattning under avsnittet visar aktuell synlighet och övergalleri. Det särskilda formuläret för att skapa och ladda upp finns kvar när fotografier ska laddas upp vid skapandet. Vid valideringsfel bevaras inskickade värden som inte är hemliga och informationen om övergalleriet så att de kan rättas.

När du har skrivit minst två tecken i titelfältet kan administrationen visa resten av titeln från ett matchande befintligt galleri som ett diskret förslag direkt i fältet. Nya titlar under det valda övergalleriet prioriteras; träffar på andra platser övervägs först när den begränsade sökningen bland syskongallerier är uttömd. Äldre titlar kan utelämnas, så avsaknad av förslag bevisar inte att en titel är unik. Formuläret begär högst åtta kandidater i stället för att läsa in hela titelkatalogen.

Med markören i slutet och utan markerad text trycker du på \codeword{Tab} utan modifieringstangent eller på \ui{Högerpil}, eller väljer förslagstexten, för att acceptera den. Modifierade kortkommandon och textinmatning med inmatningsmetod behåller sitt vanliga beteende. Det första trycket på \codeword{Escape} avvisar kompletteringen; ett senare tryck kan nå den omgivande panelen. Översatt hjälp och ett diskret tillgänglighetsmeddelande gör förslagen möjliga att upptäcka utan att ändra fältnamnet.

Unicode-matchning accepterar endast säkra gränser mellan hela tecken. Vissa förslag med tecken utanför ASCII utelämnas när servertillägg eller webbläsarstöd för segmentering saknas. En långsam, misslyckad eller inaktuell sökning hindrar aldrig vanlig inmatning. Om ett förslag accepteras ändras bara titelfältet: formuläret skickas inte, ett befintligt galleri döps inte om och dess åtkomstregler ändras inte. Samma beteende gäller efter att ett sidopanelformulär har ersatts. Se \filepath{docs/TITLE_COMPLETION.md} för matchningsbegränsningar och tekniska detaljer.

Efter skapandet med enbart namn är galleriet redan en riktig samling i galleriträdet, även innan fler redigerarfält har sparats. Den fysiska mappen kan ta emot fotografier genom webbläsaren, filsystemsöverföring, upptäckt eller ett annat uppladdningsflöde som stöds. Fyll i och spara redigerarfälten före publicering. Databasen lagrar beskrivande inställningar och presentationsinställningar; den vanliga åtkomstpolicyn avgör vad besökarna kan se.

Varje genererat formulär för skapande eller klassisk uppladdning har också en privat åtgärdsidentitet. Om anslutningen bryts efter att servern kan ha slutfört arbetet ska du försöka igen med samma formulär i stället för att öppna ett nytt: samma administratör, åtgärd och innehåll får det ursprungliga slutförda resultatet. Ett nyöppnat formulär representerar en ny avsiktlig åtgärd och kan fortfarande skapa en avsiktlig dubblett. Identiteten är inte ett lösenord och ersätter inte inloggning och CSRF-skydd.

\subsubsection{Ordning vid skapande}

\begin{enumerate}
  \item Öppna det avsedda övergalleriet och välj dess kontroll \ui{+}, eller skapa ett rotgalleri från administrationen.
  \item Ange en titel som är begriplig för besökare; acceptera eventuellt ett titelförslag.
  \item Skapa det opublicerade galleriet. Den fullständiga redigeraren öppnas i samma panel när JavaScript är tillgängligt.
  \item Lägg till beskrivning, datum eller datumperiod, källspråk, taggar och eventuella SimBrief-data under \ui{Identitet}.
  \item Granska \ui{Åtkomst}, \ui{Visning}, \ui{Medier} och galleriets mapp- och övergallerivärden under Avancerat där det är relevant.
  \item Spara galleriet, lägg till fotografier och välj ett omslag när lämpliga fotografier finns.
  \item Publicera när allt är klart och förhandsgranska sedan resultatet anonymt.
\end{enumerate}

\subsubsection{Sparade standardvärden och SimBrief-utkast}
\index{g00109@galleriskapande!s00229@sparade standardvärden}
\index{s00220@SimBrief!u00261@utkast före skapande}

Redigeraren kan komma ihåg aktuell SimBrief-identifierare och aktuellt källspråk för den inloggade administratören. Välj det närliggande alternativet \ui{Kom ihåg för framtida gallerier} bara för ett värde som du vill återanvända. Senare gallerier visar värdet förvalt med en liten anteckning om att det kommer från dina sparade standardvärden. Värdet går fortfarande att redigera; om Kom ihåg lämnas omarkerat bevaras det tidigare sparade valet. Standardvärden tillhör en administratör och publicerar inte inloggningsuppgifter eller kopierar åtkomstregler. Den vanliga migreringsåtgärden måste tillämpa \codeword{202609250001\_gallery\_creation\_preferences.php} innan Kom ihåg kan spara ett nytt standardvärde; galleriskapande fungerar även utan den valfria tabellen.

Det synliga SimBrief-fältet accepterar ett~pilot-ID eller pilotnamn. Enbart siffror behandlas som ett~pilot-ID, medan annan text behandlas som ett namn. \ui{Generera beskrivningsutkast} kan hämta senaste OFP för~ett~nytt eller~befintligt galleri och fylla i en redigerbar Markdown-beskrivning på~valt innehållsspråk. När gallerilokalisering är~aktiverad och översättningslagringen är~redo fyller utkastet även i~beskrivningsfälten för de språk som underhålls. Granska och redigera beskrivningarna före sparandet. Importen bevarar en privat referens i 30~minuter, begränsad till samma administratör och session; om identifieraren ändras blir referensen ogiltig. För ett befintligt galleri kopplas OFP och tillgängliga ruttkoordinater först efter att du sparat galleriet. Om en valfri bilaga misslyckas förblir galleriet sparat och panelen visar en varning. Importera igen när utkastet löper ut eller sessionen ändras.

\subsubsection{Arbeta i den kompakta redigeraren}
\index{g00108@galleriredigerare!s00219@sidopanel}

De vanliga arbetsområdena är \ui{Identitet}, \ui{API}, \ui{Åtkomst}, \ui{Visning} och \ui{Medier}; därefter följer bildverktyg och specialverktyg när de är tillgängliga. Panelen till höger håller \ui{Spara galleri} synlig längst ned medan inställningsformuläret rullas. Att byta flik sparar inte automatiskt. Använd knappen efter redigering och öppna sedan galleriet igen om du vill verifiera de lagrade värdena. Längre förklaringar finns bakom frågeteckenskontroller som kan nås med tangentbordet, och mer sällan använda val ligger under utfällbara \ui{Avancerat}-avsnitt. En direkt administrationsredigerare använder samma fält, medan en webbläsare utan JavaScript använder vanlig formulärinlämning.

\ui{Identitet} håller titel, beskrivning, datum, källspråk, SimBrief och taggar lätt tillgängliga. Dess avancerade avsnitt innehåller URL-namn, mapp, övergalleri, sorteringsordning och AI-galleritext när funktionen är~aktiverad. \ui{API} visar en kopierbar uppladdningsadress, nyckelns etikett och åtgärd samt en nyskapad nyckel endast en gång; befintliga nycklar visas bara när sådana finns. Omgenerering av AI-metadata, migrering och den globala API-hanteraren ligger under \ui{Avancerade API-verktyg}. \ui{Åtkomst} samlar synlighet, lösenord, delningslänkskontroller och NSFW-kryssrutan, med längre policyförklaringar under frågetecken. \ui{Visning} börjar med rutnätet och vanliga växlar för röstning och filnamn; EXIF/GPS, format för gallerikort, antalmarkering, ljuslådeläge, Bildspelet, miniatyrbildsgränser ligger under \ui{Avancerade visningsinställningar}; flygrutttext finns bredvid SimBrief under \ui{Identitet}.

\subsection{Redigera identitet och sammanhang}
\index{g00102@gallerimetadata}
\index{d00056@datumperiod}

Redigeraren låter dig förbättra ett galleri över tid. Titlar och beskrivningar kan ändras när en samling utvecklas. Datumfält kan representera en enskild dag, en resa, en säsong eller en längre historisk period. EXIF-datumförslag kan undersöka fotografier och undergallerier och sedan erbjuda en redigerbar period för godkännande.

Den fullständiga redigeraren registrerar exakt vilken gallerirevision dess fält hämtades från. Om en annan flik, administratör, uppladdning, flytt, skanning eller underhållsåtgärd ändrar galleriet först nekas en äldre sparning i stället för att tyst skriva över det nyare tillståndet. Konfliktvyn bevarar inskriven text som inte är hemlig och visar en säker jämförelse med aktuellt innehåll. Granska och kopiera de avsedda ändringarna till den senaste redigeraren; ange lösenord igen och välj filer på nytt eftersom privata värden aldrig kopieras till konfliktinformationen.

Använd datumförslag som en hjälp att komma igång. Skannade fotografier, skärmbilder, redigerade exporter och kameror med felställd klocka kan ge datum som behöver mänsklig granskning.\footnote{Att tillämpa ett EXIF-baserat förslag uppdaterar de föreslagna gallerifälten genom det befintliga redigeringsflödet. Administratören ansvarar fortfarande för de slutliga datum som visas för besökarna.}

\subsubsection{Granska datumförslag för flera gallerier}
\index{g00098@galleridatum!e00073@EXIF-granskning}

Den särskilda underhållssidan \ui{Galleridatum} tillämpar samma EXIF-underlag på många gallerier. Den kan granska hela installationen eller en vald gallerigren. Varje rad visar aktuell manuellt angiven datumperiod, föreslagen period, antalet EXIF-fotografier och gallerinoder som bidrog med underlag samt redigerbara \ui{Från}/\ui{Till}-fält. Rader utan befintligt manuellt datum är valda för tillämpning som standard; rader som redan har ett manuellt datum lämnas omarkerade så att en massgranskning inte tyst skriver över ägarens tidsuppgifter.

Endast markerade rader sparas. Omarkerade förslag ignoreras, och administratören kan redigera båda datumen innan urvalet tillämpas. Om äldre importerade fotografier ännu saknar skannade EXIF-tagningsdatum kör du först den vanliga bildskanningen eller importen för galleriet. Samma riktade förslagsåtgärd finns i en enskild galleriredigerare och kan uppdatera redigeraren på plats.

Datumfält i administrationen behåller webbläsarens inbyggda datumfält som inskickat värde men lägger till en kompakt kalenderöppnare samt hjälpknapparna \ui{Idag} och \ui{Radera}. \ui{Idag} skriver webbläsarens aktuella lokala kalenderdatum; \ui{Radera} tömmer fältet. Förbättringen tillämpas igen på formulär som läses in i administrationens sidopanel, medan det underliggande inbyggda datumfältet fortfarande går att använda utan de extra JavaScript-kontrollerna.

\subsection{Referens för galleriredigerarens kontroller}
\index{g00108@galleriredigerare!k00137@kontroller}
\index{g00096@galleri!s00235@synlighet}
\index{g00096@galleri!z00276@åsidosättningar av visning}

Galleriredigeraren är den styrande platsen för inställningar som tillhör ett fysiskt galleri. De vanliga flikarna är Identitet, API, Åtkomst, Visning och Medier, följda av Bilder och eventuella verktyg för Organiserare/Filnamnsbyte där de är tillgängliga. Avancerade utfällbara avsnitt bevarar sällan använda inställningar i det gemensamma sparformuläret. Följande tabell sammanfattar aktuella galleriägda kontroller i version~\version{}.

\begin{longtable}{>{\RaggedRight\arraybackslash}>{\hspace{0pt}}p{0.21\textwidth} >{\RaggedRight\arraybackslash}>{\hspace{0pt}}p{0.27\textwidth} >{\RaggedRight\arraybackslash}>{\hspace{0pt}}p{0.43\textwidth}}
\toprule
\textbf{Område} & \textbf{Kontroll} & \textbf{Aktuellt beteende}\\
\midrule
\endfirsthead
\toprule
\textbf{Område} & \textbf{Kontroll} & \textbf{Aktuellt beteende}\\
\midrule
\endhead
Identitet & Titel och beskrivning & Offentlig redaktionell text. Underhållna översättningar kan lagras separat utan att ändra källtexten.\\
Identitet & Manuellt datum och datumperiod & Valfritt händelse- eller fotograferingsdatum samt stödd start- och slutperiod. EXIF-baserade förslag går att redigera före sparandet.\\
Identitet & URL-namn & Styr galleriets identitet i den offentliga URL:en. Behåll etablerade URL-namn där det går eftersom bokmärken och externa länkar kan vara beroende av dem.\\
Identitet & Mappnamn & Byter namn på den fysiska gallerimappen på disken genom den vanliga ändringsvägen. Det är inte bara en visningsetikett.\\
Identitet & Övergalleri och sorteringsordning & Ändrar hierarki och syskonordning. När ett galleri flyttas flyttas dess mappträd fysiskt där den konfigurerade galleriroten tillåter det.\\
Identitet & Taggar & Återanvändbara taggar åtskilda med kommatecken. Förslag viktas efter närliggande gallerier, fotografier och mappsammanhang i stället för att vara en osorterad slumpmässig lista.\\
Identitet & SimBrief-identifierare och standardvärden & Ett synligt fält accepterar numeriskt pilot-ID eller pilotnamn. Administratören kan spara identifieraren och källspråket för framtida gallerier; förifyllda värden är markerade.\\
API & Uppladdningsadress och nycklar & Kopiera den galleribegränsade adressen, generera en nyckel med etikett och kopiera det nya råvärdet direkt. Aktiva nycklar visas bara när sådana finns; råvärdet visas inte igen. AI-omgenerering och migrering finns under Avancerat.\\
Åtkomst & Synlighet & \ui{Offentligt} listas normalt. \ui{Opublicerat} döljs från vanliga listor men kan fortfarande öppnas från sin vanliga URL när åtkomstreglerna tillåter det. \ui{Privat} är endast för administratörer, förutom genom direkt tokenåtkomst som stöds.\\
Åtkomst & Lösenordslås & Oberoende av synlighet. Ett offentligt, opublicerat eller privat galleri kan ha en egen lösenordspolicy när åtkomstschemat är tillgängligt. Ett tomt fält för nytt lösenord bevarar det befintliga lösenordet.\\
Åtkomst & Delningslänk & Kan genereras, genereras om, återkallas och eventuellt ges en utgångstid. Lagrad tokeninformation skyddas med hash; en äldre token eller en token som inte kan dekrypteras kan behöva genereras om innan en kopierbar länk kan visas igen.\\
Åtkomst & NSFW / 18+ & Markerar gallerigrenen för anonym åldersbekräftelse. Skyddet gäller även undergallerier och levererade medier enligt den effektiva åtkomstpolicyn.\\
Visning & Bildspelet & Avancerat alternativ för jämförelsespelet i denna gallerigren. För att aktivera spelet måste även bildröstning vara tillgänglig.\\
Visning & Bildröstning & Aktiverar offentliga gilla-markeringar för bilder i galleriet. Om den stängs av bevaras lagrade poäng, men nya röstningskontroller tas bort och nya röster blockeras.\\
Visning & Visa filnamn & Avstängt som standard. Ursprungliga uppladdade filnamn döljs om detta inte aktiveras; egna bildtitlar och beskrivningar visas fortfarande.\\
Visning & EXIF/GPS-visning & \ui{Använd global standard}, \ui{Tvinga på} eller \ui{Tvinga av}. Det effektiva tillståndet styr offentliga GPS-markörer, EXIF-kartor för gallerier och visning av GPS-koordinater.\\
Visning & Beskrivningsformat för gallerikort & Ärver Tema som standard och kan åsidosätta kortlayouten för detta galleri där det stöds.\\
Visning & Märke för antal bilder & Ärver Tema eller åsidosätter ikonen med staplade bilder och antalet ingående bilder på gallerikort.\\
Visning & Bläddringsläge i ljuslådan & Ärver Tema, klassisk enskild bild, bildremsa eller 3D-karusell. Helskärm och bildspel är fortsatt tillgängliga i alla lägen.\\
Visning & Eget rutnät & Första synliga kontrollen under Visning: valfria kolumner och rader för detta galleri. Utan ett uttryckligt rutnät används närmaste övergalleris rutnät som tillåter arv till undergallerier, därefter Tema.\\
Visning & Gränser för responsiva miniatyrbilder & Valfria minimi- och maximigränser för miniatyrbildsstorlek. Vid sparande kan galleriets gränser kopieras rekursivt till undergallerier; rekursivt sparande är avstängt som standard.\\
Visning & Flygrutt & Avancerad manuell rutttext stöder lokal sökning efter identifierare och uttryckliga \codeword{NAME@latitude,longitude}-punkter. SimBrief kan spara en OFP, ett beskrivningsutkast och upplösta ruttkoordinater.\\
Medier & Titelbild & Automatisk eller ett valt fotografi, inklusive lämpliga fotografier från undergallerier när väljaren erbjuder dem.\\
Medier & Uppladdad galleriminiatyrbild & Separat medieresurs för gallerikort, lagrad oberoende av vanliga gallerifotografier.\\
Medier & Bakgrundskälla & Ärver Tema, laddar upp en ny bakgrund, väljer en befintlig galleribild eller genererar ett kollage från offentliga gallerier.\\
Medier & Profileringsresurser & Valfria gallerispecifika banderoll-, logotyp- och avdelarresurser åsidosätter Temas reservvärden endast för det galleriet där de har konfigurerats.\\
Bilder & Skanna/importera & Skannar om den fysiska gallerimappen och synkroniserar medier som stöds med katalogen.\\
Bilder & Skapa alla miniatyrbilder & Återskapar miniatyrbilder för detta galleri. Välj \ui{Inkludera undergallerier} för att omfatta dess undergallerier; andra gallerigrenar ingår inte i jobbet. Webbläsargenerering väljs som standard när den är tillgänglig, med servergenerering som alternativ.\\
Verktyg & Organiserare, Filnamnsbyte, API & Datumorganiserare, fysiskt mediefilnamnsbyte, hantering av uppladdnings-API-nycklar, AI-omgenerering och gallerimigrering visas bara när deras scheman och funktioner är tillgängliga.\\
\bottomrule
\end{longtable}

\important{Gallerispecifika kontroller samlas avsiktligt inte i det globala inställningsnavet. Arv är en del av datamodellen, så ett globalt sökresultat kan länka till galleriredigeraren utan att göra det galleriägda värdet till en webbplatsövergripande inställning.}

\subsection{Välj omslag och visuell identitet}
\index{g00104@galleriomslag}
\index{g00106@galleriprofilering}

Ett omslag är galleriets visuella inbjudan. Välj en bild som känns igen som miniatyrbild och representerar samlingen rättvist. Ansikten, tydliga silhuetter, klara motiv och rena kompositioner fungerar ofta bra i liten storlek.

Galleriprofilering kan tillföra logotyp, banderoll, bakgrund, avdelare eller en presentationsåsidosättning enligt tillgängliga kontroller. Börja med det ärvda temat och lägg sedan till gallerispecifika resurser när samlingen har en egen identitet, t.ex. ett bröllop, en utställning, en kund eller ett långvarigt projekt.

\subsection{Flytta och ordna om gallerier}
\index{g00096@galleri!f00082@flyttning}
\index{g00096@galleri!z00279@ändra ordning}

Att flytta ett galleri ändrar dess plats i hierarkin och kan påverka ärvd åtkomst eller presentation. Granska destinationen innan du bekräftar flytten. Att ändra ordning ändrar följden bland syskongallerier och stöder en genomtänkt läsordning.

Flyttning av fotografier mellan gallerier använder en beständig åtgärdsjournal. Organiseraren flyttar varje original fysiskt och verifierar dess registrerade identitet innan katalogtillhörigheten ändras. Om en arbetare eller anslutning stannar halvvägs visar Systemstatus och Körningsdiagnostik ett avgränsat åtgärdstillstånd utan att avslöja privata sökvägar. En misslyckad flytt visar en säker förklaring på skärmen; betrodda administrationsloggar behåller mer detaljerad privat diagnostik. Stoppa konkurrerande lagringsarbete, bevara en säkerhetskopia och följ \filepath{docs/IMAGE_MOVE_RECOVERY.md}; återställningen behåller ändrade eller icke verifierbara filer i stället för att skriva över dem.

Ordningen kan vara kronologisk, geografisk, berättande, alfabetisk eller manuellt utvald. Inom ett övergalleri bör valet vara tillräckligt konsekvent för att besökare ska kunna förutse nästa objekt.

En inloggad administratör kan också ordna om direkta undergallerikort på den offentliga gallerisidan. Det kompakta offentliga verktygsfältet för ordning omfattar bara de kort som syns på aktuell pagineringssida. Sparbegäran innehåller sidans förskjutning och antal synliga kort, och servern vägrar ändringen om inskickade ID:n inte längre motsvarar samma aktuella delmängd. Den ändrar bara syskonens \codeword{sort_order}-värden; den ändrar inte övergallerier eller nästlar gallerier.

När en administratör är inloggad på ett offentligt galleri kan direkta undergallerier med ifyllt \ui{Från}-datum också sorteras i förhandsvisningen från äldst till nyast eller från nyast till äldst. Undergallerier utan användbart datum behåller sina befintliga positioner. Förhandsvisningen är tillfällig tills \ui{Spara} väljs; sparandet skriver resultatet som den riktiga syskonordningen som alla besökare ser. Minst två daterade direkta undergallerier krävs för att datumordningen ska kunna sparas.

\subsection{Radera ett galleri med eftertanke}
\index{g00096@galleri!r00198@radering}

Galleriradering kan påverka originalfiler, undergallerier, metadata, miniatyrbilder, länkar och besökarnas bokmärken. Granska exakt vilket galleri och vilka undergallerier som är valda, skapa en säkerhetskopia och använd det vanliga raderingsflödet i administrationen. En tillfällig synlighetsändring kan ge en säkrare förberedelseperiod när galleriet fortfarande kan behövas.

Galleriets papperskorg är~aktiverad som standard. När ett galleri raderas genom en vanlig administratörsåtgärd flyttas den valda mappen och hela dess underträd bort från den aktiva galleriroten, de metadata som krävs för återställning registreras och posten listas under \ui{Underhåll > Papperskorg}. Offentliga sidor slutar omedelbart upptäcka underträdet, men en administratör kan återställa det till ursprunglig sökväg så länge den är ledig. Återställning skriver aldrig över ett nytt aktivt galleri på samma plats.

Välj \ui{Återställ}, \ui{Radera permanent} eller \ui{Töm Papperskorg} i papperskorgspanelen. Permanent radering och Töm Papperskorg går inte att ångra. Töm Papperskorg arbetar i begränsade omgångar så att stora samlingar i papperskorgen inte kräver en enda obegränsad begäran. Radering av enskilda fotografier ingår inte i denna återställningsfunktion på gallerinivå och är fortfarande permanent. En säkerhetskopia är fortsatt ett nödvändigt skydd mot lagringsfel, administrativ permanent rensning eller skador utanför PHP Gallery.
\section{Smarta gallerier}
\index{s00227@smarta gallerier}
\index{v00263@virtuella gallerier}

Ett smart galleri är en sparad fråga, inte en mapp. Öppna \ui{Administration > Gallerier > Smarta gallerier}, lägg till villkor och kombinera dem med nästlade grupper av \ui{AND}, \ui{OR} och \ui{NOT}. Fälten omfattar källgallerier och undergallerier, taggar, tagningsdatum, EXIF/GPS, text, AI-metadata, dubblettstatus, filegenskaper och privata redaktionella betyg.

Använd \ui{Förhandsgranska} för att se aktuellt antal träffar och riktiga bildkort som matchar, och spara sedan och publicera vid behov. Bilder och filer ligger kvar i sina fysiska gallerier, så ändringar av metadata eller betyg ändrar medlemskapet omedelbart. Det redaktionella betyget är privat och oberoende av besökarnas röstning. Administratörer kan även använda \ui{Spara sökning som smart galleri} från den offentliga sökningen.

Välj \ui{Olistat} för åtkomst endast via URL, \ui{Rotgalleri} för att visa det smarta galleriet på startsidan eller \ui{Placering som undergalleri} för att låta fysiska gallerier koppla in det. I varje fysisk galleriredigerare väljer du om kopplingen hör hemma \ui{Ovanför galleriinnehållet} eller \ui{Nedanför galleriinnehållet} och anger ett ordningsvärde. Kopplingar ovanför visas först, sedan vanliga undergallerier och fotografier, och sist kopplingar nedanför. Lika ordningsvärden avgörs konsekvent med det smarta galleriets ID. Befintliga kopplingar får som standard placeringen nedanför med ordning~0.

Samma virtuella galleri kan kopplas under flera fysiska gallerier med oberoende placering och ordning i varje övergalleri. Avsnittet \ui{Används i fysiska gallerier} i det smarta galleriets redigerare visar varje övergalleri och låter dig spara dess placering och ordning eller \ui{Koppla bort} det utan att ändra andra övergallerier. När redigeraren öppnas i administrationens sidopanel stannar åtgärderna i panelen; vanlig formulärinlämning är fortfarande tillgänglig när JavaScript är avstängt. Ett offentligt smart galleri som placeras under ett privat eller annars begränsat övergalleri ärver ändå övergalleriets gräns för offentlig åtkomst för kortet.

Offentliga resultat och antal tillämpar alltid samma åtkomstpolicy för källorna som vanligt offentligt galleriinnehåll. Ett matchande fotografi räknas bara när dess fysiska galleri listas offentligt och är åtkomligt för den aktuella besökaren och fotografiet självt är offentligt. Befintliga regler för lösenord, delningslänkar, ärvd synlighet, opublicerade gallerier och åldersbegränsning gäller därför fortfarande. Rutter med frågesträngar fungerar utan omskrivning, och rena \texttt{/smart/<slug>}-URL:er är valfria.

\subsection{Säkra relationer och reparation}
\index{s00227@smarta gallerier!s00225@skydd mot cykler}

Ett smart galleri får inte leda tillbaka till ett fysiskt galleri som kopplar in det. Servern kontrollerar hela relationsgrafen när regler, kopplingar eller gallerihierarki ändras; filtrering i webbläsaren räknas inte som säkerhet. Det omfattar indirekta vägar genom fysiska undergallerier och andra inkopplade smarta gallerier. Synkronisering av hierarkin från filsystemet och reparation av övergallerier från offentliga sökvägar validerar hela den föreslagna kartan över övergallerier innan någon ny övergallerilänk skrivs. Den aktuella regelredigeraren kan referera till fysiska gallerier men erbjuder inget direkt referensvillkor till smarta gallerier.

Ett fullständigt träd som ordnas med dra och släpp i administrationen valideras en gång i sin slutliga form. De efterföljande mappflyttningarna använder den förvaliderade kartan och skjuter upp hierarkisynkroniseringen tills slutträdet finns, så att ett tillfälligt mellanläge inte kan avvisa en annars giltig omgång. Varje genomförd ändring av den fysiska hierarkin rensar den begärandelokala grafcachen för smarta gallerier innan senare grafberoende arbete fortsätter.

Om äldre eller manuellt ändrade databasdata redan innehåller en cykel stoppar den offentliga visningen den ogiltiga relationen i stället för att fortsätta rekursivt. Kopplingslistan i administrationen markerar relationen för reparation, och \ui{Koppla bort} är fortsatt tillgängligt. Både validering och genomgång vid körning använder försiktiga djup- och nodgränser samt begärandelokal deduplicering, så felaktiga data inte kan expandera samma smarta galleri upprepade gånger. Den vanliga resultatfrågan förblir paginerad och kopierar inte alla matchande fotografier till PHP enbart för att upptäcka cykler.

Migreringen \texttt{202608170002\_smart\_gallery\_attachment\_ordering.php} lägger till fälten för placering och ordning per övergalleri. Tills migreringen är tillgänglig behåller befintliga kopplingar sitt tidigare beteende nedanför, medan placeringsändringar i administrationen nekas i stället för att sparas delvis.

\subsection{Presentationsinställningar}
\index{s00227@smarta gallerier!p00190@presentation}

Som standard följer ett smart galleri aktuellt Tema och webbplatsens presentation. Aktivera \ui{Åsidosätt Temas standardvärden för detta smarta galleri} när den virtuella samlingen behöver egna kolumner, rader, paginering, storleksgränser för miniatyrbilder, miniatyrbildsrenderare, gallerikortslayout, metadataöverlägg, ljuslådebeteende, bildspel, hämtning eller röstning. Sortering och riktning förblir en del av det smarta galleriets definition.

Ett smart galleri ärver inte presentationen från det fysiska galleri där dess kort visas. Ett smart galleri kan visas under flera fysiska gallerier, så det finns ingen enda övergalleripresentation med tydligt företräde. När ingen åsidosättning finns, eller när lagrade presentationsdata saknas eller är ogiltiga, används aktuella standardvärden för Tema och webbplats som reserv. Det fysiska galleriets och fotografiets miniatyrbildsgränser har ändå företräde om en gräns i det smarta galleriet strider mot dem.

\subsection{Ljuslåda över flera resultatsidor}
\index{s00227@smarta gallerier!l00152@ljuslåda}

Den stora visaren är inte begränsad till fotografierna på aktuell HTML-sida. Korten i smarta gallerier kommer ihåg varje fotografis position i hela den behöriga resultatordningen. När besökaren går framåt eller bakåt begär webbläsaren närliggande metadata från servern i små behörighetskontrollerade grupper. Servern upprepar samma regler för det smarta galleriet, åtkomstkontroller för källor, sortering och stabila skiljeregel som sidan använder.

Det gör stora virtuella samlingar praktiska: att öppna ett smart galleri förladdar inte alla original eller all resultatmetadata. Befintliga beteenden för tangentbord, pekskärm, helskärm, zoom, stängning och återöppning samt bildspel delas med vanliga gallerier. Om bildspel är avstängt för det smarta galleriet utelämnas dess bildspelskontroller medan vanliga gallerier behåller sina standardvärden.

\subsection{Hämtningar och säkerhetsgränser}
\index{s00227@smarta gallerier!h00122@hämtningar}

När både webbplatsens hämtningar och det smarta galleriets hämtningsalternativ är~aktiverade kan servern förbereda en ZIP-fil från aktuella behöriga matchande original. Den utvärderar de sparade reglerna och åtkomstpolicyn på nytt i stället för att exponera filsystemssökvägar för webbläsaren. Arkivflödet nekar samlingar med fler än 5~000 matchande bilder och tillämpar även ett tak för originalfilernas sammanlagda byteantal. Befintliga installationer använder ett tak på 2~GiB om inte \texttt{smart\_gallery\_zip\_max\_source\_bytes} har angetts i \texttt{config.php}. Samtidiga begäranden för samma arkivsignatur körs en i taget med ett fillås; efter att låset har erhållits kontrolleras cachen igen, en unik partiell ZIP-fil skrivs och det färdiga arkivet publiceras med ett atomärt namnbyte. Felloggar bevarar endast avgränsad diagnostik för orsak och klass, inte råa undantagsmeddelanden eller privata serversökvägar.

Reglerna är läsbar versionshanterad JSON, aldrig användar-SQL: fält och operatorer tillåts enligt en uttrycklig lista, värden är parametrar i förberedda frågor, djupet är begränsat till fem och högst 50 villkor accepteras. Raderade referenser matchar ingenting, och felaktiga eller okända regelversioner avvisas. Det vanliga migreringsflödet lägger till lagring för smarta galleriers presentation; befintliga smarta gallerier ärver standardvärden för Tema och webbplats tills en åsidosättning sparas.

\section{Kooperativa gallerier}
\index{kooperativa gallerier}
Kooperativa gallerier kopplar samman oberoende installationer med offentligt nåbar HTTPS. Varje ägare behåller sina album och original. Kör väntande migreringar genom det vanliga administrationsflödet och aktivera kooperativa gallerier under Funktioner på varje server. Funktionen är avstängd från början; avstängning bevarar identiteter, vänskap, förslag och inställningar.

Para ihop varje deltagarpar direkt under Inställningar > Avancerat > Vänliga gallerier. Vänskap delar inga fotografier i sig. Välj ett offentligt, listat album utan lösenord eller NSFW under albumsamarbete och skapa dess referenskod. En deltagare kombinerar koderna till ett oföränderligt förslag. Varje administratör granskar deltagare och behörigheter och godkänner uttryckligen på sin egen installation. Leverans och verifiering sker i begränsade steg i sidopanelen. E-postinbjudningar använder den konfigurerade e-posttransporten; att öppna länken är inget samtycke.

Efter aktivering öppnas gemensamma fotografier från panelen eller länken i det lokala albumet. Varje källa behåller attribution och sidindelning. Endast befintliga auktoriserade JPEG/WebP-miniatyrer och förhandsvisningar delas. Original, skyddade album och fjärrskrivning stöds inte. Länkar fungerar utan JavaScript. Saknade miniatyrer kräver normal generering. Nya fotografier får identiteter från hela sökvägen inklusive filändelsen. Gamla namn bevaras utan massgenerering, men tvetydigt ägarskap avvisas.

Återkallat deltagande, ändrat källskydd och lokalt återkallad vänskap blockerar efterföljande auktoriserade läsningar. Onåbara deltagare begränsas till högst 120 sekunders verifieringsgiltighet; redan hämtade pixlar kan inte återkallas. Nya deltagare kräver ett separat förslag med nytt godkännande från alla. Teknisk förnyelse kan ske vid användning; schemaläggning är valfri. Saknat eller okänt obligatoriskt schema avvisar delning och ändringar. Se \filepath{docs/COOPERATIVE_GALLERIES.md} för driftsgränser.

\subsection{Simulatorposition i Windows}
WinApp 0.3.2 använder automatiskt kamerans världsposition i MSFS 2024 och faller tillbaka till flygplanets position vid ett begränsat kamerafel. MSFS 2020 använder direkt flygplanets position genom samma SimConnect-runtime. Koordinater är valfria; en otillgänglig simulator stoppar inte uppladdningen. Diagnostiken skiljer mellan positionskällorna. Installationsprogrammet begär att programmet avslutas före filbyte. Slutför pågående arbete före uppdatering. Verifiera verkliga simulatorer och uppgradering på måldatorn.

Den paketerade uppdateraren laddar Windows-API:er uttryckligen från System32 för att undvika en kollision mellan textfilen VERSION och Windows versionsbibliotek. Starta det verifierade installationsprogrammet 0.3.2 manuellt vid uppgradering från 0.3.1. Den äldre appen kopierar sin egen hjälpprocess före installationen, så den gamla hjälparen kan fortfarande rapportera DLL-felet efter att appen ersatts. Framtida uppdateringar från den korrigerade appen 0.3.2 använder den reparerade hjälparen; verifiera den installerade överlämningen i Windows.

Flygplanssvar kopplas till begärans och definitionens identifierare; begärans USER-väljare kräver inte att svarets objektidentifierare är noll. Begränsad diagnostik registrerar mottagningsmetadata, behandlingsordning och konkreta skäl till avvisade paket. Felorsaker förblir läsbara efter sökvägsredigering och aktivitetsmeddelanden visar en enda SimConnect-markering. Objektidentifierarfiltret var ett bekräftat kodfel, men det rapporterade FS2020-tidsfelet registrerade inte den returnerade objektidentifieraren. Testa korrigeringen på den berörda datorn; simulerade tester bevisar inte ett verkligt simulatorresultat.

Den installerade uppladdaren kontrollerar sin egen stabila GitHub-installationsversion högst en gång per timme; \ui{Check for updates} startar en manuell kontroll. Det senaste automatiska försöket sparas separat för den aktuella användaren och finns kvar efter programomstart. CMS-taggar bestämmer inte dess version. Den högsta numeriska versionen kräver SHA256 från GitHub-filen eller en matchande \filepath{winapp-update.json} från samma utgåva. Saknade eller motstridiga uppgifter stoppar uppdateringen; identiska installationsprogram i flera CMS-utgåvor accepteras. Galleriets inloggningsuppgifter skickas aldrig till GitHub.

\ui{Update} hämtar installationsprogrammet med förlopp och avbrottsmöjlighet, verifierar det, stoppar nytt arbete och väntar högst 120 sekunder på att bakgrundsuppgifter avslutas säkert. Ett misslyckat avslut nekar installationen och uppladdaren förblir tillgänglig. En extern hjälpprocess kontrollerar filen igen och begär Windows UAC-godkännande. Självuppdateringen tvångsavslutar aldrig arbetare; vanliga manuellt startade installationsprogram behåller sin avslutningspolicy. Det verkliga installationsförloppet kommer genom en lokal kanal för den aktuella sessionen. Efter framgång startar uppladdaren som ursprunglig användare med sparade inställningar och uppgifter. Ingen automatisk återställning till tidigare version finns. Körning från Python eller källkod erbjuder det verifierade installationsprogrammet för manuell start. Om Windows kräver omstart ska Windows startas om innan uppladdaren öppnas igen; avbruten UAC öppnar den befintliga installationen på nytt.


\section{Lägga till, beskriva och organisera fotografier}
\label{sec:photo-workflow}
\index{f00084@fotografiarbetsflöde}
\index{b00035@bildmetadata}
\index{f00085@fotoorganisering}

\subsection{Välj en uppladdningsmetod}
\index{u00256@uppladdningsmetoder}

Webbläsaruppladdning passar vardagliga tillägg. Välj ett galleri, välj filer, granska omgången och följ förloppet medan servern validerar medier och förbereder poster och miniatyrbilder. Filsystemsupptäckt passar stora befintliga mappträd eller överföringar via FTP och webbhotellsverktyg. Automatisering och WebDAV passar återkommande eller mobila arbetsflöden när deras avgränsade inloggningsuppgifter har konfigurerats.

Ladda upp i meningsfulla grupper för ett stort evenemang. Mindre grupper gör fel lättare att hitta och låter dig granska titlar, ordning och synlighet stegvis.

\subsection{Skriv användbara bildtitlar}
\index{b00040@bildtitel}
\index{b00039@bildtexter}
\index{a00008@alternativtext}

En titel identifierar fotografiet snabbt. En beskrivning eller bildtext kan förklara personer, plats, aktivitet, tekniskt sammanhang eller en berättelse utanför bildramen. Användbara exempel:

\begin{itemize}
  \item \emph{Karlsbron före soluppgången};
  \item \emph{Det första tågets ankomst till den restaurerade stationen};
  \item \emph{Familjegrupp i trädgården, juni 1998};
  \item \emph{Slutlig inflygning under lördagens uppvisning}.
\end{itemize}

Meningsfull text förbättrar sökning, tillgänglighet, senare identifiering och upplevelsen för läsare som vill ha mer än ett visuellt rutnät. Filnamn kan förbli driftsdetaljer när offentliga visningsinställningar prioriterar redigerade titlar.

\subsection{Använd taggar som ett andra organiseringslager}
\index{t00241@taggar!p00189@praktisk användning}
\index{t00242@taggningsstrategi}

Skapa ett litet ordförråd innan du lägger till många nästan likadana taggar. Föredra en stabil term, t.ex. \emph{flygplan}, framför flera oavsiktliga varianter. Taggar kan representera ämnen, personer, platser, tekniker, utrustning, årstider eller redaktionell status.

I ett researkiv kan gallerier representera resor och taggar återkommande teman. I ett familjearkiv kan gallerier representera år och händelser medan taggar identifierar personer. I en professionell portfolio kan gallerier representera kunder och taggar identifiera leveranstyp, plats och specialitet.

Taggfält kan föreslå befintliga taggar medan du skriver. Att återanvända ett förslag undviker oavsiktliga stavningsvarianter och håller taggarnas landningssidor sammanhängande. Förslaget är bara ett redigeringshjälpmedel; den vanliga sparåtgärden avgör fortfarande vilka taggar som lagras för galleriet eller fotografiet.

\subsection{Ordna visningsföljden}
\index{b00036@bildordning}
\index{b00024@berättande}

Ordning förvandlar en samling till en följd. Börja med en inbjudande bild, etablera plats eller sammanhang, variera vida vyer och närbilder, håll nära relaterade detaljer tillsammans och avsluta naturligt. Kronologisk ordning fungerar bra för evenemang, medan visuell rytm kan passa portfolior och utställningar.

De offentliga kontrollerna för ändrad ordning hjälper administratörer att justera synliga kort. Omordning av fotografier på en offentlig sida är också begränsad till den aktuella synliga pagineringsdelen och sparas genom den autentiserade adressen för bildordning. När Bildhanteraren är aktiv ändrar dragning av markerade bildkort över ett synligt undergallerikort gesten från omordning till en validerad flytt till det undergalleriet. Paginering kan dela upp en stor samling, så granska övergångar kring sidgränserna och inte bara den första sidan.

\subsection{Granska EXIF och platsuppgifter}
\index{e00072@EXIF!a00010@användargranskning}
\index{p00187@platsintegritet}

EXIF kan berika ett fotografi med tagningsdatum, kamera, objektiv, exponering och GPS-koordinater. Uppgifterna stöder kartor, datumförslag, dubblettunderlag och tekniskt sammanhang. Granska GPS-synligheten innan du publicerar fotografier tagna vid hem, skolor, privata arbetsplatser eller känsliga naturplatser.

\section{Publicering, delning och målgruppsexempel}
\label{sec:audiences}
\index{p00197@publiceringsflöde}
\index{m00170@målgrupper}
\index{a00014@användningsfall}

\subsection{Offentlig portfolio}
\index{p00188@portfolio}

En offentlig portfolio gynnas av ett litet antal starka gallerier på högsta nivån, konsekventa omslag, korta beskrivningar och omsorgsfull ordning. Använd taggar för specialiteter som går över projektgränser. Håll originalarkiv i privata eller opublicerade grenar och kopiera utvalda arbeten till offentliga presentationsgallerier när arbetsflödet kräver ett separat urval.

Testa portfolion på en telefon, en högupplöst skärm och en långsammare anslutning. Progressiv skärpning är standard och säker reserv för bildkort i det öppna galleriet; responsivt val förblir ett stött alternativ.

\subsection{Privat familjearkiv}
\index{f00074@familjearkiv}

Organisera efter år, familjegren eller händelse. Använd lösenordsskyddade övergallerier för att skapa ett begripligt privat område och placera sedan relaterade undergallerier där. Lägg till namn och datum medan minnena finns kvar. Skannade fotografier har särskild nytta av mänskliga beskrivningar eftersom inbäddade kamerametadata kan saknas.

Dela den skyddade adressen och lösenordet via en lämplig privat kanal. Förklara att mottagarna kan navigera i undergallerier utan att få ett administratörskonto. Granska regelbundet om gamla länkar och lösenord fortfarande tjänar den avsedda målgruppen.

\subsection{Kundleverans}
\index{k00141@kundgalleri}
\index{k00139@korrekturflöde}

Skapa ett skyddat galleri per kund eller uppdrag. Ladda upp korrekturbilder, ordna dem i användbara grupper och använd beskrivningar för att förklara förväntningar kring urval eller hämtning. Röstning kan ge enkel insamling av önskemål när den är~aktiverad för galleriet. Ett slutligt leveransgalleri kan innehålla godkända filer i avsedd ordning.

Delningslänkar ger en annan kontrollerad ingång där de är konfigurerade. Testa exakt den väg kunden ska följa i ett privat webbläsarfönster innan du skickar adressen.

\subsection{Evenemangs- och gemenskapsgalleri}
\index{e00071@evenemangsgalleri}

Ett evenemangsgalleri kan börja som opublicerat medan fotografierna kommer in. Använd undergallerier för dagar, scener, lag, platser eller aktiviteter. Publicera övergalleriet när det första användbara urvalet är klart och lägg sedan till senare grupper utan att ändra den välkända offentliga adressen.

Offentlig röstning kan lyfta fram publikens favoriter. Taggar knyter samman återkommande deltagare eller ämnen. Hämtningar ger ett bekvämt sätt att överföra en samling när evenemangets policy tillåter det.

\subsection{Rese- och platsarkiv}
\index{r00203@researkiv}
\index{k00133@kartor!a00013@användning för resor}

Gallerier för land, region, stad och resa ger naturlig navigering. Datumperioder etablerar kronologin medan GPS-kartor skapar en geografisk läsväg. Granska koordinater för boenden och privata platser. En beskrivande inledning kan förklara resans rutt, syfte och historiska sammanhang.

\subsection{Historiskt och institutionellt arkiv}
\index{h00121@historiskt arkiv}
\index{i00125@institutionellt arkiv}

Historiska samlingar gynnas av noggranna källanteckningar, formuleringar om osäkra datum, namn, platser och proveniens. Använd beskrivningar för att skilja bekräftade fakta från familjens eller institutionens minnen. Taggar kan koppla ihop personer, byggnader, organisationer och tidsperioder över många gallerier.

Skapa säkerhetskopior före stora metadataprojekt och exportera information genom tillgängliga verktyg när du planerar extern bevaring. Galleriet kan ge ett tillgängligt presentationslager samtidigt som ursprungliga bevaringsmetoder fortsätter vid sidan om.

\section{Förstå dina galleridata}
\label{sec:user-data}
\index{d00054@dataägande}
\index{d00049@databas!f00092@för galleriägare}
\index{f00079@filsystem!f00092@för galleriägare}

Galleriägare arbetar med två samverkande lagringssystem: vanliga filer och en databas. Att förstå deras roller underlättar säkerhetskopiering, migrering och felsökning.

\subsection{Vad som finns i gallerimapparna}

Gallerimapparna innehåller fotografier och mapphierarki, så grundsamlingen är fortfarande igenkännbar genom FTP, webbhotellets filhanterare, ett säkerhetskopieringsverktyg eller en lokal kopia. Flyttning eller namnbyte utanför programmet kan kräva upptäckt eller synkronisering så att databasen får det nya tillståndet.

\subsection{Vad som finns i databasen}

Databasen kommer ihåg titlar, beskrivningar, datum, synlighet, lösenord och åtkomstkonfiguration, offentliga sökvägar, ordning, taggar, röster, bildmått, EXIF-index, kontrollsummor, miniatyrbildsposter, administratörskonton, inställningar, granskningsinformation och arbetsflödestillstånd.\footnote{Lösenordsvärden använder hash eller skyddad konfiguration som passar deras syfte. En databasexport innehåller ändå känslig driftinformation och behöver säker lagring.}

Databasen gör sökning och administration snabba. Den låter också en fil ha en lättbegriplig titel, utförlig beskrivning, taggar, vald ordning och åtkomstregler utan att ändra originalbildens byte.

\subsection{Hur användare drar nytta av databasen}
\index{d00049@databas!f00093@fördelar}

Vardagliga fördelar:

\begin{itemize}
  \item söka efter ord i galleri- och bildbeskrivningar;
  \item öppna ett galleri via en stabil offentlig sökväg;
  \item visa fotografier i genomtänkt ordning;
  \item komma ihåg taggar och besökarröster;
  \item skydda ett galleri med ärvda åtkomstregler;
  \item hitta exakt dubblerat innehåll med lagrade kontrollsummor;
  \item förbereda lämpliga miniatyrbilder utan att läsa om varje originalfil på varje sida;
  \item registrera migreringar och underhållsbeslut i en granskningslogg.
\end{itemize}

\subsection{Säkerhetskopiera som en helhet}

En fullständig säkerhetskopia omfattar galleriträdet, databasexporten, lokal konfiguration och installationens egna anpassade resurser. Ge delarna samma säkerhetskopieringsdatum så att de beskriver ett sammanhängande tillstånd. Förvara en kopia utanför webbservern och testa återställning på en separat plats.
\section{Daglig drift och skötsel}
\label{sec:daily-care}
\index{d00048@daglig administration}
\index{m00171@månatligt underhåll}
\index{z00278@ägarens checklista}

\subsection{Efter varje viktig uppladdning}

Granska målgalleriet, antalet fotografier, ordning, titlar, synlighet och miniatyrbildernas utseende. Öppna flera fotografier i ljuslådan och kontrollera galleriet anonymt. För en skyddad samling upprepar du vägen via lösenord eller delningslänk från en ny privat session.

\subsection{Månatlig genomgång}

En månatlig underhållskontroll bör omfatta:

\begin{enumerate}
  \item Granska Systemstatus och väntande migreringar.
  \item Kontrollera nya händelser i administrationens granskningslogg.
  \item Undersök tillgängliga programuppdateringar och deras utgåveinformation.
  \item Bekräfta att schemalagda eller manuella säkerhetskopior har slutförts.
  \item Öppna representativa offentliga och skyddade gallerier.
  \item Granska miniatyrbildsunderhåll när originalfiler eller visningsinställningar har ändrats väsentligt.
  \item Ta bort eller återkalla inloggningsuppgifter och delningsmetoder som har fyllt sitt syfte.
\end{enumerate}

\subsection{Före en större offentlig lansering}
\index{l00148@lanseringschecklista}

Kontrollera titlar, beskrivningar, omslagsbilder, stavning, mobillayout, integritet, GPS-exponering, hämtningar, kartor, röstning, sökresultat och förväntningar kring kontakt. Testa med en långsam nätverksprofil om målgruppen kan använda mobilanslutningar. Skapa en säkerhetskopia omedelbart före lanseringen.

\subsection{När någon annan hjälper till med administrationen}

Ge varje administratör ett eget konto där kontomodellen stöder det avsedda arbetsflödet. Förklara galleriomfattning, integritetskrav, namnkonventioner, taggordförråd, ansvar för säkerhetskopiering och raderingsrutiner. Enskilda konton förbättrar spårbarheten i granskningsposter och håller personliga dubblettgranskningsregister åtskilda.

\section{Den offentliga besökarupplevelsen}
\label{sec:public}
\index{o00177@offentligt galleri}
\index{d00061@dirigering}

\subsection{Vad besökaren ser}
\index{b00030@besökarupplevelse}

En besökare öppnar webbplatsens startsida eller följer en direktlänk till ett galleri eller fotografi. Sidan visar webbplatsnamn, navigering, gallerititel, beskrivning, brödsmulor, taggar, undergallerier och fotografier som besökaren har rätt att se. Skyddat innehåll kräver rätt åtkomststeg innan privat material visas.

Varje galleri har en offentlig adress som kan bokmärkas eller delas. En direkt bildadress öppnar gallerisammanhanget och leder besökaren till fotografiet. Meningsfulla titlar och stabila offentliga sökvägar gör länkar begripliga i meddelanden, bokmärken och sökresultat.

\subsubsection{Ett typiskt besök}

En vanlig besökarväg:

\begin{enumerate}
  \item Öppna startsidan och välj ett intressant galleriomslag.
  \item Läs galleriets inledning och följ undergallerier eller taggar.
  \item Rulla genom mindre förhandsvisningsbilder.
  \item Öppna ett fotografi i den stora visaren.
  \item Bläddra bland närliggande fotografier med knappar, tangenter, en bildremsa eller det valda visningsläget.
  \item Öppna en karta, rösta, sök eller hämta när galleriet erbjuder dessa åtgärder.
  \item Använd brödsmulorna för att återvända till en större samling.
\end{enumerate}

\subsection{Gallerihierarki och paginering}
\index{p00182@paginering}
\index{g00100@gallerihierarki}

\subsubsection{Direkt innehåll och stabil ordning}

Ett galleri kan visa mindre gallerier först och fotografier därefter. Till exempel kan \emph{Italien 2026} innehålla \emph{Rom}, \emph{Florens} och \emph{Venedig}, och samtidigt visa några översiktsfotografier direkt på Italiensidan.

Paginering delar en lång samling i hanterbara sidor. Det hjälper första sidan att visas snabbt och gör rullningen bekväm. En liten utställning kan se bäst ut på en enda sida. Ett bröllop med 1~500 fotografier har vanligen nytta av flera sidor eller undergallerier. Den stora bildvisaren kan fortsätta genom galleriets större bildföljd och läsa in ytterligare information efter behov.\footnote{Paginering ändrar hur många kort som visas samtidigt. De underliggande fotografierna och deras ordning bevaras.}

Långa offentliga listor visar en flytande kontroll \ui{Till toppen} efter att besökaren har rullat in i gallerilistan. Den använder webbläsarens mjuka rullning, försvinner utanför den aktiva listan och förblir dold medan ljuslådan eller helskärmsvisaren är öppen så att den inte konkurrerar med bildkontrollerna.

\subsection{Sökning, taggar och navigering}
\index{s00240@sökning}
\index{t00241@taggar}

Sökning kan omfatta hela den offentliga webbplatsen eller begränsas till aktuellt galleri och dess undergallerier. En besökare som tittar på \emph{Familjearkiv} kan söka efter ett personnamn bara i den grenen, medan en besökare på startsidan kan söka i alla offentliga samlingar. Den aktuella söktjänsten börjar vid två synliga tecken i frågan, ger en begränsad blandning av galleri- och bildresultat och rangordnar starkare träffar på titel, namn och taggar före svagare textträffar. Standardbegäran ger upp till 12 sammanslagna resultat, och tjänsten begränsar en begäran till högst 30 resultat.

Sökunderlaget omfattar offentliga galleri- och bildtitlar och beskrivningar, bildfilnamn, galleri- och bildtaggar och deras offentliga beskrivningar, aktivt underhållet språkinnehåll där dess schema är tillgängligt samt intern sökbar AI-text när AI-metadata är~redo. Åtkomstfiltrering sker som en del av frågan: privat, oåtkomligt eller icke-offentligt bildinnehåll returneras inte enbart för att texten matchar. En grensökning begränsar även listsammanhanget till det galleriets underträd i stället för att söka i hela katalogen och dölja orelaterade träffar i efterhand.

Taggar ger en annan väg genom arkivet. Taggen \emph{nattfotografi} kan knyta samman Prag, London och Tokyo även när fotografierna ligger i separata resegallerier. Brödsmulorna visar aktuell plats i hierarkin, och favoritgenvägar placerar viktiga gallerier direkt i webbplatsens sidhuvud.
\section{Medie- och galleriadministration}
\label{sec:admin}
\index{a00002@administration}

\subsection{Galleriöversikt och granskade ändringar}
\index{galleri!funktionsplanering}
Galleriytan visar direkta och underordnade fotografier, undergallerier samt på/av/blandade funktionslägen. Expandera en gren och förbered Kartor, Filnamn, Röstning eller Picture Game för roten och dess undergallerier. Förberedelsen sparar ingenting. Granska exakt omfattning och beroenden och tillämpa sedan den granskade planen. Överlappande avsikter behåller ordningen. Picture Game kräver röstning; avstängd röstning stänger också av spelet. Avbryt eller kasta utkastet för att behålla lagrade värden. Granska igen om en annan ändring gjort planen inaktuell.

Snabbkontrollen för lösenord ändrar bara det aktuella galleriets eget lösenord. Ange ett nytt lösenord eller ta bort det egna. Token- och överordnat skydd, synlighet och övriga inställningar finns kvar. En gammal revision nekas. Vanliga formulär finns utan JavaScript; sidopanelens åtgärder slutförs på plats och panelen hålls öppen.

Inloggade administratörer kan skapa ett rotgalleri på den offentliga startsidan. Uppmaningen om första galleri kräver en bekräftat tom katalog. Innehåll och plats/synlighet grupperas separat; sekundära val kan fällas ut. Innehållsspråk, taggar, SimBrief-utkast och sparade standardval finns kvar.

\subsubsection{Uppladdningsinställningar och mobila anslutningar}
Öppna \ui{Inställningar > Uppladdningar} för format, namnbyte, arbetare/gränser, ZIP-omgångar, källblock för miniatyrbilder och mobil WebDAV. Kopplade gränser normaliseras tillsammans; storlekar visas i MB. Aktivera äldre menylänkar vid behov. Dolda länkar stänger inte av befintliga adresser eller tjänster.

Mobila åtgärder uppdaterar sin egen yta. Kopiera ett nytt engångslösenord genast. Om följande listläsning misslyckas förblir anslutningen sparad och lösenordet synligt i detta svar; försök läsa listan separat. Avstängda mobila uppladdningar läser inte anslutningslistan.

\subsubsection{Utseende, språk och egna stilmallar}
Utseende kombinerar grupperade fält med direkt förhandsvisning och en begränsad mus-/tangentbordsdelare på breda skärmar. Panelbredden är en lokal webbläsarinställning. Webbplatsbilder grupperar profil/bakgrund; Layout skiljer genvägar, rutnät, kort och bildvisarens standardval. Språk har Inställningar, Design, Redigerare och Diagnostik. Designprovet använder samma formulärs erbjudna språk och offentliga standardval och sparar ingenting självt.

Lämna ersättningsvalet tomt för att behålla installerad CSS. Välj uttryckligen en mall eller ladda upp en giltig CSS-fil; uppladdning har företräde. Ersättning/återställning verifierar lagring först och återställer tidigare fil om markören inte kan sparas. Behåll tillgänglig kopia och följ den säkra vägledningen vid ofullständig återställning. Driftsättning bevarar installationens \filepath{public/assets/custom.css}.

\subsubsection{Kortorientering och driftstatus}
Vertikala kort placerar bilden ovanför texten, horisontella bredvid. Den vanliga orienteringsmigreringen bevarar tidigare utseende i inställningar, fysiska och Smarta Gallerier, sidecars, importerade metadata och Papperskorg. Säkerhetskopiera databas, gallerier och Papperskorg tillsammans. Lös rapporterat problem efter avbrott och försök igen; dokumentmarkörer hindrar dubbel omvandling. Nyinstallationer använder vertikalt standardval. Alla taggar finns utan JavaScript; expansion och rullning fungerar även på kort.

Offentlig sökning är på som standard bara utan sparat val och när funktionen är tillgänglig. Uppdateringar synkroniserar versioner, notiser och API-status; automatisk metadatakontroll görs högst varje timme, och avstängda automatiska uppdateringar lämnar äldre bakgrundsjobb pausade. OurAirports använder veckovis aktualitet och en timmes felpaus. Papperskorgen visar återstående tid och säker raderingsmöjlighet. Även presentationsschemafel kräver åtgärd i Systemstatus.

Den fristående uppladdaren är fortfarande 0.3.2. Bygget levererar installationsfil och motsvarande genererade \filepath{winapp-update.json} tillsammans under \filepath{winapp/dist/0.3.2/}. Båda hör till samma publicerade utgåva. Bygget ersätter inte verifiering av installerad Windows-app eller simulator.

\subsection{Administrationens arbetsyta}

Efter inloggning ser administratören privata kontroller för översikten, gallerier, fotografier, tema, taggar, underhåll, uppdateringar, loggar och kontoinställningar. Många redigeringar öppnas i panelen till höger, så att galleriet förblir synligt medan ett objekt ändras och sparas.\footnote{En fullständig administrationssida finns kvar för arbetsflöden som behöver mer utrymme eller när webbläsarens utökade funktioner inte är tillgängliga.}

Det offentliga galleriet visar också sammanhangsberoende administrationskontroller när en vanlig administratörssession är aktiv. Pennkontroller för gallerier och fotografier öppnar samma fullständiga redigerare i sidopanelen. Kompakt galleriborttagning använder det återställningsbara papperskorgsflödet när Papperskorg är~aktiverad; om Papperskorg stängs av använder framtida galleriborttagningar åter den äldre tjänsten för permanent radering av underträd. Kompakt bildborttagning har avsiktligt etiketten \ui{Ta bort fotografi ... från CMS}: denna snabba åtgärd på den offentliga sidan tar bort bildraden i katalogen och kör inte hela tjänsten för radering av originalfiler. Använd det vanliga flödet för bild- eller massradering när även originalfilen och genererade derivat ska raderas. Anonymt förhandsvisningsläge döljer dessa sammanhangsberoende administrationskontroller.

Administrationens galleritabell kan filtrera rader efter synlighet: alla, opublicerade, offentliga eller privata. Filtreringen är ett visningshjälpmedel i webbläsaren, inte en ändring av åtkomstpolicyn. En rad som filtret döljer avmarkeras också, och huvudmarkeringen töms, så att en dold gammal kryssruta inte oavsiktligt ingår i nästa massåtgärd. Sammanfattningen anger antalet visade och möjliga gallerier.

Galleriträdet stöder utfällning och hopfällning per gren samt \ui{Fäll ihop alla} och \ui{Fäll ut alla}. Det kommer ihåg vilka gallerigrenar administratören har fällt ihop. Webbläsaren skickar de hopfällda galleriernas ID:n genom en autentiserad CSRF-skyddad åtgärd och programmet lagrar detta gränssnittstillstånd, så att samma grenar förblir stängda när du återvänder till administrationens träd i stället för att hela hierarkin öppnas igen. Undergallerier som döljs av en hopfälld gren avmarkeras också före massarbete. Omordning uppdaterar trädkontrollerna när ett galleri får eller förlorar undergallerier.

\subsection{Tillförlitliga ändringar på plats i sidopanelen}
\index{a00002@administration!z00280@ändringar i sidopanelen}
\index{s00219@sidopanel!s00226@slutförda ändringar}

När JavaScript är tillgängligt slutförs en beständig åtgärd som startas i administrationens högra sidopanel i samma panel. Panelen förblir monterad och öppen, adressen i webbläsaren är oförändrad och sidan laddas inte om enbart för att visa resultatet. Det gäller skapande, redigering, flyttning och radering av gallerier; uppladdning, redigering, radering, flyttning, omordning, synlighet, NSFW och omslagsval för bilder; skanning och import; EXIF-datumförslag; miniatyrbildsskapande; Mediefilnamnsbyte; Metadataorganiseraren; Dubblettdetektorns gransknings- och raderingsåtgärder; inloggningsuppgifter för Uppladdningsautomatisering; hämtningar till mål i Gallerimigrering; taggredigering samt skapande, redigering, placering, bortkoppling, duplicering och radering av smarta gallerier. En fristående administrationssida och vanligt POST/omdirigeringsbeteende finns kvar när JavaScript är avstängt eller rutten öppnas direkt.

Bara den senaste begäran att öppna panelen får ersätta innehållet; stängning eller öppning av ett annat objekt upphäver ett äldre svars ägarskap. Tangentbordsfokus stannar i den öppna panelen och återgår till kontrollen som öppnade den. Osparad vanlig text kan hållas i begränsat webbläsarminne för aktuell sida och återställas när samma formulär kommer tillbaka. Lösenord, token, API-nycklar, filfält, åtgärdsidentiteter och andra privata värden utesluts, och ingenting sparas i webbläsarlagring.

Det sparade databas- och filsystemstillståndet förblir styrande. Efter sparandet beskriver servern den ändrade entiteten, panelfragmentet som kan behöva renderas om och varje påverkat offentligt sammanhang. Stabila identifierare för gallerier, bilder, taggar och smarta gallerier används i stället för att förutsätta att den gamla URL:en fortfarande identifierar det sparade objektet. Det är viktigt efter namnbyte på gallerimapp eller URL-namn, ändrat övergalleri eller ändrat URL-namn för en tagg: programmet kan hämta den nya kanoniska renderingskällan samtidigt som webbläsarens synliga adress lämnas orörd. Vid flyttning beskrivs både käll- och målgalleriet; vid ändrat övergalleri beskrivs gammalt och nytt övergalleri samt det flyttade galleriet; placeringsändringar för smarta gallerier beskriver varje påverkat rot- eller fysiskt övergallerisammanhang.

En lyckad fragmentbegäran är i sig inget bevis för att ändringen är synlig. Den uppdaterade serverrenderade HTML-koden måste uppfylla ett observerbart villkor för åtgärden där ett sådant finns. Exempel är närvaro eller frånvaro av galleri eller kort och fullständiga antal, exakt effektiv synlighet, sparad gallerirevision, bildfrånvaro, synlighet eller NSFW-status, galleriets bildrevision för ändringar utanför aktuell sida, omslagsbildens identitet, sparad bildordning, taggidentitet samt smarta galleriers placering och ordning. Pagineringsmedvetna kontroller accepterar ett direkt synligt matchande kort innan sammanlagda antal används, medan ett resultat utanför sidan fortfarande väntar på att det styrande totalantalet ska stämma. Reparation av miniatyrbildscache är det avsiktliga undantaget: den uppdaterar påverkade sammanhang men har inget stabilt databasbaserat villkor för offentlig HTML.

Oberoende renderingsbegäranden kan kortvarigt returnera inaktuella data på delade webbhotell. PHP Gallery utför därför en enda centraliserad begränsad återförsöksföljd först efter att själva ändringen har lyckats och det angivna villkoret ännu inte går att observera. Varje slutförande får en ny åtgärdsgeneration; en nyare ändring avbryter eller ersätter äldre hämtningar, återförsöksväckningar och panelsvar för samma sammanhang. Strukturellt giltig men inaktuell HTML verifieras före ersättning, så att ett äldre svar inte kan skriva över ett nyare synligt tillstånd. Uppdateringar förbigår webbläsarcache och bevarar relevant tillstånd för paginering, sortering, språk, filter, aktiv flik, panelrullning och sidrullning.

Om sparandet lyckas men den uppdaterade vyn inte kan verifieras inom den begränsade återförsöksbudgeten rapporterar panelen ett synkroniseringsproblem och behåller den sparade serverändringen. Den beskriver inte felaktigt ändringen som misslyckad, laddar inte tyst om sidan, navigerar inte bort och skriver inte om webbläsarhistoriken. Fel i autentisering, behörighet, CSRF, schemaberedskap, sökvägssäkerhet och validering ger ett förväntat strukturerat fel till det utökade arbetsflödet och försöks inte igen. Dynamiskt uppdaterade kontroller förblir aktiva eftersom panelen använder delegerad hantering som fungerar genom livscykeln, i stället för bindningar som försvinner när fragment ersätts.

Webbläsarstödd skapande och uppladdning har ytterligare skydd. Att välja webbläsarbearbetning är ett uttryckligt körningsval: saknad webbläsarfunktion eller misslyckad förberedelse av en arbetare stoppar före sparandet i stället för att tyst växla till servergenerering av miniatyrbilder. PHP validerar att varje obligatorisk förberedd miniatyrbildsstorlek och varje format finns innan original lagras. När webbläsarsparandet börjar tolkas reservmetadata för direktsidan aldrig som tillstånd att spela upp den klassiska skapande- eller uppladdningsbegäran igen, vilket förhindrar dubbla gallerier och fotografier. Den konfigurerade ZIP-omgångsstorleken är ett mjukt packningsmål: ett atomärt fotografipaket som överstiger målet kan skickas ensamt i en omgång, men den effektiva PHP-uppladdningsgränsen förblir en hård gräns. Uppladdningsautomatisering i Windows kan skicka hela kompatibilitetsmatrisen JPEG+WebP; en modern server med enbart WebP accepterar obligatoriska WebP-varianter och hoppar säkert över giltiga extra JPEG-varianter, medan felaktiga, icke matchande, okända storlekar eller miniatyrbildspaket som inte stöds fortfarande ger hårda fel.

\subsection{Centralt inställningsnav}
\index{i00126@Inställningar!c00045@centralt nav}

Administrationens navigering innehåller en central sida \textbf{Inställningar} för webbplatsövergripande konfiguration. Allmänt samlar webbplatsadress, namn, administrationsspråk, offentligt språk och navigering. Övriga avsnitt omfattar Offentligt utseende, Innehåll, Medier och bläddring, Uppladdningar och automatisering, Integritet och diagnostik samt Avancerat. Äldre länkar för webbplatsadress leder till Allmänt. Stabila avsnittslänkar och vanliga formulär fungerar utan JavaScript.

Med JavaScript laddas avsnitten när de öppnas och sparas på plats. En fast sparrad visar ändrade inställningar och erbjuder att spara eller återställa utkastet. Avancerade sammanfattningar förblir hopfällda. Ändringar av adress eller administrationsspråk erbjuder en uttrycklig länk till den nya adressen eller språket. Administrationspanelen visar först sitt gränssnitt och laddar sedan Översiktens summor eller gallerilistan vid behov; misslyckade läsningar kan upprepas. Översikten räknar indexerade original utan filskanning eller sökning efter omslagsbilder.

Ett Spotlight-liknande fält söker bland inställningar medan du skriver. Det matchar namn, beskrivningar, kategorier och tekniska namn, inklusive kontroller som ägs av specialiserade administrationssidor. När du väljer ett resultat öppnas relevant avsnitt eller specialsida och matchande kontroll får fokus. Piltangenter, Enter, Escape och rensningsknappen kan användas utan att lämna tangentbordet.

Vanliga globala inställningar kan redigeras direkt från navet, däribland webbplatsnamn, offentligt språk, URL-omskrivning, offentlig sökning där den är tillgänglig, miniatyrbildsrendering, den globala standarden för EXIF/GPS-visning och utvecklingsdiagnostik. Andra inställningar stannar på sina specialsidor när de kräver mer sammanhang, inloggningsuppgifter, filuppladdningar, destruktiva åtgärder eller större redigerare. Exempel är Temas layout, galleribegränsade API-nycklar för uppladdning, telemetri, kontosäkerhet, Google/OpenAI-konfiguration, egen CSS, profilering, språkpaket och databasunderhåll.

Navet tar inte bort befintligt arv. Taggsidors layout kan fortsätta ärva det globala rutnätet och kortutformningen, medan gallerispecifika val för kort, ljuslåda och EXIF/GPS stannar hos det galleri som äger dem. Sammanhangsspecifika galleri- och bildvärden blir inte globala inställningar.

Hemligheter visas aldrig i sökresultat eller sammanfattningar. Lösenord, token, API-in\-logg\-nings\-upp\-gif\-ter och liknande värden visas bara som säkra tillstånd, t.ex. Konfigurerat, Inte konfigurerat eller en länk till sin specialsida.

\subsubsection{Vägledd installationsguide}
\index{i00126@Inställningar!i00124@installationsguide}
\index{i00124@installationsguide}

Välj \ui{Installationsguide} längst upp på \ui{Inställningar} för att starta eller återuppta ett väglett konfigurationsutkast. Guiden följer samma åtta avsnitt som navet. Korta underavsnitt håller relaterade val tillsammans; \ui{Avancerat} och \ui{Expert- och specialinställningar} visar mindre vanliga alternativ. När JavaScript är avstängt är allt underavsnittsinnehåll fortfarande läsbart och vanlig formulärnavigering fungerar.

Varje inställning som stöds har en förklaring och ett exempel. Låt kryssrutan för inkludering vara markerad för att förbereda ett värde, eller avmarkera den för att bevara ursprungsinställningen. \ui{Hoppa över} bevarar hela avsnittet. \ui{Nästa}, \ui{Tillbaka} och navigering mellan underavsnitt sparar inte programinställningar. Temaval som stöds delar den befintliga direktförhandsvisningen; en förhandsvisning är tillfällig och ändrar inte den offentliga webbplatsen.

Guiden stöder säkra centrala inställningsvärden, grundläggande temafärger och layout, rutnäts- och kortalternativ samt huvudtaggar, begränsade uppladdnings- och telemetriinställningar, förvärmning av miniatyrbilder, SEO-skydd, galleripapperskorg, automatiska uppdateringar och inställningar för Schemalagt underhåll. Inloggningsuppgifter, API-nyckelåtgärder, uppladdad profilering, rå CSS, flyttning av gallerilagring och destruktivt underhåll är endast information. Deras specialsidor öppnas inte från utkastet; lämna guiden innan du använder dessa arbetsflöden. Åsidosättningar per galleri lämnas oförändrade.

Slutgranskningen visar föreslagna ändringar först. Fäll ut oförändrade detaljer eller använd \ui{Visa alla inställningar} för att granska bevarade poster och informationsposter. Markera godkännandekryssrutan och tillämpa först efter att du granskat sammanfattningen. \ui{Avbryt} eller \ui{Börja om} kastar utkastet. Ett utkast tillhör den inloggade administratören och aktuell webbläsarsession; det är inget delat konfigurationsdokument.

Vid sparande valideras aktuella inställningar igen. Om en annan session har ändrat ett påverkat värde vägrar guiden skriva över det. Obligatorisk databaslagring måste verifieras; saknad eller okänd lagring blockerar det berörda sparandet. Program- och telemetrirader använder en transaktion, medan webbplatsadressändringar har en återställningsbar konfigurationsskrivning för felåterhämtning. Inställningar för Papperskorg och Schemalagt underhåll sparas genom sina befintliga ägare som sammanhängande grupper. Hemligheter placeras aldrig i utkastet, sammanfattningen eller guideloggen.

\subsubsection{Rätta webbplatsadressen}
\index{i00126@Inställningar!w00269@webbplatsadress}
\index{k00136@konfiguration!b00020@bas-URL}

Öppna \ui{Inställningar > Allmänt} för att rätta installationens offentliga URL efter installationen. Ange en absolut HTTP- eller HTTPS-adress, t.ex. \url{https://example.com}. Ta bara med en sökväg när galleriet faktiskt är installerat i den undermappen. Om webbhotellets identifiering lade till en oönskad mapp tar du bort den från fältet och sparar.

Sparandet ändrar endast strängen \codeword{base_url} på översta nivån i \filepath{config.php}; det skapar eller ändrar ingen databasinställning. Andra konfigurationsvärden, kommentarer och formatering bevaras. Både filen och mappen som innehåller den måste vara skrivbara. Programmet avvisar adresser med inloggningsuppgifter, frågeparametrar eller fragment och vägrar skriva om dubblerade eller beräknade konfigurationsvärden. Redigera en anpassad konfiguration manuellt när den inte kan uppdateras säkert.

Efter lyckat sparande öppnar webbläsaren samma inställningsavsnitt på den nya adressen. Efterföljande begäranden använder den uppdaterade konfigurationen. Inställningen rättar genererade adresser; den flyttar inte gallerifiler och konfigurerar inte DNS eller webbserverns dirigering.

\subsection{Administration av taggmetadata}
\index{t00241@taggar!a00002@administration}
\index{t00241@taggar!m00161@metadata}

\ui{Administration > Redigera taggar} är den centrala redigeraren för det återanvändbara taggordförrådet. Listan kan sorteras efter total användning eller alfabetiskt och visar varje taggs aktuella sammanlagda antal användningar i gallerier och fotografier. När en tagg väljs visas dess kanoniska namn med gemener, rena offentliga URL-namn och valfria offentliga beskrivning för taggens landningssida. Samma vy listar de gallerier och fotografier som för närvarande använder taggen, med länkar tillbaka till relevant offentligt sammanhang eller bildredigerare.

Att byta namn på en tagg eller ändra dess URL-namn uppdaterar den delade taggposten, så varje galleri- och bildkoppling fortsätter att referera till samma återanvändbara tagg. Namn och URL-namn normaliseras till den säkra taggkonventionen vid sparande, och ett dubblerat URL-namn avvisas. Att radera en tagg är annorlunda: programmet kopplar i en transaktion bort den från både gallerier och fotografier och tar sedan bort taggposten. Använd \ui{Visa offentlig tagg} för att granska landningssidan och \ui{Konfigurera taggvisning} för att öppna Temas kontroller för taggsidors, korts och huvudtaggars presentation.

\subsection{Funktionsväxlar}
\index{f00091@funktionsväxlar}
\index{f00088@Funktioner!a00002@administration}

\ui{Tema > Funktioner} ger globala växlar för programmets valfria delar. Registrerade funktioner är~aktiverade som standard så att en uppgradering inte tyst tar bort beteenden som en befintlig installation redan använder. En administratör kan stänga av enskilda besökarfunktioner, kartor och flygsimuleringsverktyg, AI- och uppladdningsintegrationer eller specialiserade administrationsverktyg när en enklare installation föredras.

En avstängd funktion döljs från vanlig navigering och dess skyddade rutter nekar funktionens arbetsflöde. Att stänga av en växel raderar inte funktionens kod, fotografier, databasposter, inloggningsuppgifter eller konfiguration. Återaktivering återställer därför vanliga ingångar, med samma kontroller av schema, åtkomst och konfiguration som tidigare. Exempel på växlingsbara funktioner är offentlig sökning, ljuslådans bläddringslägen, Bildhanteraren, Administration direkt på offentliga sidor, Smarta gallerier, röstning, Bildspelet, hämtningar, galleri- och flygkartor, navigeringsdata, SimBrief, OpenAI-hjälp, lokala AI-metadata, API-uppladdningar, mobil WebDAV, gallerimigrering, Mediefilnamnsbyte och anonym telemetri.

\subsection{Gallerihantering}

Galleriadministration omfattar samlingens hela liv: skapande, upptäckt, namngivning, beskrivning, datum, taggar, integritet, lösenord, profilering, omslag, layout, ordning, flyttning och slutlig radering. Ett undergalleri kan ärva användbara val från sitt övergalleri. Till exempel kan varje galleri i en privat familjegren följa övergalleriets skydd, och ett enskilt undergalleri kan använda egna presentationsval där redigeraren erbjuder en åsidosättning. Skärmen för gallerimassåtgärder ger motsvarande flergalleriåtgärder när samma ändring passar flera samlingar samtidigt.

Galleribeskrivningar kan visa säkra externa länkar skrivna som vanliga Markdown-länkar, \codeword{[link=URL]label[/link]}, \codeword{[url=URL]label[/url]} eller kortformer där den synliga texten också är målet. Endast HTTP- och HTTPS-mål blir länkar; andra scheman förblir overksam text. Kända tjänster använder medföljande lokala varumärkessymboler. För andra offentliga värdar kan sparande av galleriet hämta en liten validerad favicon till den lokala gallericachen. Hämtningen dedupliceras efter värdnamn, begränsas av storlek, tid, omdirigeringar och återförsök, avvisar privata eller reserverade nätverksmål och sker aldrig medan en besökare visar den offentliga beskrivningen. Om cachen eller dess databastabell är otillgänglig förblir länken användbar utan ikon.

När lagring för rena offentliga sökvägar är tillgänglig bygger \ui{Generera om rena offentliga sökvägar} om lagrade galleri- och bildsökvägar från aktuell hierarki och rapporterar separata antal för gallerier och bilder. Omordning och vissa sparningar i redigeraren uppdaterar också sökvägar när ändringarna kräver det. Underhållsåtgärden är användbar efter hierarkireparation, migrering eller äldre importer; den aktiverar inte URL-omskrivning av sig själv, så dirigering med frågesträngar förblir reservvägen när omskrivning är avstängd eller inte stöds av värden.
\subsection{Metadataorganiserare}
\index{m00162@metadataorganiserare}
\index{e00072@EXIF!d00055@datumorganiserare}
\index{g00096@galleri!a00018@automatisk organisering}

Galleriredigeraren innehåller ett arbetsflöde \ui{Organiserare} för att göra en platt samling direkt placerade fotografier till datumbaserade undergallerier. Den aktuella implementationen grupperar efter EXIF-tagningsdatum som redan är lagrade i bilddatabasen. Platsbaserad sekundär gruppering är reserverad för framtida användning och är inte en aktiv organiseringsstrategi i version~\version{}.

Välj tidigaste och senaste EXIF-datum och skapa sedan en förhandsvisning. Förhandsvisningen läser databasmetadata i små omgångar och skannar inte, döper inte om och flyttar inte originalfiler. Utkastet visar antalet direkt placerade fotografier, matchande kandidater, föreslagna undergallerier och fotografier som ignoreras eftersom de saknar användbart EXIF-datum eller ligger utanför vald period. Ett tidigaste datum är särskilt användbart för att utesluta fotografier från en kamera vars klocka inte var inställd.

För varje accepterat kalenderdatum är det föreslagna undergalleriets mappnamn det stabila ISO-datumet \codeword{YYYY-MM-DD}. Visningstiteln använder galleriets formatterare för datumvisning och faller tillbaka på ett lättläst dag/månad/år-värde om formatteraren inte kan ge en etikett. Innan ett nytt undergalleri föreslås söker organiseraren efter ett befintligt direkt undergalleri med samma organiserartitel och återanvänder det. Därmed lägger upprepade förhandsvisningar och tillämpningar till material i samma dagsdestination i stället för att skapa ännu ett likvärdigt dagsgalleri.

Att tillämpa utkastet kräver en separat bekräftelse. PHP Gallery skapar eller återanvänder de föreslagna datumundergallerierna och flyttar matchande fotografier genom samma fysiska flyttjänst som den vanliga massflyttaren för bilder använder. Original och kända genererade derivat flyttas därför tillsammans med uppdateringen av databastillhörigheten. Stora tillämpningar delas i begränsade webbläsardrivna begäranden i stället för en lång PHP-begäran. Granska utkastet och ha en säkerhetskopia innan du organiserar om ett stort galleri.

\subsection{Mediefilnamnsbyte}
\index{m00159@Mediefilnamnsbyte}
\index{f00078@filnamn!n00172@namnbyte}
\index{m00160@medieunderhåll}

\ui{Gallerier > Mediefilnamnsbyte}, och panelen för filnamnsbyte som finns för enskilda gallerier, kan normalisera fysiska bildfilnamn utifrån gallerisammanhang och aktuell bildordning. Det är en filsystemsåtgärd, inte en kosmetisk titeländring, så verktyget börjar alltid med en provkörningsplan. Förhandsvisningen listar varje gammalt och föreslaget filnamn och markerar filer som redan stämmer, saknas, skulle krocka eller måste hoppas över enligt en säkerhetsregel.

Endast bildfiler direkt i den valda gallerimappen kan behandlas. Saknade original, sökvägar utanför galleriet och osäkra mål döps inte om. Om det föredragna målnamnet redan är upptaget men ett säkert deterministiskt suffix kan lösa kollisionen visar förhandsvisningen justeringen. Att tillämpa en plan kräver en bekräftelse på att förhandsvisningen har granskats.

Standardmönstret är \codeword{\{gallery_path\}_\{seq4\}}. Ett eget mönster kan använda \codeword{\{gallery_path\}}, \codeword{\{gallery_title\}}, \codeword{\{photo_title\}}, \codeword{\{old_name\}}, \codeword{\{old_stem\}}, \codeword{\{image_id\}} och sekvensplatshållare från \codeword{\{seq\}} till \codeword{\{seq5\}}. De numrerade sekvensvarianterna fyller ut till två, tre, fyra eller fem siffror. Textkomponenter normaliseras till filsystemssäkra ASCII-liknande namn med gemener och understreck, originalets filändelse bevaras i rensad form och ett tomt eget mönster faller tillbaka på standarden.

Efter ett lyckat fysiskt namnbyte håller PHP Gallery databasraden och referenser till offentliga sökvägar samordnade, uppdaterar härledda titlar där verktyget ansvarar för dem, flyttar eller ogiltigförklarar genererade derivat efter behov och tar bort inaktuella cacheposter för ZIP-arkiv som refererar till originalets gamla identitet. Varningar vid derivatrensning motiverar inte att ett annars slutfört namnbyte av original återställs, eftersom återskapbara derivat kan genereras igen. För webbplatsövergripande urval driver webbläsaren arbetet framåt i begränsade delar och visar resultatet för varje galleri.

\subsection{Bildhantering}
\label{sec:media}
\index{b00031@bilder}
\index{u00254@uppladdningar}

Fotografier kan laddas upp i webbläsaren, kopieras till mappar och upptäckas, tas emot genom konfigurerad automatisering, överföras från en annan kompatibel installation eller läggas till genom ett avgränsat mobilt arbetsflöde. När ett fotografi har kommit in registrerar galleriet dess storlek och tillgängliga kamerauppgifter och förbereder de mindre visningskopior som används på hela webbplatsen.

Den enskilda bildredigeraren omfattar för närvarande titel, beskrivning, underhållna språkvarianter, synlighet, sorteringsordning, taggar, NSFW / 18+-status, det privata redaktionella betyget som används av regler för smarta gallerier, valfria responsiva miniatyrbildsgränser per fotografi och skrivskyddade EXIF/GPS-uppgifter. Där lokala AI-metadata finns visar redigeraren genererad modell och källa, sökbar intern text och rå metadata-JSON separat från den offentliga bildtexten. OpenAI-kontroller infogar granskningsbara utkast i redigerbara textfält och ersätter inte sparad redaktionell text automatiskt.

Massåtgärder för bilder är avsiktligt mer begränsade än den enskilda redigeraren och kontrollerar tillhörigheten igen på servern. Aktuella massåtgärder omfattar:

\begin{itemize}
  \item flytta valda fotografier till ett befintligt galleri;
  \item skapa ett nytt målgalleri och flytta urvalet dit, med rensning av ett tomt nyskapat mål om flytten misslyckas innan filer har lagrats slutgiltigt;
  \item radera valda bildposter tillsammans med deras tillhörande original och kända genererade derivat genom den vanliga raderingstjänsten;
  \item ange ett valt fotografi som galleriets titel- eller omslagsbild;
  \item ändra bildsynlighet mellan de stödda tillstånden utkast, offentligt och privat;
  \item aktivera eller stänga av NSFW-skydd per fotografi när dess schema är känt;
  \item skapa miniatyrbilder för valda fotografier.
\end{itemize}

Bildordningen kan ändras genom att dra i radhandtaget, och galleriet kan sorteras efter filnamn eller EXIF-fotodatum. Foton utan fotodatum ligger kvar sist; ändringar sparas genom den vanliga ordningsrutten.

Fliken \ui{Bilder} i galleriredigeraren samlar vanliga bildåtgärder i en kompakt lista. Tillgängliga miniatyrförhandsvisningar kan öppnas med pekaren, tangentbordsfokus eller klick; radmarkering stöder intervall med Skift-klick. Varje rad har direktkontroller för synlighet, val av galleriets titelbild, redigering och radering av fotot. Åtgärderna använder fortfarande de befintliga autentiserade arbetsflödena på servern. \ui{Visa filnamn} ändrar bara redigerarens visning; administratörer kan spara ett valfritt webbläsarstandardvärde, medan varje galleri kommer ihåg valet under den aktuella redigerarsessionen.

När den valfria funktionen \ui{Bildhanteraren} är~aktiverad kan en vanlig inloggad administratör välja fotografier och direkta fysiska undergallerier i det offentliga gallerirutnätet. Bockar väljer enskilda kort; Skift-klick väljer ett intervall; Ctrl-klick eller Cmd-klick växlar ett kort; \ui{Markera alla}, \ui{Rensa} och Ctrl/Cmd+A stöder större urval på den synliga sidan. Smarta gallerikort ingår inte i detta filbaserade urval. Att dra ett omarkerat fotografi börjar med det fotografiet som urval, medan dragning av ett befintligt bildurval flyttar de valda fotografierna tillsammans. Ett synligt undergallerikort blir ett släppmål, och släpp där utför samma validerade bildflytt som destinationsväljaren. Fysiska gallerikort kan inte själva dras.

Från verktygsfältet kan administratören radera ett blandat urval eller skapa ett riktigt fysiskt galleri från det. Skapandet accepterar ett sökbart val av övergalleri, kopierar valda fotografier till den nya roten och kopierar valda fysiska undergalleriträd under den samtidigt som källorna lämnas oförändrade. Servern skriver aktuella galleri-sidofiler, avvisar rekursiva mål och befintliga målmappar, reserverar varje målrot före kopiering, bygger om kopierade gallerirader från filsystemet och tar bort det nya partiella resultatet om åtgärden misslyckas. Blandad radering tar bort valda fotografier genom vanlig bildrensning och tillämpar den konfigurerade återställningsbara papperskorgspolicyn på valda galleriträd när den är tillgänglig.

Flyttning och kopiering till befintligt galleri är fortfarande åtgärder enbart för fotografier. Om ett fysiskt galleri är valt avmarkerar du det innan du använder dessa kontroller. Bildurval kan också delas eller hämtas. \ui{Dela} hämtar de behöriga bildversioner som webbläsaren kan visa som \codeword{File}-objekt i webbläsaren och öppnar den inbyggda Web Share-panelen när webbläsaren och enheten stöder delning av flera filer. Webbläsaren väljer vilka mottagande appar som är tillgängliga, så PHP Gallery kan inte tvinga fram Instagram Story, Reel eller något annat särskilt mål. Om inbyggd delning saknas, inte kan ta emot filerna eller misslyckas efter förberedelsen startar hanteraren den autentiserade ZIP-hämtningen för valda fotografier som reserv. Samma ZIP-adress kan därför användas både för uttrycklig hämtning av valda fotografier och som reserv för delning. Vanliga besökare får aldrig dessa kontroller, och servern kontrollerar på nytt det inloggade kontot, CSRF-token, källtillhörighet, direkta undergalleriers tillhörighet och målgalleri för varje ändring eller ZIP-begäran.

\subsection{Uppladdningar och upptäckt}
\index{u00258@upptäckt}

Filsystemsupptäckt är användbar när fotografier redan har kopierats till servern via FTP, webbhotellets filhanterare eller en återställd säkerhetskopia. Upptäcktsskärmen visar vad programmet hittade och låter administratören föra in mapparna och filerna i katalogen. Upptäckt körs som ett sessionsbaserat webbläsarjobb i små AJAX-omgångar i stället för att skanna hela trädet under den första begäran för administrationssidan. Aktuell standard undersöker upp till 80 mappar per omgång, accepterar en begränsad omgångsstorlek upp till 300 och låter övergivna upptäcktsjobb löpa ut efter två timmar. Det håller stora mappträd responsiva på långsammare webbhotell och visar ändå förlopp och en slutlig resultatlista som går att arbeta med.

Skannern går under den konfigurerade galleriroten och följer inte symboliska länkar till mappar. Dolda mappar och kända genererade eller cachelagrade mappnamn som \filepath{cache}, \filepath{thumbs}, \filepath{thumbnail}, \filepath{thumbnails}, \filepath{preview} och \filepath{previews} ignoreras. En ny kandidatgren måste innehålla minst en bild som stöds någonstans under sig. Tomma strukturella övermappar kan ändå bli gallerier när de behövs för att bevara hierarkin ovanför undergallerier som innehåller bilder. En mapp som bara innehåller en \filepath{gallery.json}-sidofil och ingen stödd bild någonstans i grenen rapporteras som enbart metadata i stället för att importeras som galleri i det aktuella upptäcktsflödet. Redan indexerade gallerimappar tas bort från kandidatlistan, och en importkandidat vars visningstitel skulle dubblera en befintlig syskontitel skapas inte tyst.

För varje ohanterat resultat kan administratören importera mappträdet på plats, flytta stödda fotografier från valda ohanterade mappar till ett befintligt galleri eller radera valda ohanterade mappar från disken. Importen utökar urvalet i ordning från övergalleri till undergalleri och kan inkludera saknade över- och undergallerier som innehåller bilder. Flyttåtgärden undviker att skriva över ett befintligt målfilnamn genom att välja ett deterministiskt unikt namn, skannar målgalleriet efter lyckade flyttar och tar bort källmappar endast när de blir tomma. Upptäcktens särskilda diskradering är avsiktligt begränsad till säkra ohanterade sökvägar under galleriroten som inte är roten själv och nekar mappar som redan representeras av databasgallerier. Använd det vanliga galleriraderingsflödet för ett importerat CMS-galleri.

\subsubsection{Metadata-sidofilen gallery.json}
\index{g00114@gallery.json}
\index{s00218@sidofilsmetadata}

PHP Gallery håller sin aktiva katalog och sina relationer i databasen, medan en valfri \filepath{gallery.json}-fil intill en gallerimapp bevarar användbara redigerbara metadata med filsystemet. Vanliga gallerisparningar skriver tillbaka sidofilen till gallerimappen. Aktuell sidofilsutmatning omfattar titel, beskrivning, normaliserade galleritaggar, synlighet, ordning, röstnings- och filnamnsvisningstillstånd, underhållet innehållsspråk och översättningar när de är tillgängliga, manuella galleridatum, valda åsidosättningar för kort, antal, ljuslåda och rutnät, responsiva miniatyrbildsgränser, åtkomst- och listningspolicy, omslagsreferenser samt konfigurerade profileringsresurser för galleriet där dessa scheman är tillgängliga. Sidofilsskrivning följer samma schemakontroller som motsvarande gallerifunktioner.

Vid upptäckt och reparation av saknade övergallerier kan PHP Gallery läsa stödda värden från sidofilen i stället för att härleda allt från mappnamnet. Importerade rader får försiktigt standardvärdet opublicerat när ingen synlighet anges. Den aktuella importvägen återställer de sidofilsfält som importören uttryckligen stöder, inklusive titel, beskrivning och lokalisering, synlighet, manuella datum, röstnings- och filnamnspresentation, valda åsidosättningar för kort, antal, ljuslåda, rutnät och miniatyrbilder, åtkomst- och listningspolicy, ordning och profileringssökvägar. Sidofilen bör därför behandlas som flyttbara gallerimappsmetadata, inte som en fullständig databassäkerhetskopia: relationstillstånd, inloggningsuppgifter, loggar, röster, speltillstånd, definitioner för smarta gallerier och andra databasägda poster kräver fortfarande den vanliga säkerhetskopieringsrutinen.

\securitynote{Åtgärden \ui{Radera från disk} på upptäcktsskärmen är en riktig rekursiv filsystemsradering av ohanterade kandidatmappar. Granska de valda relativa sökvägarna innan du bekräftar. Skyddet mot kända importerade gallerisökvägar är en ytterligare säkerhetsgräns, inte en ersättning för säkerhetskopiering.}

Webbläsaruppladdning är det mest lättanvända valet för vanligt arbete. Den kan ladda upp till ett befintligt galleri eller skapa ett undergalleri och ladda upp valfria fotografier i samma sidopanelflöde. Välj destination, välj filer och följ förloppsvisningen. Programmet kontrollerar destinationen och filinformationen innan varje resultat accepteras.

Enbart komprimerad filstorlek bevisar inte att en bild är säker att avkoda. Innan GD eller valfri Imagick allokerar ett helt raster kontrollerar PHP Gallery format, mått, pixelberäkningar, konfigurerad minnesbudget och körningsminnesbudget samt pipelinegränser. En fil som inte kan behandlas säkert nekas innan dess original eller genererade derivat lagras slutgiltigt i målgalleriet. Policyn gäller klassisk och webbläsarförberedd inläsning; den påstår inte att varje extern RAW-avkodare har samma minnesmodell i processen.

Uppladdningsinställningarna skiljer den vanliga servervägen från webbläsarstödd förberedelse. \ui{Tillåt alla format som servern stöder} håller väljaren i linje med serverns upptäckta funktioner. \ui{Föredra telefonrenderad JPG/PNG/WebP, ingen RAW/DNG} ber mobilwebbläsare om en renderad bildrepresentation när det går, vilket är användbart när ett iPhone ProRAW/DNG-flöde ger olämpliga färger. Webbläsarens anvisning garanterar ingen omkodning, så servern validerar fortfarande de medier som faktiskt laddas upp.

Webbläsarstödd förberedelse aktiveras oberoende och är markerad som standard i uppladdningsflödet. Webbläsaren kan förbereda optimerade miniatyrbilder och begränsade ZIP-omgångar innan de skickas. Om webbläsarförberedelsen misslyckas innan någon skrivning på servern börjar kan vanliga bildurval automatiskt falla tillbaka på klassisk serveruppladdning. Administratörer kan begränsa standard- och maxantalet webbläsararbetare, bilder per omgång, förhållandet mellan ZIP-storlek och PHP-uppladdningsgräns, högsta byteantal för ZIP-omgångar och originaldelstorleken som används vid webbläsarbaserad ombyggnad av miniatyrbilder. Normaliseraren tillämpar också ett hårt tak för arbetare så att en oavsiktlig inställning inte kan skapa obegränsade parallella begäranden på delade webbhotell.

Aktuella standardvärden för webbläsarpipelinen är 8 arbetare, ett mål på 80~\% av den upptäckta PHP-uppladdningsgränsen för varje ZIP-omgång utan komprimering, högst 8 bilder per ZIP-omgång och ett absolut tak på 24~MiB per ZIP-omgång. De validerade intervallen är 1--32 arbetare, 10--95~\% av PHP-uppladdningsgränsen, 1--64 bilder per omgång och 1--128~MiB för det absoluta ZIP-taket. Webbläsarbaserad ombyggnad av miniatyrbilder använder som standard en original-ZIP-del på 512~MiB, begränsad till 16~MiB--3~GiB, och originalbyggaren begränsar normalt en del till 96 objekt med en intern hård gräns på 512. Detta är säkerhetstak och standardvärden, inte en rekommendation att höja ett långsamt delat webbhotell till maximum.

När webbläsarstödd uppladdning är~aktiverad accepterar väljaren även ZIP-arkiv, t.ex. en export från iCloud Photos. Webbläsaren undersöker arkivet lokalt, lägger till stödda JPEG-, PNG-, WebP- och GIF-bilder i den vanliga förberedelsekön och hoppar över mappar, dolda metadata och poster som inte stöds. Den valda ZIP-filen laddas aldrig själv upp till PHP. Arkiv utan komprimering och med Deflate stöds; krypterade arkiv, flerdiskarkiv, ZIP64-arkiv, skadade arkiv, misstänkt komprimerade arkiv eller alltför stora arkiv avvisas innan någon serveruppladdning börjar. ZIP-import har ingen klassisk PHP-reservväg, så låt webbläsarstödd förberedelse vara aktiverad när du väljer ett arkiv.

Det globala alternativet för automatiskt filnamnsbyte vid uppladdning kan köra samma namngivningspolicy efter skanning som Mediefilnamnsbyte använder, efter inläsning genom webbläsare, API eller WebDAV. Låt det vara avstängt när externa filnamn ingår i din arkivkonvention; aktivera det när alla ankomstvägar ska samlas kring galleriets normaliserade namnschema.
\subsection{Genererade metadata och automatisering}
\index{a00007@AI-metadata}
\index{o00180@OpenAI}
\index{u00255@uppladdningsautomatisering}
\index{w00271@WebDAV}

Den valfria lokala eller arbetarbaserade bildanalyskön kan registrera visuella metadata separat från mänskligt skrivna titlar, beskrivningar, taggar och EXIF-fält. Administratörer kan granska dessa metadata i bildredigeraren och använda dem som sökbart underlag utan att redaktionell text ersätts tyst. Konfigurerad OpenAI-hjälp kan generera en bild- eller galleribeskrivning, förbättra befintlig text, sammanfatta sammanhang, beskriva synligt innehåll från säkra genererade miniatyrbilder eller föreslå en översättning. Åtgärderna infogar granskningsbara utkast; de publicerar eller sparar inte text automatiskt, utan administratören måste slutföra den vanliga sparåtgärden.

För återkommande överföringar skapar du en avgränsad token för uppladdningsautomatisering till det avsedda galleriet och skickar stödda filer till uppladdningsadressen. Token kan återkallas oberoende av varandra och måste behandlas som lösenord. Den råa API-nyckeln visas bara när den skapas; därefter lagras endast dess hash, så kopiera den nya nyckeln till den avsedda klienten vid skapandet. Den globala \ui{API-hanteraren} listar automatiseringsuppgifter över gallerier och ger ett centralt återkallningsflöde utan att exponera ursprungliga hemligheter. Gallerikontroller finns kvar för samma avgränsade uppgifter.

Det kompletterande Windows-programmet använder samma galleribegränsade uppladdningskontrakt. Det har två vanliga uppladdningslägen. \ui{Bevaka mapp} ignorerar filer som redan fanns när bevakningen börjar, väntar på att nya filer blir stabila, kommer ihåg bekräftade uppladdningar lokalt och försöker igen vid tillfälliga fel med ökande väntetid. Det kan vid behov bifoga aktuell latitud, longitud och höjd för kameran eller flygplanet i Microsoft Flight Simulator genom SimConnect omedelbart före varje bevakad uppladdning och radera ett bevakat original först efter att galleriet bekräftat en lyckad uppladdning. \ui{Manuell uppladdning} accepterar ett filurval oberoende av den bevakade mappen och kan antingen generera responsiva miniatyrbilder på datorn med arbetarprocesser eller be galleriservern att generera dem. Miniatyrbildsarbete och nätverksuppladdningar överlappar i en pipeline i stället för att först förbereda ett helt stort urval. Stöd för meddelandefältet låter bevakaren eller det manuella jobbet fortsätta medan huvudfönstret är dolt.

Samma Windows-program kan vid behov köra en lågprioriterad AI-metadataarbetare. Den tar ett autentiserat analysjobb, hämtar bara den tilldelade resursen, skapar interna sökbara metadata med vald lokal bakände och modellgeneration och skickar tillbaka resultatet till den befintliga adressen för uppladdningsautomatisering. Arbetaren är åtskild från mänskliga bildtexter. Galleriets åtgärd \ui{Tvinga omgenerering av AI-metadata} rensar lagrade genererade rader och gammalt kötillstånd för den gallerigrenen och köar nytt arbete; den tunga analysen körs fortfarande på Windows-arbetaren.

Mobila WebDAV-uppgifter är galleribegränsade och använder HTTP Basic-autentisering samt en sökvägstoken som inte går att gissa. När en anslutning skapas visas dess genererade lösenord en gång; servern behåller sedan endast en lösenordshash. Administrationssidan listar anslutningsetikett, målgalleri, server-URL och tid för senaste användning och kan återkalla anslutningen. WebDAV-adressen stöder det minimala beteendet för \codeword{OPTIONS}, \codeword{PROPFIND}, \codeword{MKCOL} och \codeword{PUT} som behövs för bildöverföringsklienter som PhotoSync. Aktuell WebDAV-inläsning accepterar JPG, PNG, GIF och WebP. Konfigurera mobilklienten att konvertera HEIC till JPEG före överföring där det är möjligt. Uppladdade medier går sedan genom samma skanning, valfria automatiska filnamnsbyte, databas- och miniatyrbildspipeline som andra uppladdningar.
\section{Generering av miniatyrbilder och offentlig rendering}
\label{sec:thumbnails}
\index{m00165@miniatyrbilder}
\index{w00272@WebP}
\index{j00130@JPEG}
\index{b00032@bildförhandsvisningar}
\index{s00217@sidhastighet!m00165@miniatyrbilder}

\subsection{Vad en miniatyrbild är}
\index{m00164@miniatyrbild!d00057@definition}

En miniatyrbild är en mindre visningskopia av ett fotografi. Originalfotografiet kan vara flera tusen pixlar brett och uppta många megabyte. Ett gallerikort på en telefon kanske bara behöver några hundra pixlar. Att skicka hela originalet för varje litet kort skulle låta besökarna vänta på detaljer som de inte kan se i den storleken.\footnote{Miniatyrbildsgenerering behåller originalfotografiet som arkivresurs och för visning i full storlek. Den genererade kopian är ett utbytbart bläddringshjälpmedel.}

Föreställ dig ett galleri med 60 fotografier, vart och ett med en originalfil på 12~megabyte. Att läsa in varje original för det första rutnätet kan kräva hundratals megabyte. Små miniatyrbilder kan minska den första visningens belastning kraftigt. Originalet förblir tillgängligt för behörig fullstor visning eller hämtning, medan rutnätet får effektiva förhandsvisningar.\footnote{Den verkliga besparingen beror på bildinnehåll, kvalitet, format, mått och den kandidat som webbläsaren väljer. Exemplet visar storleksordningen snarare än att lova ett fast komprimeringsförhållande.}

Miniatyrbilder fyller flera vardagliga syften:

\begin{itemize}
  \item gallerirutnät visas snabbare;
  \item rullning använder mindre minne och nätverkskapacitet;
  \item mobilbesökare får en bild som ligger närmare skärmens storlek;
  \item galleriomslag och sökresultat förblir kompakta;
  \item den stora visaren kan börja med en lämplig förhandsvisning medan en större källa förbereds;
  \item återbesök kan återanvända filer som redan ligger i webbläsarens cache.\footnote{Webbläsarcachen lagrar nyligen använda webbresurser på besökarens enhet enligt serverns och webbläsarens policy. Återanvändning kan göra ett senare besök mycket snabbare.}
\end{itemize}

\subsection{Varför flera miniatyrbildsstorlekar finns}
\index{m00164@miniatyrbild!s00233@storlekar}
\index{e00070@enhetens pixeltäthet}

En enda liten storlek passar inte varje skärm. Ett smalt telefonkort kan se tydligt ut med en miniatyrbild på 300~pixlar. Ett brett datorkort kan behöva 600 eller 800~pixlar. En skärm med hög pixeltäthet använder fler bildpixlar för att rita samma fysiska yta och kan ha nytta av en större kandidat.

PHP Gallery kan förbereda vanliga storlekar som 300, 600, 800, 960, 1280 och 1600~pixlar, beroende på konfiguration och originalbildens begränsningar. Talet beskriver den genererade kopians maximala långsida. Stående och liggande fotografier behåller sina ursprungliga proportioner.\footnote{För ett liggande fotografi är långsidan vanligen bredden. För ett stående fotografi är det vanligen höjden. Den kortare sidan beräknas från originalets bildförhållande.}

De mindre kandidaterna används för rutnät och omslag. Större kandidater används för breda kort, skärmar med hög pixeltäthet, sökvisningar, kartförhandsvisningar och den stora visaren. Webbläsaren kan välja ur den tillgängliga gruppen i stället för att få en fast storlek för varje situation.

\subsection{Varför JPEG och WebP är tillgängliga}
\index{m00164@miniatyrbild!f00083@format}
\index{j00130@JPEG!a00012@användning för miniatyrbilder}
\index{w00272@WebP!a00012@användning för miniatyrbilder}

JPEG har bred kompatibilitet och är välkänt för fotografier. WebP ger ofta liknande visuell kvalitet med färre överförda byte. PHP Gallery kan generera båda så att en kompatibel webbläsare kan använda WebP medan JPEG förblir ett tillförlitligt alternativ.\footnote{Komprimeringseffektiviteten varierar med fotografiet och kvalitetsinställningarna. Fin struktur, brus, toningar och skarpa grafiska detaljer kan ge olika resultat.}

Besökaren väljer inte format manuellt. Sidan talar om för webbläsaren vilka versioner som finns, och webbläsaren väljer ett format den förstår. Originalfotografiet behåller sitt källformat och är åtskilt från dessa bläddringskopior.

\subsection{När miniatyrbilder skapas}
\index{m00164@miniatyrbild!g00115@generering}
\index{m00164@miniatyrbild!u00251@underhåll}

Miniatyrbilder kan förberedas när fotografier importeras eller laddas upp, byggas om genom administrationsunderhåll eller genereras genom en skyddad reparationsprocess när den offentliga sidan upptäcker att en förväntad kopia saknas. Att förbereda dem nära uppladdningstillfället ger senare besökare den smidigaste upplevelsen eftersom de nödvändiga filerna redan finns före första besöket.\footnote{Stora importer kan kräva betydande bearbetningstid och lagring. Att köra förberedelsen som en avsiktlig administrationsuppgift ger ägaren en tydligare möjlighet att övervaka slutförandet.}

En galleriägare bör granska miniatyrbildsunderhåll efter att ha:

\begin{itemize}
  \item importerat ett stort befintligt mappträd;
  \item ändrat aktiverade miniatyrbildsstorlekar eller kvalitet;
  \item återställt en säkerhetskopia med original och databasdata men ofullständig miniatyrbildscache;
  \item flyttat en installation mellan servrar med olika bildformatstöd;
  \item sett upprepade återgångar till originalfiler eller saknade förhandsvisningar i Systemstatus.
\end{itemize}

Genererade miniatyrbilder är återskapbara cachedata och kan byggas om från originalfotografierna. Själva originalen behöver permanent säkerhetskopieringsskydd.

\subsection{Visningsmaster för DNG och RAW}
\index{d00064@DNG}
\index{r00199@RAW-fotografier}

DNG-fotografier kan förbli arkivoriginal medan PHP Gallery skapar en visningsmaster för webbläsaren för miniatyrbildsgenerering och offentlig visning. Visningsmastern är genererade derivatdata, inte en ersättning för den ursprungliga DNG-filen. Om det nödvändiga lokala konverteringsverktyget saknas eller konverteringen misslyckas identifierar administrationens uppladdnings- och miniatyrbildsrapporter det misslyckade DNG-derivatet. Miniatyrbildsunderhåll kan bygga om mastern senare; radering av genererade derivat raderar inte original-DNG:n.

\ui{Körningsdiagnostik} visar den aktiva DNG-konverteringspolicyn tillsammans med serverns faktiska funktioner för GD, Imagick/ImageMagick och format. Källpolicyn kan använda automatisk reserv med full RAW-avkodning först och inbäddad kameraförhandsvisning därefter, uttryckligen föredra full RAW-avkodning eller föredra den inbäddade förhandsvisningen. Färghanteringen kan tvinga webbläsarsäker sRGB, bevara det avkodade utseendet så nära som möjligt eller använda kamerans vitbalans när den är tillgänglig. Valen påverkar endast genererade webbläsarderivat. De skriver aldrig om arkivets DNG-original, och tillgängligheten beror fortfarande på de bildhanterare och minnesgränser som webbhotellet tillhandahåller.

\subsection{Vad besökarna får}
\index{m00164@miniatyrbild!w00268@webbläsarval}

Galleriet skickar webbläsaren en lista över miniatyrbildskopior som faktiskt finns för ett kort. Webbläsaren tar hänsyn till tillgängligt utrymme på sidan, skärmens skärpebehov, det stödda formatet och sina egna inläsningsbeslut. Den begär sedan en lämplig fil.

Anta till exempel att ett fotografi har kopior på 300, 600, 800 och 960~pixlar:

\begin{itemize}
  \item ett litet telefonkort kan få kopian på 300~pixlar;
  \item ett medelstort kort kan få kopian på 600~pixlar;
  \item ett stort kort eller ett kort på en skärm med hög pixeltäthet kan få kopian på 800 eller 960~pixlar;
  \item den stora visaren kan begära en ännu större förhandsvisning som förberetts för ändamålet.
\end{itemize}

Det exakta valet kan skilja sig mellan webbläsare och skärmar. Flexibiliteten är användbar eftersom servern inte behöver gissa en universell storlek.

\subsection{Responsivt webbläsarval}
\index{r00205@responsiv renderare}
\index{s00232@srcset}
\index{m00164@miniatyrbild!r00206@responsivt läge}

\ui{Responsivt webbläsarval} är det stödda äldre läget under \ui{Administration > Tema > Layout}. Sidan visar tillgängliga miniatyrbildskandidater omedelbart och webbläsaren väljer den som bäst matchar det renderade kortet och skärmens pixeltäthet. JavaScript behövs inte för detta val.\footnote{Den bakomliggande webbstandarden kallar kandidatlistan \emph{srcset}. Galleriägare kan förlita sig på det visuella beteendet utan att lära sig denna märkningsterm.}

Responsivt val överför vanligen en vald miniatyrbild för varje inläst kort och är det vanliga valet när förutsägbart inbyggt webbläsarbeteende är viktigast.

\begin{tabularx}{\textwidth}{>{\RaggedRight\arraybackslash}p{0.30\textwidth} X}
\toprule
\textbf{Situation} & \textbf{Sannolikt resultat}\\
\midrule
Smalt mobilkort & En mindre miniatyrbildskandidat räcker vanligen.\\
Brett datorkort & Webbläsaren kan välja en medelstor eller större kandidat.\\
Skärm med hög pixeltäthet & En kandidat med högre upplösning kan väljas för samma CSS-storlek.\\
JavaScript saknas & Responsivt kandidatval fungerar fortfarande.\\
\bottomrule
\end{tabularx}

Kopian på 300~pixlar fungerar som reserv när den finns. Dess närvaro tvingar inte en modern webbläsare att hämta den först eftersom de större alternativen presenteras samtidigt.

\subsection{Progressiv skärpning av miniatyrbilder}
\index{p00193@progressiv renderare}
\index{p00194@progressiv skärpning}
\index{m00164@miniatyrbild!p00195@progressivt läge}

\ui{Progressiv skärpning av miniatyrbilder} är standard och säker reserv under \ui{Administration > Tema > Layout}. Galleriet visar först en liten miniatyrbild, vanligen 300~pixlar, och ersätter senare relevanta kort med större kopior. På långsammare anslutningar kan detta göra samlingen igenkännbar tidigare, till priset av att både den lilla och den ersättande miniatyrbilden kan överföras.

\subsubsection{Vad som händer på skärmen}

\begin{enumerate}
  \item Sidan reserverar kortets slutliga geometri.
  \item Små miniatyrbilder börjar fylla synliga kort.
  \item Relevanta kort köas för en större kandidat.
  \item Den större kopian läses in medan den lilla förblir synlig.
  \item Kortet ersätts när den större kopian är klar.
\end{enumerate}

Större uppgraderingar köas i stället för att startas för varje kort samtidigt, och synliga kort prioriteras framför avlägset innehåll. Det aktuella progressiva alternativet gäller bildkort i ett öppet galleri. Galleriomslag, kollagerutor, gallerikort på startsidan och administrationsbilder fortsätter använda responsivt webbläsarval. Administrationsgränssnittet märker progressiv skärpning som standard och responsivt val som det stödda äldre alternativet.

\subsubsection{Avvägningen i överföring}

Progressivt läge kan överföra två filer för ett fotografi: den första lilla kopian och den skarpare ersättningen. Responsivt läge överför ofta en kopia vald av webbläsaren. Progressivt läge prioriterar därför tidigare visuell fyllning, medan responsivt läge vanligen minskar total överföring för samma kort.

\subsection{Vilket läge bör ägaren välja?}
\index{m00164@miniatyrbild!v00265@välja renderare}

Välj \ui{Responsivt webbläsarval} när effektivt inbyggt webbläsarval och vanligen en miniatyrbildsöverföring per inläst kort passar samlingen.

Behåll \ui{Progressiv skärpning av miniatyrbilder} för en liten första bild med efterföljande begränsade uppgraderingar av relevanta kort. Testa båda stödda renderarna med representativa gallerier och enheter när du väljer rätt presentation för besökarna.

Användbara frågor:

\begin{itemize}
  \item Ser besökarna de första fotografierna tidigare?
  \item Ser den lilla versionen acceptabel ut under skärpningen?
  \item Hur mycket extra data överförs vid ett typiskt besök?
  \item Skärps bilder på telefoner och skärmar med hög pixeltäthet i en behaglig takt?
  \item Passar det valda läget målgruppens sannolika anslutningskvalitet?
\end{itemize}

\subsection{Inläsningsprioriteringar och stabil layout}
\index{l00149@lat inläsning}
\index{l00150@layoutstabilitet}
\index{m00164@miniatyrbild!l00149@lat inläsning}

De första synliga fotografierna får tidigare uppmärksamhet vid inläsningen eftersom de formar besökarens första intryck. Fotografier längre ned på sidan använder lat inläsning, vilket innebär att webbläsaren kan vänta tills de närmar sig det synliga området. Det undviker att lägga anslutningskapacitet på slutet av ett långt galleri medan besökaren fortfarande tittar högst upp.\footnote{Tidpunkten för lat inläsning styrs delvis av webbläsaren. Webbläsare tar hänsyn till rullning, avstånd till visningsområdet, anslutningstillstånd och sin egen implementationspolicy.}

Galleriet känner till indexerade fotografiers bredd och höjd. Det använder proportionerna för att reservera stabila kortformer innan bildinläsningen är klar. Text, knappar och närliggande kort behåller därför sina positioner när miniatyrbilder visas.

Besökare som föredrar minskad rörelse får en lugn presentation utan obligatorisk kontinuerlig inläsningsanimation. Bildlänkar går att använda under hela inläsningen och skärpningen.

\subsection{Miniatyrbildskvalitet och lagring}
\index{m00164@miniatyrbild!k00142@kvalitet}
\index{m00164@miniatyrbild!l00146@lagring}

Högre kvalitet bevarar fler fina detaljer och kan skapa större filer. Lägre kvalitet sparar utrymme och överföring men ger mer synlig komprimering. Den bästa inställningen beror på motiv, skärmanvändning, webbhotellskapacitet och målgrupp. Detaljerat lövverk, hår, text och fin arkitektur kan avslöja komprimering tydligare än stora mjuka bakgrunder.

Varje ytterligare aktiverad storlek och varje format använder lagring eftersom de skapar ännu en kopia för varje fotografi. Fördelen är mer precis leverans till olika skärmar. Systemstatus och lagringsstatistik hjälper ägaren att förstå den resulterande cachen. Ett stort originalarkiv kan ge en omfattande miniatyrbildssamling, men filerna går fortfarande att återskapa från originalen.

\subsection{Kompatibilitet och underhållslägen för miniatyrbilder}
\index{m00166@miniatyrbildskompatibilitet}
\index{m00167@miniatyrbildsunderhåll}
\index{w00266@webbläsarbaserad ombyggnad av miniatyrbilder}

Underhållspanelen har två kompatibilitetspolicyer. \ui{Modern} genererar WebP-miniatyrbilder där servern kan skriva WebP. \ui{Äldre kompatibilitet} genererar JPEG-reservminiatyrbilder utöver WebP. När en installation permanent har växlats till Modern raderar \ui{Ta bort äldre JPG-miniatyrbilder} endast genererade JPEG-miniatyrbildsfiler; original, WebP-derivat och DNG-visningsmaster bevaras.

En fullständig körning av \ui{Kontrollera saknade miniatyrbilder} undersöker importerade fotografier och registrerar gallerier med genererade varianter som saknas, är inaktuella, saknar metadata eller har fel bildförhållande, utan att skapa någonting. Det lagrade resultatet aktiverar \ui{Skapa saknade miniatyrbilder}, som bara behandlar rapporterade bilder. \ui{Skapa alla miniatyrbilder} under Underhåll omfattar alla gallerier; samma åtgärd i ett galleris bildredigerare omfattar det galleriet och, när \ui{Inkludera undergallerier} är valt, dess undergallerier. Dessa administrationsåtgärder förblir det främsta avsiktliga underhållsflödet. Offentlig bläddring startar ingen obegränsad reparationsgenomgång, men den aktuella renderaren har en liten skyddad bakgrundsförvärmning för specifika kort som behövde falla tillbaka på behöriga originalmedier eftersom ett förväntat derivat saknades.

\ui{Uppdatera miniatyrbildsdatabasen} stämmer av databasmetadata mot befintliga genererade filer och tar bort genererade filer med fel bildförhållande i stället för att publicera dem. \ui{Radera alla miniatyrbilder} är en separat bekräftad destruktiv cacheåtgärd; originalfotografier förblir orörda och derivaten kan byggas om senare.

\ui{Webbläsarbaserad ombyggnad av miniatyrbilder} väljs som standard för dessa omgenereringsåtgärder när webbläsarbearbetning är tillgänglig. Avmarkera kryssrutan för att använda servergenerering. Servern skickar originalfiler till den autentiserade webbläsaren i begränsade original-ZIP-delar, webbläsaren genererar miniatyrbildsvarianter och förberedda miniatyrbilds-ZIP-omgångar laddas upp tillbaka för servervalidering och installation. Gränser för originaldelar och omgångar konfigureras med webbläsaruppladdningsinställningarna så att arbetsflödet inte beror på en enda överstor PHP-begäran.

\subsubsection{Skyddad offentlig förvärmning av miniatyrbilder}
\index{f00095@förvärmning av miniatyrbilder}
\index{m00164@miniatyrbild!s00222@självreparation}

När servern har renderat ett offentligt tillgängligt fotografi genom reservvisning av originalmedier kan den markera exakt det fotografiet och de saknade stödda storlekarna som kandidat för bakgrundsförvärmning. Webbläsaren väntar på ledig tid, stannar när sidan är dold, skickar högst två kandidater per begäran, gör högst fyra begäranden från en sidinläsning och lämnar ungefär 1.8~sekunder mellan begärandena. Servern begränsar oberoende den normaliserade inskickade listan till 24 bild-ID:n, men den faktiska genereringsbudgeten är mindre: en anonym begäran behandlar högst två originalbilder inom ungefär fyra sekunder, medan en autentiserad administratörsbegäran behandlar högst fyra inom ungefär åtta sekunder.

En kandidat har en HMAC-token knuten till bild-ID, galleri-ID, filnamn, relativ sökväg och gallerimapp som servern ursprungligen renderade. Innan något genereras läser adressen om bilden och galleriet, verifierar token, upprepar aktuella besökaråtkomstkontroller, normaliserar begärda storlekar till den konfigurerade miniatyrbildsuppsättningen, konsulterar miniatyrbildsmetadata när de finns och hoppar över varianter som redan kan visas. En väntetid på 60~sekunder per bild förhindrar upprepade kostsamma återförsök. Ett enda globalt arbetarlås utan väntan innebär att samtidig offentlig trafik inte kan stapla flera bildavkodningsjobb; en upptagen begäran rapporterar bara att en annan förvärmningsarbetare är aktiv. Förvärmningen är~aktiverad som standard genom den interna programinställningen \setting{thumbnail_background_warmup_enabled} och skriver avgränsade diagnostikhändelser till administrationsloggen när logglagring är tillgänglig.

\securitynote{Bakgrundsförvärmning gör inte ett privat eller annars oåtkomligt fotografi till ett offentligt reparationsmål. Den signerade kandidaten bevisar att servern renderade bilden som reservkandidat, och adressen upprepar ändå aktuell åtkomstpolicy innan originalet öppnas eller derivat skrivs.}

\subsection{Om en miniatyrbild saknas}
\index{m00164@miniatyrbild!s00209@saknad}
\index{m00164@miniatyrbild!o00178@ombyggnad}

En saknad miniatyrbild kan visas som en långsammare reserv med originalfil, en tom förhandsvisning eller en underhållsvarning, beroende på bilden och tillgängliga filer. Öppna miniatyrbildsunderhållsverktygen, undersök berört galleri eller berörda storlekar och bygg om saknade varianter. Granska skrivbehörigheter och bildformatstöd när ombyggnad misslyckas upprepade gånger.

Efter reparationen öppnar du galleriet med tom webbläsarcache och kontrollerar att omslag, kort och den stora visaren får förväntade förhandsvisningar.
\section{Ljuslåda, kartor, EXIF och röstning}
\label{sec:lightbox}
\index{l00152@ljuslåda}
\index{e00072@EXIF}
\index{g00117@GPS}
\index{r00208@röstning}

\subsection{Asynkrona ljuslådemetadata}

Ljuslådan är den stora bildvisaren som öppnas ovanpå gallerisidan. Den låter besökare koncentrera sig på ett fotografi och samtidigt behålla kontroller för framåt, bakåt, stängning, helskärm, bildspel, karta och annat som är tillgängligt.

Galleriet förbereder visaren när den behövs och hämtar information om närliggande fotografier i hanterbara grupper. Besökaren kan därför bläddra genom ett stort galleri utan att första sidan behöver innehålla varje titel, förhandsvisning, kartpunkt och röst samtidigt.\footnote{Den aktuella visaren begär normalt information för grupper av närliggande fotografier. Åtkomstkontroller upprepas för varje begäran så att privata galleriregler fortfarande gäller.}

Vänster- och högerpiltangenterna går ett fotografi bakåt eller framåt. \codeword{Shift+Left} och \codeword{Shift+Right} flyttar tio positioner för snabbare förflyttning genom en stor resultatuppsättning. Hoppet följer samma stabila ordnade följd som vanlig navigering, går runt vid ändarna och fungerar även när den aktiva resultatuppsättningen är ett smart galleri. Den översatta tangentbordshjälpen listar dessa kortkommandon; inga separata grafiska hoppknappar läggs till.

Visaren börjar vanligen med en förhandsvisning i lämplig storlek. En större källa eller originalkälla kan användas där åtgärden och åtkomstpolicyn kräver det. Närliggande förhandsvisningar kan förberedas diskret så att nästa fotografi visas smidigt.

Medan ett nytt fotografi förbereds använder visaren samma förloppsområde som rapporterar inläsning i full kvalitet. Vid snabb förflyttning framåt eller bakåt får alltid det senast begärda fotografiet kontroll över indikatorn, så att en äldre förhandsvisnings- eller metadatabegäran inte kan lämna visaren permanent upptagen eller ersätta aktuellt urval. Om ett fotografi inte kan visas ger förloppsområdet ett felmeddelande som går att återhämta sig från; besökaren kan fortsätta med kontrollerna för framåt, bakåt, stängning eller återöppning. Dessa inläsningstillstånd ändrar inte regler för galleriåtkomst, skyddade medier, smarta gallerier, kartor, bildspel eller användning utan JavaScript.

\subsection{Zooma och panorera ett fotografi}
\index{l00152@ljuslåda!z00274@zoom}
\index{z00275@zoomkontroller}
\index{n00176@nypgest}

Ljuslådans zoomintervall är 100--400~\%. Använd minus- och pluskontrollerna, eller välj procentkontrollen för att återgå till 100~\%. Tangentbordsgenvägarna är \codeword{+} eller \codeword{=} för att zooma in, \codeword{-} eller \codeword{\_} för att zooma ut och \codeword{0} för återställning. Fokusindikatorer och aktuell procentsats är fortsatt tillgängliga för hjälpmedel.

Zoomning med mushjul och styrplatta följer pekaren: delen av fotografiet under markören stannar på samma plats när skalan ändras, även under flera snabba zoomsteg. Nypzoom använder gestens mittpunkt. När fotografiet är förstorat kan det dras vågrätt och lodrätt. Vid 100~\% på pekenheter fortsätter vanlig vågrät svepning att byta fotografi. Hjulinmatning med Control/Command förblir tillgänglig för webbläsarens sidzoom.

Vid 100~\% passas fotografiet in och centreras i visaren. Helskärmsläget använder samma zoom- och panoreringsbeteende, och dess Stäng-kontroll och övriga kontroller går att klicka på när bilden är förstorad.

Zoomning över 100~\% begär omedelbart den skyddade källan i full kvalitet för det aktiva fotografiet. Den befintliga förhandsvisningen förblir synlig tills den större bilden är klar, och aktuell zoom- och panoreringsposition bevaras. I vanligt läge, helskärmsläge och mobilt helskärmsläge visar överföringen samma förloppsrad med procenttal och byteantal. Att byta fotografi, stänga visaren eller starta en nyare begäran om full kvalitet avbryter föregående överföring. En misslyckad eller avbruten begäran om full kvalitet lämnar den skyddade förhandsvisningen på plats.

Navigering, bildspelsstart, stängning och återöppning återställer visaren till centrerade 100~\%. Att gå in i eller lämna helskärm behåller aktuell skala när själva fotografiet inte har ändrats och anpassar panoreringen till det nya visningsområdet. Inställningar för minskad rörelse tar bort den korta stegvisa zoomövergången.

Zoom påverkar endast presentationen. Den ändrar inte det lagrade fotografiet och kringgår inte mediebehörighet. Regler för lösenord, delningslänk, NSFW, smarta gallerier, karta, röstning, paginering och ljuslådemetadata gäller fortfarande.

\subsection{EXIF- och kartpolicy}

EXIF är information som många kameror och telefoner sparar inuti ett fotografi. Den kan innehålla tagningsdatum, kamera, objektiv, exponering, mått, orientering och GPS-koordinater. PHP Gallery kan använda informationen för datumförslag, kartor, tekniska detaljer, ordningshjälp och dubblettgranskning.

Kartor läses in när besökaren begär dem, vilket sparar arbete för besökare som bara är intresserade av fotografierna. Ett galleri kan visa enskilda bildplatser och en översikt över tillgängliga punkter. Över- och undergalleriers inställningar avgör vilka gallerier som visar informationen. Den globala inställningen för offentlig EXIF/GPS-visning är reserv för gallerier vars åsidosättning är satt till arv. Administrationen visar också hur många gallerier som för närvarande har en uttrycklig åsidosättning och kan återställa alla dessa till arv i en enda bekräftad inställningssparning. Återställningen är schemasäkrad: om åsidosättningstillståndet inte kan verifieras nekar programmet ändringen och hänvisar administratören till Systemstatus i stället för att gissa.

Att välja \ui{Öppna fotografi} i en gallerikartmarkör öppnar fotografiet i aktuell ljuslåda när det hör till det aktiva galleriet. Det fungerar också när fotografiet ligger utanför den pagineringssida som visas: visaren begär bara ett begränsat metadatafönster kring valt fotografi. I helskärmsläge med delad karta förblir kartan öppen medan bildrutan ändras. Från ett vanligt kartöverlägg stängs kartan och visar valt fotografi i visaren. Om den utökade visaren inte kan hitta urvalet följer samma kontroll den kanoniska bildsidan; den länkar inte direkt till den råa bildfilen.

\securitynote{GPS-data kan avslöja känsliga platser. Administratörer bör granska ärvda kartinställningar, bildmetadata och offentlig åtkomst innan personliga fotografier publiceras.}

\subsection{Flygkartor och navigeringsdata}
\index{s00220@SimBrief}
\index{f00081@flygkartor}
\index{n00173@navigeringsdata}

Flyginriktade samlingar kan använda den valfria SimBrief-integrationen för att hämta senaste operativa färdplanen för ett~pilot-ID eller pilotnamn. Den kompakta redigeraren accepterar en synlig identifierare: enbart siffror väljer vägen för pilot-ID; annan text väljer vägen för pilotnamn. Befintliga integrationer som skickar separata fält förblir kompatibla, med pilot-ID som förstahandsval när båda anges. Redigeraren skapar ett redigerbart Markdown-beskrivnings- och ruttutkast; OFP och tillgängliga ruttkoordinater kopplas först efter att galleriet sparats. En misslyckad SimBrief-begäran lämnar det vanliga beskrivningsfältet och den manuella ruttredigeraren tillgängliga. Utkastet före skapande är privat och tillfälligt enligt avsnittet om galleriskapande.

SimBrief-förhandsvisningen fyller den befintliga ruttredigeraren med hela den inlämnade rutten. Saknade avgångs- och ankomstflygplatser läggs till en gång vardera, medan ruttinstruktioner och flygplats-/banangivelser bevaras. Beskrivningsutkast på alla språk som underhålls behåller hela rutten, inklusive destinationen vid långa rutter. Statusen förklarar att OFP och tillgängliga kartpunkter kopplas när du sparar galleriet; själva importen sparar dem inte. Om du ändrar rutt, identifierare, beskrivning eller källspråk medan begäran pågår bevaras formulärvärdena och svaret ber dig importera igen. Efter sparandet behåller kartan de ursprungliga OFP-koordinaterna.

Manuell rutttext tolkas i första hand offline. Importerade OurAirports-databasrader kontrolleras först, sedan den lilla medföljande CSV-reserven. Ett giltigt cachelagrat leverantörsresultat kan användas därefter; när en administratör har kopplat ett konfigurerat Navigraph-konto och fjärrsökning är tillåten kan en okänd identifierare slås upp genom Navigraph och cachelagras under känd AIRAC-cykel. Hanteraren för Navigeringsdata visar antal i lokal databas och medföljande data, ger ett uppslagstest för identifierare såsom flygplatser, navigeringshjälpmedel och fixpunkter och kan uppdatera det lokala OurAirports-underlaget för flygplatser och navigeringshjälpmedel. Lagring för Navigraph OAuth och konton är valfri och kan kopplas bort utan att lokala data tas bort.

Med JavaScript aktiverat kontrolleras lokala navigeringsdatas aktualitet i bakgrunden även när en rutt som inte är tom matas in eller får fokus, när ett SimBrief-utkast begärs eller när ett formulär med rutt sparas. Både \ui{Flygkartor} och \ui{Navigeringsdata} måste vara aktiverade. OurAirports-data uppdateras när den senaste lyckade importen är minst en vecka gammal, med en timme mellan misslyckade försök. Gallerisparandet väntar inte på denna begäran; redigeringsutkastet, den öppna sidopanelen och webbläsaradressen bevaras. Galleriet kan sparas även om uppdateringen misslyckas.

Efter kontrollen kan sparade manuella rutter få saknade koordinater från aktuella navigeringsdata. Befintliga koordinater och SimBrief OFP-geometri bevaras, och en nyare ruttändring eller radering hindrar bakgrundsresultatet från att skriva över detta tillstånd. Endast sparade rutter kompletteras; osparad webbläsartext skickas inte med uppdateringen. Berörda offentliga kartor uppdateras genom det gemensamma sidopanelsflödet. Import och ruttskrivningar kräver verifierad databaslagring; saknad eller osäker lagring hindrar den valfria skrivningen medan vanlig galleriredigering förblir tillgänglig.

Offentlig kartvisning anropar aldrig SimBrief, Navigraph eller någon annan aktiv planeringstjänst. Den visar endast koordinater som redan lagrats med galleriet. Sparade SimBrief OFP-koordinater förblir därför användbara även om fjärrleverantören senare är otillgänglig. Manuell rutttext kan också helt kringgå identifieraruppslag med uttryckliga \codeword{NAME@latitude,longitude}-punkter.

\subsection{Röstning}

Offentliga kort och ljuslådekontroller delar bildröstningstillstånd. Den aktuella offentliga rösten är en återkallelig positiv röst: att trycka på den aktiva positiva rösten igen tar bort den. Den visade poängen är summan av lagrade bildröster; det finns ingen separat adress för röstning på hela gallerier i den aktuella implementationen. Inloggade röster kopplas till kontot, medan anonyma röster använder programmets integritetsskyddande besökarhash. Nya positiva röster är frekvensbegränsade; återkallning av en befintlig röst tillåts omedelbart. Serverrenderade formulär ger direkt funktion, medan JavaScript skickar vanliga interaktioner asynkront och synkroniserar poängen mellan matchande vyer. Servern validerar bildidentitet, galleripolicy, åtkomst, schemaberedskap, CSRF och begärans integritet.

\subsection{Bildspelet}
\index{b00038@Bildspelet}
\index{b00034@bildjämförelse}

Bildspelet är ett valfritt offentligt jämförelseläge. En administratör aktiverar det för en gallerigren; lämpliga offentliga fotografier från galleriet och dess undergallerier kan då delta enligt vanliga åtkomst- och synlighetsregler. Det offentliga spelet visar två fotografier sida vid sida, registrerar vilket par som visades och låter besökaren välja den föredragna bilden innan nästa jämförelse visas. Galleriet kan också presentera ledande resultat som samlats in genom spelet.

Eftersom visning av ett nytt par skapar beständigt speltillstånd redan innan besökaren väljer en vinnare kräver Bildspelet att dess databasschema är känt och tillgängligt. Om den nödvändiga lagringen inte kan verifieras nekas skrivningen i stället för att tyst falla tillbaka på en oregistrerad jämförelse. Den globala funktionsväxeln kan dölja Bildspelet utan att radera befintliga speldata eller aktiveringsinställningar per galleri.

\section{Dubblettdetektor för fotografier}
\label{sec:duplicates}
\index{d00066@Dubblettdetektor för fotografier}
\index{d00067@dubblettregister}
\index{s00216@SHA-256}

Administrationens Dubblettdetektor för fotografier analyserar som standard en vald gallerigren och kan utökas till global omfattning när administratören aktiverar motsvarande alternativ. Den rapporterar exakta och möjliga dubblettpar, ger länkar till galleri- och bildsammanhang, registrerar granskningsbeslut och överlåter bekräftad radering till den vanliga tjänsten för galleribildsradering.\footnote{Detektorn förblir ett granskningsverktyg. Underlaget stöder administratörens bedömning eftersom matchande filstorlek och fotografiska metadata kan beskriva olika fotografier i vissa arbetsflöden.}

\subsection{Underlagskategorier}

\subsubsection{Exakta och möjliga fynd}

En exakt dubblett innehåller samma filinnehåll. Programmet känner igen detta genom en kontrollsumma, som fungerar som ett mycket detaljerat digitalt fingeravtryck beräknat från filen. Matchande fingeravtryck är ett starkt tecken på att två filer innehåller samma byte.

En möjlig dubblett har stödjande likheter utan en exakt fingeravtrycksmatchning. Detektorn kan jämföra filstorlek och tillgängliga kamerauppgifter såsom mått, tagningstid, kamera, objektiv, exponering, bländare, brännvidd, ISO och orientering. Två separat exporterade kopior av ett fotografi kan skilja sig som filer men ändå behålla tillräckligt mycket gemensam information för att behöva granskas.

Saknade kamerauppgifter ger helt enkelt detektorn färre ledtrådar. Resultaten förblir förslag för en människa att undersöka. Öppna båda bildlänkarna, jämför gallerisammanhang och kvalitet och radera bara den kopia vars borttagning passar samlingen.

\subsection{Beständigt granskningsregister}

Registret tillhör den autentiserade administratören. Parposter lagrar kanoniska bildidentifierare, med den lägre identifieraren i den låga kolumnen och den högre i den höga kolumnen. Galleriposter gäller exakt det lagrade galleriet. Beslut för över- och undergallerier förblir oberoende.

Tillgängliga granskningsåtgärder:

\begin{itemize}
  \item ignorera valt par;
  \item ignorera alla fynd från den vänstra bildens exakta galleri;
  \item ignorera alla fynd från den högra bildens exakta galleri;
  \item rensa den aktuella administratörens register;
  \item radera en bekräftad dubblett genom validerad vanlig bildradering.
\end{itemize}

AJAX-åtgärder uppdaterar detektorfragmentet i den befintliga högra administrationssidopanelen. Direkta POST-begäranden använder omdirigering som reserv. Serverägda skanningsjobb bevarar omfattningen mellan fortsättningsbegäranden, och radering tar bort den berörda bilden från aktuella jobbresultat.
\section{Åtkomstkontroll och säkerhet}
\label{sec:security}
\index{s00237@säkerhet}
\index{c00047@CSRF}
\index{a00017@autentisering}
\index{n00174@NSFW}

\subsection{Åtkomstbedömning i flera lager}

Integriteten bedöms innan skyddad utmatning sammanställs. Gallerisynlighet och ärvda åtkomstregler kombineras med synlighet per fotografi, lösenords- och delningslänkstillstånd samt åldersbegränsningsregler där de är konfigurerade. Samma behörighetskontroll gäller gallerisidor, miniatyrbilder, förhandsvisningar, originalfiler, ljuslådemetadata, kartor, sökresultat och hämtningar.

Administrationsåtgärder kräver ett autentiserat konto. Formulär innehåller en privat begärandetoken, och servern validerar målgalleriet, fotografiet och den begärda åtgärden innan beständiga data ändras.\footnote{Det tekniska namnet på den privata formulärtoken är CSRF-token. Galleriägare möter den normalt bara genom vanliga administrationsformulär.}

\subsection{Lösenord och delningslänkar}
\index{l00155@lösenordsskydd}
\index{d00059@delningslänkar}

Skyddade gallerier lagrar lösenordshashar och ger sessionsbegränsad åtkomst efter lyckad verifiering. Delningslänkar använder servergenererat tokentillstånd och utgångspolicy. Lösenordsåterställning använder engångstoken som lagras hashade, en valfri återställningsadress, konfigurerad giltighetstid och vald e-posttransport. Använd HTTPS, starka administratörsuppgifter, begränsade databaskonton och skyddade e-post- och API-hemligheter.

\subsection{Besökarkonton och registrering endast via inbjudan}
\index{b00028@besökarkonton}
\index{b00026@besökarinbjudningar}
\index{h00119@hantering av besökarkonton}
\index{b00027@besökarinloggning}

Besökarkonton är ett valfritt identitetslager för personliga gallerifunktioner. De är åtskilda från administratörskontot och låser inte i sig upp ett privat eller lösenordsskyddat galleri. Befintliga gallerilösenord, delningslänkar, åtkomst till smarta gallerier och mediebehörighet avgör fortfarande vad besökaren får se.

Hela delsystemet för besökarkonton omges av en huvudväxel i \ui{Administration > Funktioner}. Denna huvudfunktion är avstängd som standard. När den är avstängd visar det offentliga galleriet inga kontroller för besökarinloggning eller konto, besökarrutter är otillgängliga och posten \ui{Besökarkonton} döljs från administrationens kontomeny. Befintliga besökaridentiteter, favoriter och privata samlingar finns kvar i lagringen; växeln styr tillgängligheten och raderar inte data.

Efter att huvudfunktionen aktiverats kan administratören öppna \ui{Konto > Besökarkonton}. Sidan behåller den separata lägesväxeln \ui{Besökarkonton (endast via inbjudan)}. Lägesväxeln styr besökarinloggning och inbjudningar inom det aktiverade delsystemet, medan huvudväxeln under Funktioner styr om delsystemet alls visas.

Administratören hanterar besökaridentiteter från \ui{Konto > Besökarkonton}. Det befintliga flödet med enbart inbjudningar finns kvar: administratören kan vid behov knyta en inbjudan till en avsedd e-postadress och kopiera dess hemliga länk som visas en gång. Att öppna länken skapar inget konto. Mottagaren skickar in en e-postadress, får ett verifieringsmeddelande genom den konfigurerade e-posttransporten, bekräftar verifieringen i webbläsaren och väljer sedan ett lösenord med minst 15~tecken. Det finns ingen allmän sida för \ui{Registrera} eller \ui{Skapa konto}.

Samma administrationssida kan också skapa ett besökarkonto direkt. Ange besökarens e-postadress och antingen ett tillfälligt lösenord som uppfyller besökarlösenordspolicyn eller lämna lösenordet tomt så att galleriet genererar ett starkt lösenord. Det tillfälliga lösenordet visas en gång efter skapandet. Om den valfria aviseringen om skapat konto skickas innehåller meddelandet den betrodda inloggningsadressen och instruktioner men aldrig det tillfälliga lösenordet; ge lösenordet separat genom en betrodd kanal. Med huvudfunktionen aktiverad är direkt kontoskapande tillgängligt även när det underordnade offentliga besökarläget är avstängt, så att konton kan förberedas i förväg, men dessa besökare kan inte logga in förrän besökarläget aktiveras.

En direkt skapad besökare kan inte använda det tillfälliga lösenordet för varaktig inloggning. Efter den första lyckade kontrollen av inloggningsuppgifterna skickas webbläsaren till den privata sidan för första inloggning och måste välja ett annat lösenord som uppfyller kraven. Tills bytet lyckas etablerar galleriet inte den vanliga besökaridentiteten, utfärdar inte \ui{Kom ihåg mig} och tillåter inte besökarbehörighet för favoriter eller samlingar. Den vanliga återställningen av glömt lösenord kan också ersätta den tillfälliga uppgiften genom sitt oberoende verifierade återställningsflöde. En samtidig administratörsinloggning i samma webbläsare förblir oberoende.

Befintliga besökarkonton listas på administrationssidan med kontotillstånd och säkerhetskontroller. En aktiv besökare kan stängas av med \ui{Stäng av} eller loggas ut överallt; en avstängd besökare kan återaktiveras med \ui{Återställ}. Avstängning blockerar besökarautentisering och ogiltigförklarar besökarens befintliga session, kom-ihåg-uppgifter, återställningsbehörighet, väntande verifiering och e-poständring samt förberedd behörighet för samlingsdelning genom den centrala livscykeltjänsten för besökare. Återställning gör kontot aktivt igen men återupplivar ingen inloggningsuppgift eller session som fanns före avstängningen, så besökaren måste autentisera sig igen på vanligt sätt. Ett tvingat lösenordsbyte vid första inloggning krävs fortfarande efter avstängning och återställning. \ui{Logga ut överallt} lämnar ett aktivt konto aktivt men återkallar dess besökarinloggningsbehörighet på alla enheter. Dessa säkerhetskontroller är tillgängliga för administratören när huvudfunktionen är~aktiverad men det underordnade offentliga besökarläget är avstängt. Om huvudfunktionen stängs av döljs även besökarhanteringen. Ingen av dessa växlar eller säkerhetskontroller tar bort favoriter, privata samlingar, samlingsobjekt, fotografier, gallerier, smarta gallerier, galleridelningslänkar eller administratörsautentisering.

Radering som utförs av administratören förblir en separat destruktiv åtgärd med uttrycklig bekräftelse. Den tar bort besökaridentiteten och besökarägda tillstånd för autentisering, livscykel, favoriter och samlingar genom de förberedda databasrelationerna. Den raderar inte fotografier, gallerier, smarta gallerier, galleridelningslänkar eller administratörskonton och administratörssessioner.

Det offentliga sidhuvudet visar \ui{Logga in} endast när besökarkonton är~aktiverade. Efter inloggning visar det \ui{Konto}. Kontosidan visar verifierad e-postadress, länkar till de privata sidorna \ui{Favoriter} och \ui{Samlingar}, \ui{Byt lösenord}, \ui{Byt e-postadress}, \ui{Radera konto} och en \ui{Logga ut}-åtgärd enbart för besökaren. Att logga ut en besökare eller ändra besökarkontots behörighet loggar inte ut en administratör som råkar vara autentiserad i samma webbläsarsession.

\index{f00076@favoriter}
En inloggad besökare kan markera ett behörigt fotografi som favorit med den lilla hjärtkontrollen på ett gallerikort eller i ljuslådan. Samma fotografi hålls synkroniserat mellan kortet och ljuslådan, och den privata sidan \ui{Favoriter} listar sparade fotografier som besökaren fortfarande får öppna. Att ta bort en favorit raderar bara besökarens sparade referens; det raderar eller ändrar inte fotografiet.

En favorit ger aldrig åtkomstbehörighet. Om källgalleriet senare blir privat, dess lösenords- eller delningsbehörighet löper ut eller fotografiet annars inte längre är tillåtet för den begäran kringgår den sparade favoriten inte reglerna och källmetadata visas inte på Favoriter-sidan. En webbläsare som samtidigt är inloggad som administratör och besökare behöver ändå besökarens oberoende källbehörighet innan besökaren kan spara bilden som favorit.

\index{b00029@besökarsamlingar}
En inloggad besökare kan också skapa privata namngivna \ui{Samlingar}. En samling är en ordnad lista med referenser till befintliga fotografier; den är inte ett galleri, kopierar inte fotografier och ändrar inte den administratörsägda gallerihierarkin. Från ett behörigt gallerikort, smart gallerikort eller den privata sidan \ui{Favoriter} kan besökaren använda \ui{Lägg till i samling} och välja en av sina egna samlingar. Det privata samlingsområdet stöder att skapa och byta namn på en samling, ta bort fotografier, flytta synliga fotografier uppåt eller nedåt och radera själva samlingen.

Ägande av samlingen och åtkomst till fotografier är separata beslut. Att spara ett fotografi lagrar bara dess interna bildreferens och ordning. Samlingen lagrar inte gallerilösenord, galleridelningstoken, filsystemssökväg, miniatyrbildssökväg eller något ihågkommet åtkomstbeslut. Varje gång samlingen öppnas utvärderar galleriet aktuell behörighet för källgalleriet och fotografiet på nytt. Om ett tidigare sparat fotografi inte längre är tillåtet utelämnas det utan att avslöja titel, filnamn, gallerititel, miniatyrbild, EXIF-data eller orsaken till nekad åtkomst. Den sparade referensen kan finnas kvar så att fotografiet blir synligt igen om vanlig källåtkomst senare återkommer.

Regeln gäller också när samma webbläsare är inloggad både som administratör och besökare. Administratörens synlighet blir inte besökarbehörighet för samlingar, och en samling kan aldrig användas som en andra väg runt gallerilösenord, delningsbehörigheter, åldersbegränsningar eller medieleveranskontroller. Favoriter och samlingar är oberoende: ett fotografi kan finnas i endera, båda eller ingen, och radering av en samling raderar aldrig en favorit eller det underliggande fotografiet.

\index{s00212@samlingsdelning}
En besökare kan skapa en olistad skrivskyddad delningslänk från en egen samling. Länken löper ut efter 30~dagar och kan ersättas eller återkallas när som helst. Den fullständiga hemliga länken visas bara direkt efter skapande eller ersättning; senare visar samlingssidan endast att en delning är aktiv samt tidpunkterna för skapande och utgång. Att ersätta en delning gör den tidigare länken oanvändbar. Att återkalla en delning tar endast bort delningsbehörigheten och raderar inte samlingen, dess objekt, favoriter eller fotografier.

Mottagaren behöver inget besökarkonto. Att öppna den hemliga länken validerar den, etablerar endast en smal sessionsbehörighet för den enda samlingen och omdirigerar till en ren adress utan hemligheten. Länken kan återanvändas tills den löper ut eller återkallas, så automatiska länkskannrar i e-post eller chattar förbrukar den inte. Den delade sidan är olistad, kan inte cachelagras och är markerad för att inte sökindexeras.

Delningen ger endast åtkomst till samlingsbehållaren. Den låser aldrig upp ett källgalleri eller fotografi. Vid varje delad visning kontrolleras varje lagrad bildreferens igen mot mottagarens aktuella källgalleribehörighet. En lösenordsskyddad källa förblir dold tills mottagaren själv låser upp den genom den vanliga gallerimekanismen; en oberoende befintlig galleridelning kan på samma sätt ge vanlig behörighet. Även när samma webbläsare också är inloggad som administratör används inte administratörens undantag för att fylla den delade samlingen. Direkta miniatyrbilds- och medierutter fortsätter tillämpa sina egna befintliga källregler. Oåtkomliga objekt utelämnas utan att avslöja dold titel, filnamn, sökväg, galleri eller orsak.

Att stänga av samlingsägaren återkallar aktiva samlingsdelningar, och återställning av kontot återupplivar dem inte. \ui{Logga ut överallt} återkallar bara besökarinloggningsbehörighet och återkallar inte en annars giltig delning. Att radera ägaren eller samlingen ogiltigförklarar dess delning utan att radera originalfotografier. Den globala funktionsväxeln för Besökarkonton styr också samlingsdelning och är fortfarande avstängd som standard.

Besökarens \ui{Kom ihåg mig} använder en särskild roterande besökaruppgift i stället för administratörens token för beständig inloggning. Att återställa den uppgiften loggar bara in besökaren och räknas inte som nylig återautentisering med lösenord. Begäranden om glömt lösenord ger ett generiskt svar och använder återställningsbekräftelse i två steg så att automatiska länkskannrar inte kan återställa ett lösenord genom att öppna en URL.

Känsliga kontoändringar kräver ett nyligt bevis med besökarens lösenord. En besökare som endast återställts genom \ui{Kom ihåg mig} ombeds därför bekräfta aktuellt besökarlösenord före lösenordsbyte, påbörjad e-postadressändring eller kontoradering. Returvägen är begränsad till dessa kontoåtgärder och kan inte omdirigera till en godtycklig sida.

Att byta besökarens lösenord använder samma besökarlösenordspolicy som aktivering och återställning. Ett lyckat byte ogiltigförklarar den gamla lösenordsbehörigheten, besökarsessioner och besökarens kom-ihåg-uppgifter enligt säkerhetslivscykeln och ber sedan besökaren logga in igen. Sparade favoriter förblir kopplade till samma besökarkonto. En administratör som är inloggad genom samma webbläsarsession förblir inloggad.

Byte av besökarens e-postadress sker stegvis. Att skicka in en ny adress ersätter inte den aktuella verifierade adressen. Galleriet skickar först en verifieringslänk till den föreslagna adressen inom besökarens gränser mot e-postmissbruk. Att öppna länken validerar den bara och förbereder ett kortlivat bekräftelsetillstånd; det ändrar inte kontot. Den slutliga ändringen sker först efter att besökaren skickat in det separata skyddade bekräftelseformuläret. Ett lyckat byte loggar ut besökaren så att den nya verifierade adressen kan användas vid nästa inloggning.

\ui{Radera konto} är en permanent åtgärd endast för besökaren. Efter nylig lösenordsbekräftelse måste besökaren uttryckligen bekräfta raderingen. Åtgärden tar bort besökarkontot och besökarägda data såsom favoriter och privata samlingar enligt den förberedda databaslivscykeln. Den raderar inte gallerifotografier, gallerier, smarta gallerier eller deras befintliga delningslänkar, och den förstör inte en samtidig administratörsinloggning.

Besökarsäkerhet och privata sidor för konto, favoriter, samlingar och livscykel levereras som privata svar som inte kan cachelagras. Inbjudnings-, verifierings-, återställnings- och e-poständringslänkar använder programmets konfigurerade betrodda ursprung. E-post för besökarverifiering, återställning och e-poständring omfattas fortfarande av missbruksbudgetar för besökaradress, besökare och global användning innan den konfigurerade PHP~mail- eller SMTP-transporten används.

Olistad skrivskyddad samlingsdelning finns endast inom funktionen Besökarkonton. Offentlig upptäckt av samlingar, offentliga besökarprofiler, mottagarredigering och samarbete, besökaruppladdningar och kommentarer, öppen registrering, social besökarinloggning, TOTP, passkeys och inloggning med magisk länk är fortfarande otillgängliga.

\subsection{Administratörskonto och e-postinställningar}
\index{a00006@administratörskonto}
\index{s00228@SMTP}
\index{g00116@Google-inloggning}

Kontoområdet stöder byte av användarnamn och lösenord, en valfri återställningsadress, giltighetstid för lösenordsåterställning och ett testmeddelande. Återställningspost kan använda PHP \codeword{mail()} eller en autentiserad SMTP-transport med inställningar för värd, port, kryptering, användarnamn och lösenord. Behandla SMTP-uppgifter som databasuppgifter och ta aldrig med dem i exporter eller diagnostikskärmbilder.

Google-inloggning är valfri. Administratören konfigurerar OAuth-klienten och godkänd återanrops- eller omdirigerings-URI och kopplar sedan en Google-identitet endast när administratören redan är autentiserad med en vanlig lösenordssession. En okopplad Google-identitet avvisas på inloggningssidan och får inte skapa eller göra anspråk på en administratörsprofil. Efter koppling kan bara den lagrade Google-identiteten logga in genom Google-knappen för profilen. Att koppla bort länken kräver aktuellt administratörslösenord. Extern inloggning ersätter inte lösenordsinloggning; om OAuth-konfiguration eller lagring för identitetskoppling saknas eller inte kan verifieras förblir vanlig lösenordsautentisering oberoende.

Inloggningsformuläret accepterar administratörens användarnamn eller konfigurerade återställningsadress. Autentiseringsförsök begränsas både efter besökare och normaliserad inskickad identifierare när frekvensbegränsningsschemat är tillgängligt. Begränsningstabellen lagrar hashade subjekt i stället för råa IP-adresser eller råa inskickade identifierare. Aktuell policy begränsar upprepade misslyckade inloggningar inom ett 15-minutersfönster och lösenordsåterställningsbegäranden inom längre begränsade fönster; lyckad inloggning rensar relevant tillstånd för inloggningsbegränsning.

Lösenordsåterställning ger avsiktligt ett generiskt begäranderesultat så att det offentliga formuläret inte bekräftar om ett inskickat användarnamn eller en e-postadress finns. Återställningslänkar använder engångstoken som lagras hashade och en konfigurerad giltighetstid på 15 till 1440~minuter. Att slutföra en återställning ändrar lösenordet, markerar token som använd och återkallar token för beständig inloggning för kontot. Kontosidan kan skicka ett testmeddelande för återställning genom konfigurerad PHP \codeword{mail()} eller SMTP-transport.

Alternativet \ui{Håll mig inloggad} använder en hashad webbläsartoken så att en administratör kan förbli inloggad genom vanlig PHP-sessionsrensning på delade webbhotell utan att den råa beständiga uppgiften lagras i databasen. Beständig inloggning förblir valfri. Den återkallas av relevanta säkerhetsflöden för konto och utloggning och kan vara otillgänglig när dess schematillstånd inte kan verifieras, utan att aktuell vanlig PHP-session stängs av.
\subsection{Kontrakt för databasberedskap}
\label{sec:db-readiness-contract}
\index{s00236@Systemstatus!d00051@databasberedskap}
\index{d00049@databas!b00023@beredskap}
\index{a00001@503-svar}

Funktioner som beror på databasschema använder ett gemensamt beredskapskontrakt. Den viktiga skillnaden är mellan ett schema som har undersökts framgångsrikt och konstaterats vara äldre, och ett schema som för närvarande inte kan undersökas.

\begin{longtable}{>{\RaggedRight\arraybackslash}>{\hspace{0pt}}p{0.12\textwidth} >{\RaggedRight\arraybackslash}>{\hspace{0pt}}p{0.20\textwidth} >{\RaggedRight\arraybackslash}>{\hspace{0pt}}p{0.35\textwidth} >{\RaggedRight\arraybackslash}>{\hspace{0pt}}p{0.17\textwidth}}
\toprule
\textbf{Tillstånd} & \textbf{Betydelse} & \textbf{Programbeteende} & \textbf{Ägarens åtgärd}\\
\midrule
\endfirsthead
\toprule
\textbf{Tillstånd} & \textbf{Betydelse} & \textbf{Programbeteende} & \textbf{Ägarens åtgärd}\\
\midrule
\endhead
Tillgängligt & Nödvändig lagring hittades och kontrollerades. & Vanliga läsningar och skrivningar för funktionen. & Ingen.\\
Saknas & Undersökningen lyckades, men ett nödvändigt objekt saknas. & En dokumenterad väg för äldre schema kan finnas kvar; annars kräver funktionen en migrering. & Säkerhetskopiera och tillämpa väntande migrering.\\
Okänt & Schemat kunde inte undersökas tillförlitligt. & Skyddade läsningar kan nekas av säkerhetsskäl eller valfri skrivskyddad utmatning utelämnas. Beroende skrivningar blockeras. & Reparera databasanslutning, valt schema eller metadatabehörigheter.\\
Avstängt & En valfri funktion är avstängd. & Funktionen används inte. & Ingen, om inte funktionen ska aktiveras.\\
\bottomrule
\end{longtable}

\ui{Systemstatus} och \ui{Körningsdiagnostik} använder dessa tillstånd konsekvent. Resultatet \ui{Saknas} pekar normalt på migreringsarbete; \ui{Okänt} är ett driftproblem och ska inte behandlas som bevis för att en tabell eller kolumn aldrig funnits. Diagnostik kan visa funktionsnamn, status, validerade databasobjektnamn, föreslagna kontroller och en begärandereferens. Den utesluter SQL-text, råa databasundantag, inloggningsuppgifter, token, kakor och privata filsystemssökvägar.

\subsection{Säkerhetskänsliga undantag}
\index{n00175@NSFW-skydd}
\index{l00156@lösenordsåterställning!d00051@databasberedskap}
\index{g00116@Google-inloggning!d00051@databasberedskap}
\index{d00059@delningslänkar!d00051@databasberedskap}

Några funktioner har strängare eller mer specifikt beteende:

\begin{description}
  \item[NSFW-skydd.] Om undersökningen bekräftar det äldre schemat utan NSFW-fält behåller webbplatsen sitt beteende från före funktionen tills migrering sker. Om fälten är \ui{Okända} nekas offentliga begäranden som kan exponera åldersbegränsade medier av säkerhetsskäl, och NSFW-redigering pausas.
  \item[Galleriåtkomst.] En databas behandlas som det äldre formatet från före skyddet endast när undersökningen lyckas och alla grundläggande åtkomstfält saknas. En blandad struktur är en ofullständig migrering, så skyddade rutter förblir otillgängliga tills den repareras. Det äldre publiceringsvärdet \codeword{draft} skrivs endast när fältdefinitionen är känd att kräva det; aktuella scheman använder \codeword{unpublished}.
  \item[Återkallning av delningslänk.] Återkallning kan använda den mindre uppsättningen kända fält som behövs för att rensa den validerande hashen eftersom åtgärden minskar åtkomsten. Den stannar ändå när dess egna nödvändiga fält är \ui{Okända}.
  \item[Beständig administratörsinloggning.] En saknad tabell för ihågkommen inloggning stänger av beständiga webbläsartoken men hindrar inte en vanlig lösenordsinloggning från att skapa aktuell PHP-session.
  \item[Lösenordsåterställning.] Återställning kräver både återställningsadressfältet och lagring för återställningstoken. Saknad lagring ger migreringsvägledning; \ui{Okänd} lagring ger ett tillfälligt fel i stället för ett felaktigt konto- eller tokenresultat.
  \item[Google-identitetskopplingar.] Extern inloggning kräver både giltig OAuth-konfiguration och känd lagring för identitetskoppling. Om lagringen saknas eller är okänd förblir lösenordsinloggning tillgänglig oberoende.
\end{description}

\subsection{Skrivningar, underhåll och inloggningsuppgifter}
\index{d00049@databas!z00281@ändringssäkerhet}
\index{u00254@uppladdningar!d00051@databasberedskap}
\index{u00253@uppdateringar!d00051@databasberedskap}

Åtgärder som ändrar filer, inloggningsuppgifter eller beständiga poster utför sin beredskapskontroll före det första steget som ändrar målet. Det omfattar galleri- och bildändringar, uppladdningar, WebDAV-skrivningar, gallerimigrering, miniatyrbildsunderhåll, databasreparation och rensning samt aktivering av programuppdateringar. Resultatet \ui{Okänt} stoppar åtgärden innan målet ändras. Ett bekräftat resultat \ui{Saknas} får använda en kompatibilitetsväg för äldre schema eller en migreringsuppstart endast där beteendet är dokumenterat.

Återkallning av inloggningsuppgifter följer samma restriktiva princip som återkallning av delningslänkar: den kan använda en mindre känd fältuppsättning när fälten räcker för att stänga av åtkomst. Att skapa eller lita på en uppgift kräver fortfarande hela lagringen som autentiseringsvägen använder.

När en skrivning blockeras som \ui{Okänt} reparerar du databasundersökningen innan du försöker igen. När tillståndet blir \ui{Saknas} säkerhetskopierar du och tillämpar den nödvändiga migreringen. När det blir \ui{Tillgängligt} försöker du åtgärden en gång till och kontrollerar både filsystemets och databasens resultat. Återuppta ett befintligt återställningsbart migrerings- eller uppdateringsjobb i stället för att starta ett dubblerat jobb.

\subsection{Valfria funktioner}
\index{k00133@kartor!d00051@databasberedskap}
\index{b00038@Bildspelet!d00051@databasberedskap}
\index{r00208@röstning!d00051@databasberedskap}
\index{o00180@OpenAI!d00051@databasberedskap}
\index{s00220@SimBrief!d00051@databasberedskap}
\index{t00243@telemetri!d00051@databasberedskap}

Valfria presentationsfunktioner kan få begränsad funktion utan att huvudgalleriet stängs av när det inte kan avslöja skyddat innehåll eller ge åtkomst. En karta kan t.ex. utelämnas om dess valfria lagring inte kan läsas, men att spara en GPS- eller ljuslådeåsidosättning kräver ändå att det relevanta fältet är känt.

Både röstning och Bildspelet skriver beständigt tillstånd. Bildspelet registrerar ett visat par innan en vinnare väljs, så inläsning av ett nytt par är också en skrivgräns och kräver känd lagring. AI-köer, telemetriunderhåll, navigeringstoken, sparade SimBrief-ruttdata och liknande valfria skrivningar följer samma regel. En funktion kan ge ett skrivskyddat resultat utan valfri lagring endast där gränssnittet tydliggör begränsningen.

\subsection{Driftsäkerhet}

Rekommenderade arbetssätt:

\begin{itemize}
  \item Leverera programmet via HTTPS.
  \item Begränsa databasuppgifterna till programmets schema och nödvändiga åtgärder.
  \item Håll \filepath{config.php}, säkerhetskopior, cache och genererade arkiv utanför offentlig upptäckt där webbhotellets upplägg tillåter det.
  \item Tillämpa migreringar och uppdateringar genom autentiserade underhållsflöden.
  \item Granska granskningsloggar och säkerhetsdiagnostik.
  \item Testa skyddade gallerier genom anonyma besökarsessioner.
  \item Ha samordnade säkerhetskopior av filsystem och databas.
\end{itemize}
\section{Tema, presentation och lokalisering}
\label{sec:theme}
\index{t00244@tema}
\index{l00154@lokalisering}

Temaområdet styr webbplatsens identitet, färger, mått, sidbredd, bakgrunder, profilering, favicon-resurser, språk, gallerikortspresentation, standardvärden för ljuslådan, GPS-presentation, favoritnavigering och offentlig miniatyrbildsrendering. Inställningar använder normaliserade skalära värden eller avgränsade tjänstemodeller.

Underfliken \ui{Animationer} ställer in öppnings- och stängningstiden för den offentliga informationsrutan på gallerikort och Admin-sidopanelen. Båda värdena kan vara 0--800~ms; 0 stänger av respektive rörelse. Standardvärdena är 320~ms respektive 260~ms. Den offentliga panelen innehåller galleriets titel, beskrivning, datum och hela tagglistan. Den inbyggda utfällningen fungerar även utan JavaScript; med JavaScript hålls panelen inom visningsområdet och kan stängas genom att klicka utanför eller trycka på Escape. Interaktionen följer webbläsarens inställning för minskad rörelse.

Temaredigeraren använder ingen separat växel för ljust eller mörkt läge. I stället definierar sparade färgkontroller och ett valfritt utseende med Egen CSS direkt den offentliga färgpaletten. Den ställer också in layouter för gallerikort och undergallerier och kan konfigurera upp till tre genvägar i sidhuvudet. Varje genväg kan peka på startsidan, ett galleri eller förbli tom. Anonyma besökare ser bara konfigurerade gallerigenvägar som fortfarande är offentliga och listade; dubbla eller raderade gallerival tas bort när inställningarna sparas. Om alla tre platserna är tomma döljs dessa favoritgalleriknappar. Gallerikort kan också visa ett konfigurerbart märke för antalet ingående bilder. När märket gäller för ett galleri förblir samma grenantal synligt i huvudpanelen när galleriet öppnas, och räknar bilder i aktuellt galleri och dess åtkomliga undergallerier enligt samma offentlighets- och åtkomstregler som kortet. Presentation per galleri kan ärva webbplatsens standardvärden eller åsidosätta de val som galleriredigeraren erbjuder.

\subsection{Referens för Temas kontroller}
\index{t00244@tema!k00137@kontroller}
\index{p00182@paginering!t00245@temainställningar}
\index{f00075@favicon}
\index{e00068@Egen CSS}

Den aktuella temasidan är uppdelad i \ui{Utseende}, \ui{Profilering \& medier}, \ui{Layout}, \ui{Språk} och \ui{Egen CSS}. Följande referens beskriver kontrollerna som ägaren ser och den normalisering som den aktuella programversionen tillämpar.

\begin{longtable}{>{\RaggedRight\arraybackslash}>{\hspace{0pt}}p{0.20\textwidth} >{\RaggedRight\arraybackslash}>{\hspace{0pt}}p{0.27\textwidth} >{\RaggedRight\arraybackslash}>{\hspace{0pt}}p{0.41\textwidth}}
\toprule
\textbf{Område} & \textbf{Kontroll} & \textbf{Aktuellt beteende}\\
\midrule
\endfirsthead
\toprule
\textbf{Område} & \textbf{Kontroll} & \textbf{Aktuellt beteende}\\
\midrule
\endhead
Utseende & Webbplatsidentitet och färger & Webbplatsnamn samt färger för accent, mörk accent, sidbakgrund, panel, öppet galleris panel, sidhuvudstitel och gallerititel. Den ''mörka accenten'' är en färg för hovring och konturer, inte en väljare för mörkt läge.\\
Utseende & Hörn och typsnitt & Hörnradien är 0--32~pixlar. Typsnittsläget är \ui{Klassiskt med seriffer} eller \ui{Rent utan seriffer}. Direktförhandsvisningen uppdateras när kontrollerna redigeras.\\
Utseende & Offentlig sidbredd & \ui{Standard}, \ui{Bredare}, \ui{Full bredd} eller \ui{Egen}. Egen bredd begränsas till 1024--2048~pixlar och får standardvärdet 1440 när giltigt värde saknas.\\
Utseende & GPS-markör & När GPS-kartfunktionen är~aktiverad kan bildkortens markörer och deras bakgrundsplatta visas eller döljas separat. Markörstorleken är 14--48~pixlar med 26 som standard. Bakgrundsplattans reglage visas som 0--48~pixlar, men aktuell servernormaliserare behandlar 0 eller ett annat icke-positivt värde som ej angivet och återställer 22~pixlar; använd den separata kryssrutan för att dölja bakgrunden.\\
Utseende & Offentliga taggsidor & Särskilda kolumner och rader för taggsidor använder samma säkra maxgränser som andra offentliga rutnät: 1--12 kolumner och 1--50 rader. Taggresultatkort kan använda lodrät eller vågrät gallerikortsutformning oberoende av vanliga gallerilistor.\\
Utseende & Taggar i galleriets huvudpanel & Sortera efter användning eller alfabetiskt. Visa antingen alla direkt eller 1--200 taggar före kontrollen för utökning på plats; standardtröskeln är 20. Valfri rullning använder 1--12 synliga rader, standard 5, och aktiveras endast när de visade taggarna faktiskt radbryts utöver konfigurerat antal rader.\\
Utseende & Gränssnittsanimationer & Tiden för den offentliga informationspanelen på gallerikort är 0--800~ms (standard 320~ms) och tiden för Admin-sidopanelen 0--800~ms (standard 260~ms). Värdet 0 stänger av rörelsen.\\
Layout & Genvägar i sidhuvudet & Tre ordnade platser. En plats kan vara tom, peka på startsidan eller ett galleri som hittas med galleriväljaren. Offentliga besökare ser bara offentliga och listade gallerimål.\\
Utseende & Gallerikortsutformning & \ui{Lodrätt system} placerar bilden ovanför beskrivningen; \ui{Vågrätt system} bredvid den. Nyinstallationer börjar lodrätt; migreringen bevarar äldre installationers utseende. Ett fysiskt galleri kan ärva eller åsidosätta Temas val där presentationsschemat är tillgängligt.\\
Layout & Bildantalsmärke & Det globala märket för ingående bilder är aktiverat som standard. Enskilda gallerier kan ärva det eller spara en åsidosättning. När det är aktiverat för det öppna galleriet visar huvudpanelen också aktuellt galleris grenantal för bilder, inklusive åtkomliga undergallerier.\\
Layout & Offentlig miniatyrbildsrenderare & \ui{Progressiv skärpning av miniatyrbilder} är standard och säker reserv. \ui{Responsivt webbläsarval} förblir det stödda äldre valet med fullständiga serverrenderade kandidater och beteende utan JavaScript. Omslag och kollage behåller sin responsiva pipeline.\\
Layout & Paginering & Global paginering är avstängd som standard. Kolumner har standardvärdet 3 och begränsas till 1--12; rader har standardvärdet 3 och begränsas till 1--50. Effektivt antal objekt per sida är kolumner multiplicerat med rader.\\
Layout & Startsiderutnät & Startsidan har eget rutnät med 1--12 kolumner och 1--50 rader. Det ändrar inte rutnäten i gallerier.\\
Layout & Återställning av rutnät per galleri & \ui{Återställ alla egna gallerirutnät} rensar databasåsidosättningar, återställer standardtillståndet för undergalleriers arv och tar bort matchande rutnätsnycklar från \filepath{gallery.json}-sidofiler så att en senare upptäcktskörning inte kan återimportera gamla rutnätsinställningar.\\
Layout & Ljuslådestandard & \ui{Enskild bild}, \ui{Bildremsa} eller \ui{3D-karusell} medan funktionen för ljuslådelägen är~aktiverad. Det klassiska enbildsläget är reserv och det enda läge som visas när funktionen är avstängd.\\
Profilering & Offentlig sidhuvudsbanderoll & Valfri banderoll på temanivå i JPG/PNG/GIF/WebP, högst 8~MB. Den ersätter den synliga webbplatstiteln när ingen gallerispecifik banderoll är konfigurerad.\\
Profilering & Sidhuvudsavdelare & Valfri avdelare i JPG/PNG/GIF/WebP, högst 8~MB. Bredd 0 följer den responsiva sidbredden; en positiv bredd begränsas till 160--3840~pixlar. Höjd har standardvärdet 72 och begränsas till 8--512~pixlar. Normalt läge bevarar bildförhållandet och behandlar höjden som ett maximum; sträckläge använder exakt den konfigurerade visningsrutan och kan avsiktligt förvränga bilden.\\
Profilering & Favicon & Ladda upp JPG/PNG/GIF/WebP, dra en kvadratisk beskärning och zooma beskärningen 1--3 gånger. Sparandet skapar PNG-favicon-varianter på 32, 48 och 180~pixlar.\\
Profilering & Temabakgrund & Bevara originaluppladdningen och generera vid behov ett WebP-derivat när GD/WebP-stöd finns. Den optimerade längsta sidan är 1024--3840~pixlar, standard 1920, med reglagesteg på 128~pixlar. Derivatet kan genereras om eller tas bort utan att originalet raderas.\\
Profilering & Bakgrundssynlighet och reserv & Bakgrundssynlighet är 0--100~procent, standard 65. Galleriets reservkälla kan vara ingen, en ny temauppladdning, en befintlig galleribild eller ett kollage genererat från offentliga gallerier. Separata åtgärder återställer alla bakgrundsval per galleri, tar bort temabakgrunden eller tar bort favicon.\\
Egen CSS & Förvalt utseende eller uppladdning & Förvalen är CSS-filerna som för närvarande finns i \filepath{custom_css/}; när ett väljs kopieras det till \filepath{public/assets/custom.css}. En manuellt uppladdad \filepath{.css}-fil ersätter den aktiva filen. Den aktiva egna formatmallen läses in efter inbyggd CSS och sparade temakontroller.\\
Egen CSS & Återställningsåtgärder & \ui{Återställ egen CSS} raderar den aktiva egna formatmallen och rensar dess förvalsmarkör. \ui{Återställ till CSS} rensar sparade visuella åsidosättningar för palett, typsnitt, radie, bakgrundskälla och GPS så att underliggande CSS blir styrande igen; strukturella sidbreddsinställningar bevaras avsiktligt.\\
Språk & Administrationsspråk och offentligt språk & Administrationsspråket lagras för administratörens session och webbläsare; det offentliga språket är webbplatsens anonyma standard. Besökarväljaren är en oberoende besökaråsidosättning som beskrivs nedan.\\
Språk & Paketunderhåll & Stödda paket visar antal strängar och täckning. JSON-redigeraren validerar ett objekt med strängvärden, importerar och exporterar JSON och kan rensa autentiserad diagnostik om saknade nycklar. Engelska förblir käll- och reservkatalog.\\
\bottomrule
\end{longtable}

\subsection{Ljuslådans bläddringslägen}
\index{l00152@ljuslåda!b00043@bläddringslägen}
\index{b00037@bildremsa}
\index{a00000@3D-karusell}

Temaredigeraren definierar den offentliga ljuslådans standardpresentation. \ui{Enskild bild} behåller den klassiska visaren. \ui{Bildremsa} lägger till närliggande miniatyrbildsförhandsvisningar under det aktiva fotografiet. \ui{3D-karusell} placerar en liten uppsättning närliggande fotografier bakom det aktiva för att ge mer visuellt sammanhang under bläddringen.

Ett fysiskt galleri kan ärva Temas standard eller spara en egen åsidosättning av ljuslådeläge. Smarta galleriers presentation kan också bestämma ett effektivt ljuslådeläge utifrån egna presentationsinställningar och webbplatsens standardvärden. Lägena delar samma behörighetskontrollerade pipeline för ljuslådemetadata och navigering, så helskärm, bildspel, zoom, kartor, röstning och progressiv kvalitet fortsätter följa sina vanliga regler. Att stänga av den globala funktionsväxeln för ljuslådelägen döljer dessa bläddringskontroller utan att ändra lagrade fotografier.

\subsection{Välja gränssnittsspråk}
\label{sec:language-selection}
\index{s00230@språk}
\index{l00154@lokalisering!s00231@språkval}
\index{z00282@översättningspaket}

PHP Gallery skiljer programmets gränssnittsspråk från texten du skriver själv. Att ändra gränssnittsspråket ändrar knappar, menyposter, hjälptext, statusmeddelanden, dialoger och annan programtext. Galleri- och bildtitlar och beskrivningar kan ha valfria underhållna språkversioner, men de redigeras och sparas separat; taggar och annat ägarinnehåll förblir oförändrade.

Öppna \ui{Tema > Språk}. Två oberoende väljare finns:

\begin{itemize}
  \item \ui{Administrationsgränssnittets språk} ändrar det språk administratören ser i aktuell webbläsare.
  \item \ui{Offentligt besökarspråk} anger standardgränssnittsspråket för offentliga sidor.
\end{itemize}

Den valfria \ui{Besökarens språkväljare} låter enskilda besökare åsidosätta den offentliga standarden i sin egen webbläsare. Administratören väljer vilka underhållna språk som visas i den offentliga väljaren; om funktionen stängs av återgår besökare till webbplatsens offentliga språk. Administrationsspråk och offentligt språk behöver inte vara samma.

Offentliga sidor använder en kompakt språkkontroll med språkens egna namn och lokalt levererade flaggikoner. Att välja ett språk håller besökaren på samma offentliga sida och kommer ihåg valet för senare besök i den webbläsaren. \ui{Använd webbplatsens standard} tar bort den personliga åsidosättningen. Motsvarande parameter \codeword{?lang=} finns kvar för direktlänkar och användning utan JavaScript.

\subsubsection{Anpassa besökarväljaren}

Väljaren har fem förval: \ui{Klassisk}, \ui{Fyllda kapslar}, \ui{Kontur}, \ui{Mjuka kort} och \ui{Minimal}. \ui{Inställningar > Allmänt} ger grundläggande förvals- och flaggval, medan \ui{Tema > Språk} innehåller detaljerade utformningskontroller och direktförhandsvisning. Dessa utseendeinställningar ändrar inte erbjudna språk, webbplatsstandard, administrationsspråk eller besökarens personliga val.

\subsubsection{Flerspråkig galleri- och bildtext}
\index{l00154@lokalisering!g00101@galleriinnehåll}
\index{l00154@lokalisering!b00041@bildtitlar}

Galleri- och bildredigerare behåller befintlig titel och beskrivning som källinnehåll. Använd \ui{Språk för aktuell titel och beskrivning} för att klassificera texten som engelska, tjeckiska, tyska eller svenska, eller lämna den \ui{Ej angivet}. Befintliga installationer klassificeras inte automatiskt som engelska.

Öppna \ui{Andra språk} med jordglobskontrollen för att ange valfria översatta titlar och beskrivningar. Gallerititlar och beskrivningar faller tillbaka oberoende av varandra. En sparad bildspråksrad behandlas i stället som en enda bildtextvariant: tomma fält förblir dolda i stället för att blanda källspråkstext under översatt bildtext. Om båda översatta fälten lämnas tomma tas språkraden bort. Sparandet sker genom det vanliga galleri- eller bildformuläret och stannar i den befintliga administrationssidopanelen.

Offentliga sidor använder besökarens valda underhållna språk. Matchande text visas på kort, i galleribeskrivningar, bildöverlägg, ljuslådemetadata, sökresultat, alternativtext och SEO-metadata. Saknade gallerifält faller tillbaka på källinnehåll; en befintlig bildtextvariant visar endast sina sparade översatta fält. Översättningar ändrar inte URL:er, URL-namn, mappar, filnamn, ordning, synlighet, lösenord, NSFW-regler eller mediebehörighet.

Om OpenAI-texthjälp är konfigurerad skickar \ui{Föreslå översättning med OpenAI} källfältet och infogar ett granskningsbart utkast i målfältet. Utkastet sparas eller publiceras aldrig automatiskt. Manuell översättning är helt tillgänglig utan leverantör.

\subsubsection{Tillgängliga språk och deras aktuella status}
\index{s00230@språk!e00069@engelska}
\index{s00230@språk!t00248@tjeckiska}
\index{s00230@språk!t00250@tyska}
\index{s00230@språk!s00234@svenska}

Engelska är det kanoniska källspråket, standard vid nyinstallation och reservspråket. Tjeckiska, tyska och svenska underhålls också som fullständiga kataloger. Dessa fyra är de gränssnittsspråk som kan väljas för närvarande.

\begin{longtable}{>{\RaggedRight\arraybackslash}>{\hspace{0pt}}p{0.20\textwidth} >{\RaggedRight\arraybackslash}>{\hspace{0pt}}p{0.10\textwidth} >{\RaggedRight\arraybackslash}>{\hspace{0pt}}p{0.58\textwidth}}
\toprule
\textbf{Språk} & \textbf{Kod} & \textbf{Aktuell roll}\\
\midrule
\endfirsthead
\toprule
\textbf{Språk} & \textbf{Kod} & \textbf{Aktuell roll}\\
\midrule
\endhead
English & \codeword{en} & Fullständig kanonisk källa, standard vid nyinstallation och reservspråk.\\
Čeština & \codeword{cs} & Fullständig underhållen valbar tjeckisk översättning.\\
Deutsch & \codeword{de} & Fullständig underhållen valbar tysk översättning.\\
Svenska & \codeword{sv} & Fullständig underhållen valbar svensk översättning.\\
\bottomrule
\end{longtable}

\important{Endast engelska, tjeckiska, tyska och svenska kan väljas för närvarande. Engelska förblir reserv om en underhållen översättning saknas.}

\subsubsection{Så fungerar översättningspaket}
\index{z00282@översättningspaket!r00204@reserv}
\index{z00282@översättningspaket!j00131@JSON}
\index{z00282@översättningspaket!d00060@diagnostik}

Ett översättningspaket är helt enkelt en ordbok med namngivna gränssnittsmeddelanden. Programmet begär en stabil meddelandenyckel såsom ''save gallery'' eller ''delete selected photos'', och det aktiva språkpaketet tillhandahåller formuleringen som ska visas. Den interna nyckeln är densamma på varje språk; bara den visade texten ändras.

Tabellen \ui{Upptäckta språkpaket} visar täckningen för stödda språk, och underavsnittet \ui{Diagnostik} registrerar saknade nycklar som observeras under en autentiserad administrationssession. Engelska tillhandahåller reservtext när en underhållen översättning saknar en nyckel.

\ui{Paketredigerare} är avsedd för omsorgsfullt underhåll eller import av en förberedd översättning. Språkdata lagras som JSON. I redigeraren är texten på vänster sida i varje JSON-par den stabila nyckeln och ska inte översättas. Texten på höger sida är värdet som får översättas. Platshållare som \codeword{\{count\}}, \codeword{\{gallery\}} eller \codeword{\{time\}} representerar aktuella värden som PHP Gallery infogar och måste behålla samma namn på varje språk.

Ett språkpaket kan också exporteras som JSON, översättas utanför PHP Gallery och importeras igen. Import accepteras endast när filen är ett JSON-objekt med textvärden, så en översättare kan arbeta med paketet utan åtkomst till galleriservern.

För utvecklare och underhållare finns de underhållna katalogerna under \filepath{app/lang/}. JSON är det kanoniska redigerbara formatet. PHP-ordböcker för engelska, tjeckiska, tyska och svenska behålls endast som kompatibilitetsreserver för installationer som saknar motsvarande JSON-fil. Serverrenderade sidor och webbläsar-JavaScript använder samma modell med aktivt språk och engelsk reserv, så en översättning behöver inte dubbleras i separata klient- och serverkataloger.

Egen CSS stöder installationsspecifik visuell finjustering. Temagenererad CSS exponerar validerade variabler och beräknade selektorer genom en särskild rutt. Galleriprofilering kan åsidosätta utvalda resurser på webbplatsnivå samtidigt som effektivt reservbeteende bevaras.

\subsection{Taggar i galleriets huvudpanel}
\index{t00241@taggar!g00099@galleriets huvudpanel}
\index{t00244@tema!g00111@galleritaggar}

Underavsnittet \ui{Utseende > Galleritaggar} styr den kompakta taggsamlingen under titeln i ett öppet galleri. Taggområdet använder hela huvudpanelens bredd, så taggar kan fortsätta radbrytas mot högerkanten i stället för att begränsas av den vanliga lättlästa styckebredden. Direkta galleritaggar och taggar som samlas från ingående innehåll förblir separata innehållsgrupper, men varje grupp följer samma globala visningspolicy.

Samma underavsnitt styr också den offentliga tagglandningssidan. \ui{Gallerier per rad} och \ui{Rader per sida} definierar taggresultatens gallerirutnät och sidkapacitet, medan \ui{Gallerikortsutformning} väljer lodrät eller vågrät kortpresentation för resultaten. Värdena påverkar offentliga sidor som öppnas via en tagg; de ändrar inte vanliga gallerisidor. Om de särskilda värdena inte har sparats förblir Temas globala rutnät och gallerikortsutformning reserv. Den globala pagineringsväxeln avgör fortfarande om ett långt taggresultat delas i sidor. Från \ui{Redigera taggar} öppnar \ui{Konfigurera taggvisning} detta underavsnitt direkt, och \ui{Hantera taggmetadata} ger den omvända sammanhangslänken.

Som standard visar webbläsaren först 20 taggar. Gränsen kan ändras från 1 till 200 med antingen ett reglage eller ett exakt nummerfält. När fler taggar finns visas \ui{Visa alla taggar} och utökar de befintliga taggelementen på plats med JavaScript; ingen sidomladdning eller ytterligare serverbegäran används. Samma kontroll kan fälla ihop samlingen igen. Administratörer som inte vill ha stegvis visning kan aktivera \ui{Visa varje tagg direkt}. Alla taggar finns fortfarande i den serverrenderade HTML-koden i varje läge, så avstängt JavaScript lämnar hela taggsamlingen synlig och användbar i stället för att permanent dölja innehåll.

Taggordning kan konfigureras som \ui{Mest använda först} eller \ui{Alfabetiskt}. \ui{Mest använda först} är standard. Användning är det totala antalet direkta taggtilldelningar till gallerier plus direkta taggtilldelningar till fotografier i installationen. Lika användningsantal faller tillbaka på naturlig skiftlägesokänslig alfabetisk ordning och därefter en stabil identifierarregel. Direkta galleritaggar sorteras inom sin grupp och ingående taggar inom sin grupp; grupperna blandas inte.

Långa taggsamlingar kan också använda en responsiv intern rullningslist. Alternativet är aktiverat som standard och har fem synliga rader som standard. Radtröskeln kan konfigureras från 1 till 12 med ett reglage eller exakt nummerfält. Webbläsaren mäter raderna efter att taggarna faktiskt har radbrutits vid aktuell huvudpanelsbredd och aktiverar rullning endast när det uppmätta radantalet överstiger konfigurerad tröskel. Storleksändring av sidan räknar om måttet. Att stänga av rullningslistalternativet tar bort höjdbegränsningen och låter huvudpanelen växa naturligt.


\section{Hämtningar, delning och överföring}
\label{sec:sharing}
\index{h00122@hämtningar}
\index{z00273@ZIP}
\index{w00271@WebDAV}

Gallerihämtningar förbereder behörigt innehåll som ZIP-arkiv och strömmar det färdiga arkivet till webbläsaren. En vanlig gallerihämtning går igenom det tillåtna galleriunderträdet; hämtningar från smarta gallerier utvärderar den sparade frågan och aktuella åtkomstregler på nytt; urvalshämtningar i Bildhanteraren paketerar bara aktuellt valda tillhörande fotografier; administratören kan också begära ett arkiv med alla gallerier. Arkivposternas namn normaliseras för att förhindra sökvägstraversering och behålla deterministiska sökvägar.

I en modern webbläsare begär en hämtning från ett offentligt galleri eller smart galleri först ett privat, begränsat manifest. Webbläsaren begär sedan varje behörigt original genom en oberoende medieadress och bygger ZIP-filen lokalt, visar förlopp och tillåter avbrott eller återförsök utan att hålla en lång PHP-arkivbegäran öppen. Varje originalbegäran upprepar aktuella kontroller av galleri, medlemskap i smart galleri, besökarsynlighet, medieversion och sökvägens inneslutning; en bildidentifierare ensam ger aldrig åtkomst. Gränser för filantal, sammanlagt originalbyteantal, dubbla namn, minne och ZIP64 håller arbetsflödet begränsat. Direktlänkar och klienter utan JavaScript behåller den befintliga serverbaserade ZIP-reserven.

Genererade ZIP-filer för gallerier, alla gallerier och urval använder en innehållssignatur som härleds från relevant galleri- och bildmetadata, så att ett fortfarande aktuellt cachelagrat arkiv kan återanvändas och en ändrad samling får en ny cacheidentitet. Denna \emph{innehållssignatur} är metadata för cacheogiltigförklaring, inte en offentlig HMAC-token för hämtningsbehörighet. Hämtningskontrollern kontrollerar behörigheten innan ett arkiv levereras.

\subsection{Migrering mellan gallerier}
\index{g00103@gallerimigrering}
\index{a00015@API-migrering}

Galleriredigerarens API-område stöder båda överföringsriktningarna mellan kompatibla PHP Gallery-installationer. \ui{Exportera detta galleri till en annan instans} behandlar aktuellt galleri som källa och använder en galleribegränsad API-nyckel som skapats för det mottagande övergalleriet. \ui{Importera till detta galleri från en annan instans} behandlar aktuellt galleri som mottagande övergalleri och använder en API-nyckel som skapats för fjärrkällgalleriet. Importen skapar ett nytt undergalleri under övergalleriet; den skriver aldrig över eller flyttar valt mottagande galleri.

\ui{Inkludera undergallerier} är aktiverat som standard. När det väljs beskriver det versionshanterade manifestet källgalleriet och hela dess underhierarki med övergallerier först. Målet återskapar hierarkin under den nya importerade roten och bevarar stödda galleri- och bildmetadata, översättningar, flygkartstillstånd, original, befintliga miniatyrbilder och galleriprofileringsresurser. Att stänga av alternativet överför bara valt källgalleri. Den aktuella kompatibilitetspolicyn kräver att källa och mål kör exakt samma programversion.

Den mottagande servern skapar en deterministisk paketplan begränsad av dess uppladdningsgränser. Varje original hålls samman med sina redan genererade miniatyrbilder, medan resurser på gallerinivå kan paketeras oberoende. ZIP-poster använder stabila resursnycklar, lagrar befintliga byte utan att komprimera om fotografier och kontrolleras mot manifestets storlekar och SHA-256-kontrollsummor före installationen. Exakt API-nyckelomfattning, galleriträdsstruktur, målsökvägar, databasberedskap och resursernas fullständighet valideras vid respektive ändringsgräns.

Administratören väljer ett återanslutningsintervall från 5 till 300~sekunder. Efter tidsgräns, omladdning eller avbruten anslutning frågar den aktiva sidan det beständiga måljobbet vilka resursnycklar som redan finns och hoppar över färdiga paket i stället för att blint skicka dem igen. Mappkopplingar och mottagna resurser sparas stegvis så att samma överföring kan återupptas säkert. Slutförande är ett separat steg och nekas tills varje angiven resurs finns. Att avbryta webbläsarflödet stoppar fortsatta begäranden; inloggningsuppgifter bör ändå återkallas när en tillfällig migreringsnyckel inte längre behövs.

\subsection{Mobil WebDAV-överföring}
\index{w00271@WebDAV!m00168@mobil uppladdning}

\ui{Inställningar > Uppladdningar} skapar en WebDAV-kompatibel anslutning för ett valt galleri. Ange en igenkännbar etikett, välj galleriet, skapa anslutningen och kopiera genererad server-URL, användarnamn och lösenord till mobilappen. Lösenordet visas endast en gång och lagras som en hash. Befintliga rader visar senaste användningstid och kan raderas omedelbart från administrationen.

Använd måltypen WebDAV i klienter som PhotoSync. Aktuell serverinläsning accepterar JPG, PNG, GIF och WebP, så aktivera HEIC-till-JPEG-konvertering i klienten där den finns. WebDAV-uppladdningar kringgår inte galleriets inläsningsregler: uppgiften är galleribegränsad, autentisering krävs vid varje skrivning, galleri- och bildschemat måste vara känt och accepterade medier skannas sedan, döps eventuellt om automatiskt och behandlas genom den vanliga miniatyrbildspipelinen.

\subsection{Offentlig upptäckt och verktyg för valda fotografier}
\index{r00207@robots.txt}
\index{s00221@sitemap.xml}
\index{b00042@bildwebbplatskarta}
\index{b00033@bildhanterare}

Det offentliga svaret \codeword{/robots.txt} beskriver gränser för sökrobotar, medan \codeword{/sitemap.xml} listar lämpliga offentliga galleri- och bild-URL:er med bildmetadata och information om senaste ändring. Opublicerat, privat, lösenordsskyddat, åldersbegränsat eller annars otillåtet innehåll görs inte offentligt genom dessa utdata. Kanoniska URL:er, sociala metadata och strukturerade bilddata använder samma åtkomstmedvetna offentliga sökvägsmodell.

Bildhanteraren är ett överlägg för inloggade administratörer på den offentliga gallerivyn, inte en besökarfunktion. Administratören kan välja fotografier och fysiska undergallerier, radera ett blandat urval eller kopiera det till ett nytt fysiskt galleri under ett valt övergalleri. Urval med enbart fotografier kan också dras till ett synligt undergalleri, flyttas eller kopieras genom den sökbara destinationsväljaren, delas genom enhetens inbyggda delningspanel där den stöds eller hämtas. Inbyggd delning faller tillbaka på ZIP för valda fotografier när webbläsaren inte kan dela de förberedda filerna. Smarta gallerier utesluts från fysiska underträdsåtgärder. Urvalet självt är bara webbläsartillstånd och ger ingen behörighet. Varje ändrings- och ZIP-begäran kontrollerar på nytt vanligt inloggat konto, CSRF-token, källgalleri, vald tillhörighet och destination där det är relevant. Vanliga besökare får inte dessa kontroller. Bildhanteraren är avsiktligt oberoende av funktionen Administration direkt på offentliga sidor: om den senare stängs av tas direkta kontroller på offentliga sidor för redigering, tillägg, radering och omordning bort, medan en oberoende aktiverad Bildhanterare fortfarande kan erbjuda sina flöden för urval, radering, flyttning, kopiering, skapande av fysiska gallerier och urvalsexport.
\section{Underhåll, uppdateringar och återställning}
\label{sec:maintenance}
\index{u00251@underhåll}
\index{u00253@uppdateringar}
\index{s00238@säkerhetskopiering}

\subsection{Migreringar}

Schemaändringar ligger i tidsstämplade PHP-filer under \filepath{database/migrations/}. Migreringsköraren förvaliderar väntande definitioner, tillämpar dem i filnamnsordning och registrerar lyckade versioner. Nytt schemaarbete får en ny migreringsfil så att befintliga installationer behåller en deterministisk historik.\footnote{Vissa migreringar innehåller reparationsfunktioner utöver SQL. Kompatibilitetstester skyddar partiella uppdateringsfall där en äldre körare möter nyare migreringsdefinitioner.}

Kontroller av databasberedskap sparas bara för aktuell PHP-begäran. Det minskar normalt upprepade metadatafrågor, men en migrering kan ändra svaret medan samma administrations-, installations-, uppdaterings- eller kommandoradsprocess fortfarande körs. Migreringsköraren rensar därför de sparade svaren efter schemaändrande satser och efter lyckade reparationsfunktioner. En kontroll efter migreringen läser då den nya databaskatalogen i stället för att återanvända ett svar från före migreringen. Vanliga migreringssatser som enbart ändrar data rensar inte beredskapscachen eftersom de inte kan skapa eller ta bort en tabell, ett fält eller ett index.

Version~0.106 lade till tabellen \codeword{maintenance\_jobs} genom en vanlig migrering. Kör den som vanligt från galleriadministrationssidan om den inte redan har tillämpats. Den kräver installationens vanliga behörighet att skapa tabeller, men kräver inte databasutlösare, lagrade rutiner, \codeword{SUPER}, ändring av en global servervariabel eller separat ingripande av en databasadministratör. Tills migreringen är tillgänglig och verifierad vägrar Underhållscentret starta ett jobb i stället för att försöka underhålla utan beständiga kontrollpunkter. Migreringarna för kooperativa gallerier och miniatyridentiteter använder samma normala uppgraderingsflöde.

\subsection{Så kontrollerar galleriet databasberedskap}
\index{d00049@databas!b00023@beredskap}
\index{s00214@schemaundersökning}
\index{m00163@migreringar!s00210@saknade objekt}
\index{s00236@Systemstatus!d00053@databasundersökning}
\index{f00077@felsökning!d00050@databasanslutning}

De gemensamma beredskapstillstånden och ägarens åtgärder definieras i avsnitt~\ref{sec:db-readiness-contract}. Underhåll använder samma kontrakt när det kontrollerar tabeller, fält och index som behövs för lösenordsåterställning, NSFW-skydd, uppladdningsautomatisering, miniatyrbildsmetadata och andra schemaberoende funktioner.

Vid en uppgradering kan \ui{Saknas} vara väntat när nya programfiler kommer innan deras migrering har körts. \ui{Okänt} betyder däremot att PHP Gallery inte kunde undersöka vald databas tillförlitligt. Vid \ui{Saknas} säkerhetskopierar du och tillämpar väntande migrering. Vid \ui{Okänt} kontrollerar du databasanslutning, valt schema och behörigheter för metadataundersökning innan schemat ändras.

Efter reparationen laddar du om \ui{Systemstatus} och testar berörd funktion igen. Spara eventuell begärandereferens tillsammans med tid och programversion när du ber webbhotellet eller en underhållare om hjälp.

\subsection{Galleripapperskorg och återställning}
\index{g00096@galleri!p00183@papperskorg}
\index{g00096@galleri!z00277@återställning}
\index{p00183@papperskorg!a00019@automatisk rensning}

Öppna \ui{Underhåll > Papperskorg} för att granska återställningsbara galleriraderingar. Varje aktiv rad visar ursprunglig sökväg, raderingstid, galleri- och bildantal, lagrad storlek, livscykelstatus och tidsgräns för automatisk rensning när policyn är aktiv. Problemposter förblir synliga för undersökning; PHP Gallery gissar inte att ett tvetydigt filsystems- och databastillstånd är säkert att förstöra.

Papperskorgens huvudväxel styr framtida galleriraderingar och är~aktiverad som standard. Att stänga av den döljer eller raderar inte befintliga poster: återställning och manuell permanent radering finns kvar. Automatisk rensning är ett separat aktivt val och är avstängd som standard. Dess lagringstid är som standard 30~dagar, underhållsomgången är begränsad och automatisk radering pausas när Papperskorg är avstängd. Om automatisk rensning aktiveras efter en paus får aktuella återställningsbara poster en ny full lagringsperiod innan schemalagd radering kan börja.

Migreringarna \filepath{202609070001_gallery_trash_bin.php} och \filepath{202609070002_gallery_trash_state_machine.php} lägger till beständiga åtgärdsposter. När Papperskorg är~aktiverad kontrollerar raderingen detta schema före första filsystemsflytten eller databasradraderingen. Bekräftat saknad eller okänd lagring stoppar återställningsbar radering och hänvisar administratören till att tillämpa eller kontrollera migreringar; den faller inte tyst vidare till permanent radering. Schemalagt webbplatsunderhåll stämmer försiktigt av gamla pågående åtgärder och kör utgången automatisk rensning endast när både Papperskorg och automatisk rensning är~aktiverade.

Papperskorgsinnehåll ligger i skyddad beständig \filepath{data/gallery-trash/}-lagring, inte i cachen eller det aktiva galleriträdet. Driftsättningspaket innehåller mappens åtkomstpolicyfil men utesluter lagrat papperskorgsinnehåll. Ta med denna körningsmapp och databasen i planeringen för säkerhetskopiering och återställning; ingen av halvorna är ensam en fullständig återställningsbar post.

\subsection{Underhållscenter}
\index{u00252@Underhållscenter}
\index{u00251@underhåll!c00044@centralt arbetsflöde}

Öppna \ui{Underhåll > Underhållscenter} för ett väglett flöde med \ui{Analysera}, \ui{Granska}, \ui{Utför} och \ui{Verifiera} över vanliga underhållsansvar. Analysera är skrivskyddat för galleri-, telemetri-, logg-, papperskorgs-, cache-, miniatyrbilds- och databasinnehåll; det skriver bara den beständiga kontrollpunkt för Underhållscentret som behövs för att återuppta undersökningen. Granska-steget grupperar uppgifter efter område, märker obligatoriska automatiska steg och placerar tillgängliga valfria uppgifter utan upptäckt arbete i en infälld grupp. Obligatoriska, förvalda och valfria uppgifter förblir tydliga före rensning; djup medieverifiering och fullständig databasoptimering går fortfarande att hitta när de är tillgängliga, även om de inte väljs automatiskt.

Den granskade planen binds till aktuell programversion, tillämpad migreringsrevision, Underhållscentrets registerrevision och planinnehåll. Om någon av dessa indata ändras vägrar körningen den gamla planen och begär en ny analys. Servern kontrollerar uppgiftsberoenden, funktionstillstånd och databastabellers lämplighet igen vid körningen; webbläsarkryssrutor är presentationsindata, inte behörighet.

Körningen går framåt genom begränsade kontrollpunkter, så du kan ladda om eller stänga sidan och återuppta samma jobb senare. Bara ett körande eller pausat centralt jobb äger den beständiga ändringsreservationen, och två webbläsarflikar kan inte driva ett jobb framåt samtidigt. Automatiskt webbplatsunderhåll avstår före sin första ändring medan reservationen hålls. Offentliga skrivskyddade galleribegäranden förblir tillgängliga.

Pausa och Avbryt samarbetar med körningen: de får effekt mellan begränsade åtgärder och ångrar inte redan genomfört arbete. En databasåtgärd \codeword{ANALYZE TABLE} eller \codeword{OPTIMIZE TABLE} som redan körs måste återvända innan avbrottet kan fortsätta. Centret behandlar högst en för närvarande lämplig tabell per webbläsarsteg, hoppar över skyddade, okända, överstora eller nyligen olämpliga tabeller och registrerar ett osäkert atomärt utfall i stället för att blint köra en tabell igen efter ett förlorat svar.

Slutrapporten skiljer mellan borttagna rader och filer, analyserade eller optimerade tabeller, överhoppningar, varningar och fel. Den rapporterar fysisk lagring före och efter endast när inventeringsunderlaget stöder slutsatsen. Det generiska arbetsflödet tillämpar aldrig migreringar, gör schemareparation, installerar uppdateringar, raderar administrationsloggarkiv, tömmer alla papperskorgsposter eller genererar om medier; använd specialsidorna för dessa avsiktliga åtgärder.

Rensning av hämtningar och cache går genom manifestcachen, äldre hämtningsartefaktcache och genererad ZIP-cache som separata begränsade steg. Dessa cacher är oberoende optimeringar: om en är otillgänglig på aktuellt webbhotell bevarar centret redan slutförd rensning, registrerar en säker varning och fortsätter med övriga ansvarsområden och senare uppgifter. Varningen kan identifiera åtgärden och en stabil felkategori men avslöjar inte råa undantagsmeddelanden, SQL, inloggningsuppgifter, privata sökvägar eller stackspår. Efter att en uppdatering ändrat detta körningskontrakt är en äldre färdig plan avsiktligt inaktuell; kör Analysera igen innan du fortsätter.

\subsection{Miniatyrbildsunderhåll}

Underhållsverktyg undersöker genererade varianter, identifierar saknade eller gamla derivat och bygger om valda storlekar. Webbläsarstödd ombyggnad kan förbereda originaldelar där det är konfigurerat. Bakgrundsförvärmning hanterar små luckor i offentlig rendering med begränsade begäranden, medan uttryckligt administrationsunderhåll förblir huvudvägen för storskalig reparation.

\subsection{SEO-begärandeskydd och offentliga frågesträngar}
\index{s00215@SEO-begärandeskydd}
\index{f00086@frågesträngsspam}
\index{k00132@kanonisk URL}

Underhållsområdet innehåller ett valfritt offentligt SEO-begärandeskydd, aktiverat som standard. För anonyma offentliga GET-begäranden håller det en uttrycklig lista över tillåtna frågeparametrar per rutt. Kända analysparametrar såsom UTM-fält och vanliga annonsklickidentifierare ignoreras eftersom de inte ändrar renderat innehåll. Oväntade parametrar avvisas före den vanliga offentliga renderaren med 404, \codeword{noindex,nofollow} och ett svar utan cachelagring så att sökrobotgenererat parameterspam inte mångfaldigar indexerbara galleri-URL:er.

Administrationsrutter och maskinadresser såsom installation, WebDAV, uppladdningsautomatisering, mottagare för gallerimigrering och underhållets cron-adress ligger utanför detta offentliga frågeskydd. Inloggade administratörer är också undantagna så att diagnostik- och administrationsflöden inte blockeras av offentliga sökrobotregler. Loggning av avvisade frågor kan aktiveras oberoende och har en provtagningsgräns per UTC-dygn för att förhindra att en sökrobotöversvämning fyller administrationsloggen. Samma tjänst ger kanoniska reserv-URL:er när en offentlig sida inte redan har skickat en egen kanonisk tagg.

\subsection{Databasunderhåll}

Databasundersökning inventerar schemaobjekt, referenser, ansvarskategorier, lagringspolicy och tillförlitligheten för rensning. Logisk rensning använder uttryckligen tillåtna regler med deterministiskt urval, begränsade omgångar och granskningsposter. Schemareparation ger en provkörning och en idempotent migreringsväg. \codeword{ANALYZE TABLE} kan uppdatera optimerarstatistik för valda tabeller. Fysisk \codeword{OPTIMIZE TABLE} finns bara som förhandsvisning följd av en separat bekräftad åtgärd för valda tabeller; den är aldrig ett automatiskt rensningssteg.\footnote{Lagringsmotorns allokeringsvärden såsom uppskattningar av ledigt datautrymme beskriver motorns tillstånd och garanterar inte ett exakt antal filsystemsbyte som frigörs efter optimering.}

\subsection{Lagringsstatistik och fullständig rapport}
\index{l00147@lagringsstatistik}
\index{f00087@fullständig rapport}
\index{g00107@gallerirapport}

Detaljerad \ui{Lagringsstatistik} beräknas bara när administratören begär den, så att den vanliga översikten inte gör en fullständig skanning av derivatfilsystemet vid varje besök. Filvyn skiljer indexerade originalfotografier från genererade JPEG/WebP-miniatyrbilder och DNG-visningsmaster för webbläsaren och visar fördelningar av originalfiltyper och storlekar, största gallerier och original, okända storlekar och skanningsvarningar. En cachelagrad ögonblicksbild markeras som inaktuell när galleridata ändras. Uppdatering av ögonblicksbilden skannar förväntade genererade filer i små webbläsardrivna omgångar. Samma område länkar till databaslagringsdetaljer och databasunderhåll.

Databasfliken läser lagringsstatistik från informationsschemat separat från filsystemsskanningen. Den skiljer Gallery-ägda tabeller från hela den valda databasen där det går och visar de största tabellerna i stället för att tvinga in en fullständig tabellundersökning i den vanliga översiktsbegäran. \ui{Uppdatera allt} uppdaterar filskanningen, förnyar tabelluppskattningar med begränsade \codeword{ANALYZE TABLE}-begäranden och kör sedan den skrivskyddade databasinspektionen. Förloppet sparas i ett serverägt arbetsflöde och den visade fliken uppdateras när stegen är klara utan att sidan lämnas; en ny start under samma administratörs aktiva arbetsflöde fortsätter från det sparade steget. Partiella fel rapporteras per steg och skrivs till administrationsloggen. Åtgärden rensar inte data, reparerar inte scheman, optimerar inte tabeller och kör inte en plan från Underhållscentret.

\ui{Underhåll > Fullständig rapport} bygger en fristående HTML-rapport som omfattar lagring, databasstruktur och storlek, galleristruktur, bild- och EXIF-statistik, GPS-översikt, telemetri, driftloggar, funktions- och inställningstillstånd samt körningsdiagnostik. Bild- och genererat mediearbete behandlas i begränsade webbläsardrivna omgångar så att rapporten förblir praktisk på delade webbhotell. Tillfälliga tidsgräns-, frekvensbegränsnings- och serverfel försöks igen med korta ökande fördröjningar, medan EXIF- och GPS-sammanfattningar med många olika värden förblir begränsade för att styra sessionslagring och minnesanvändning i stora bibliotek. Vald telemetriperiod påverkar bara telemetriavsnittet; bildstatistik omfattar alltid hela bildbiblioteket. GPS-platssammanfattningar använder lokal ungefärlig klustring och offline-referensdata; sannolika simulator- eller spelskärmbilder utesluts från verkliga platsfrekvenser där tillgängligt underlag stöder den klassificeringen.

Den färdiga rapporten returneras till webbläsaren som en hämtning och sparas inte som rapportfil på servern. För att behålla en PDF-kopia öppnar du den genererade HTML-filen och använder webbläsarens Skriv ut eller Spara som PDF; rapporten har utskriftsformatering. Behandla exporterade rapporter som administrativa data eftersom de kan innehålla detaljerad drift- och samlingsstatistik.

\subsection{Telemetri och integritetskontroller}
\index{t00243@telemetri}
\index{i00127@integritet!t00243@telemetri}

Det valfria området \ui{Telemetri} registrerar integritetskontrollerade driftmätningar och presenterar dem som administrationsrapporter. Tillgängliga kontroller omfattar huvudväxeln för telemetri, offentliga användningsmätningar, prestanda, cache- och databashändelser, respekt för Do Not Track, uteslutning av administratörer, maximal räknad bildvisningstid och lagringstid för rådata, timdata och dygnsdata. Rapporter innehåller mått såsom anonyma sessioner, sidvisningar, öppnade fotografier och begränsad visningstid, klientfel, uppmätta mediebyte, webbläsarfördelning och cache- och databasaktivitet där kategorierna är~aktiverade.

Telemetridelsystemet är utformat för att inte lagra råa IP-adresser, råa user-agent-strängar från webbläsare, råa hänvisnings-URL:er, namn, e-postadresser, kontoidentifierare, begärandekroppar eller exakta besökarplatser. Exporter innehåller sammanställd anonym telemetri och sessionshashar i stället för råa besökaridentifierare. Underhåll sammanställer och rensar telemetri enligt konfigurerad lagringspolicy, och administrationsskärmen kan exportera en fristående HTML-rapport för offline-granskning. Granska integritetskontrollerna innan du aktiverar insamling på en offentlig webbplats.

\subsection{Körningsdiagnostik, integritet och galleriprestandatest}
\index{k00145@Körningsdiagnostik}
\index{i00128@integritetskontroll}
\index{g00105@galleriprestandatest}

\ui{Körningsdiagnostik} ger en kopierbar miljö- och supportvy för PHP, databasberedskapsklasser, webbhotellsgränser, GD, Imagick/ImageMagick och relevant medieformatstöd. Den visar också den tidigare beskrivna DNG-konverteringspolicyn. Beredskapspanelerna skiljer presentations- och rapporteringsberoenden, autentiserings- och säkerhetsberoenden samt ändringskänsliga beroenden så att administratören kan skilja en valfri saknad funktion från ett tillstånd som måste blockera en skrivning.

Samma avgränsade modell rapporterar beredskap för galleriets redigeringsrevision och olösta bildflyttar. Ett saknat eller okänt obligatoriskt schema är ett åtgärdstillstånd och stoppar berörd ändring innan målfiler eller rader ändras. Sammanfattningar av väntande flyttar innehåller bara begränsade åtgärds- och galleriidentiteter samt säkra tillståndskategorier; använd det dokumenterade kommandoradsflödet för återställning av en uttrycklig åtgärd i stället för att förvänta dig att diagnostiksidan flyttar filer.

Den separata växeln \ui{Utvecklingsdiagnostik} är endast för administratörer. När den aktiveras för en inloggad administratör lägger den till mätinstrumentering i visaren för förladdning, cache, minne, nätverk och bildrutetider i den offentliga upplevelsen, ljuslådan och helskärmen. Den aktiverar inte diagnostiken för vanliga anonyma besökare. Låt den vara avstängd vid vanlig drift om du inte undersöker rendering eller prestanda; galleriprestandatestet beror också på detta utvecklingsdiagnostiksammanhang.

Administrationens \ui{Integritet}-flöde kontrollerar integritetshanterade kärnfiler mot det installerade manifestet och hjälper till att identifiera ofullständiga eller oväntade driftsättningsändringar. Det är ett verifieringsverktyg vid körning; samordnade säkerhetskopior behövs fortfarande eftersom ett integritetsmanifest inte innehåller användarfotografier eller ersätter återställningsmedier.

När utvecklingsdiagnostik är tillgänglig kan en inloggad administratör starta ett \ui{Prestandatest för offentlig galleriinläsning} för aktuellt galleri. Webbläsaren läser in galleriet flera gånger i en dold ram med samma ursprung och anonym förhandsvisning, kombinerar PHP-renderingsräknare med webbläsarens navigeringstider och erbjuder sedan testloggen som JSON. Använd samma galleri, presentationsläge och nätverks- och cacheförhållanden när du jämför två ändringar; testet är mest användbart som en före- och eftermätning snarare än ett universellt poängtal.

\subsection{Schemalagt webbplatsunderhåll}
\index{s00213@schemalagt underhåll}
\index{c00046@cron!w00270@webbplatsunderhåll}

Vanligt underhåll kan köras i begränsade delar i stället för en lång begäran. Underhållstjänsten länkar säkra webbegäranden när trafik finns och kan också anropas från webbhotellets cron med CLI-kommandot som visas i administrationen. Den kan kontrollera miniatyrbildstillstånd, uppdatera metadata, behandla rensningsuppgifter, sammanställa telemetri och arkivera lämpliga administrationsloggdagar enligt deras konfigurerade policyer. Den dolda webb-cron-token kan roteras; rotation ogiltigförklarar tidigare publicerade webb-cron-URL:er. Underhåll är icke-destruktivt som standard och ersätter inte uttrycklig bekräftelse för databasoptimering eller andra åtgärder med betydande följder.
\subsection{Administrationsloggens historik och arkiv}
\index{a00004@administrationsloggar}
\index{l00153@loggarkiv}

Dagligt administrationsarbete skapar användbar historik: uppladdningar, galleriändringar, säkerhetsbeslut, underhållsresultat, varningar och andra drifthändelser kan registreras i administrationsloggen. Historiken är värdefull när du vill förstå vad som hände, när det hände och vilket konto eller vilken begäran som var inblandad. Det är också normalt att tabellen blir stor i en aktiv installation. Version~\version{} behandlar därför loggskötsel som en egen uttrycklig underhållsuppgift i stället för att låta generell databasrensning oväntat ta bort granskningshistorik.

Administrationsloggskärmen kan visa enskilda poster eller samla upprepade händelser i en sammanfattning. En grupp kan representera upprepade miniatyrbildsvarningar, uppladdningsresultat eller underhållsmeddelanden; när den öppnas visas de enskilda posterna. Filter och kontroller för arbetsflödesstatus hjälper till att avgränsa driftgranskningen, medan begränsade exporter bevarar vald historik utan att läsa in hela tabellen i en begäran. Stora listor, utökning till enskilda råposter och exporter behandlas i begränsade delar, så att en begäran om fullständig historik inte kräver att webbläsaren eller PHP håller hela tabellen i minnet samtidigt.

Ägaren kan välja hur länge nya loggar ska finnas kvar i den aktiva databasen. \emph{För alltid} stänger av automatisk arkivering. En begränsad aktiv lagringsperiod kastar inte omedelbart bort äldre poster. Underhållsflödet väljer först hela kalenderdygn utanför den aktiva lagringsperioden, skriver ett skyddat arkiv med både maskinläsbar JSON och en läsbar statisk HTML-rapport, kontrollerar att arkivet innehåller förväntat antal poster och tar först därefter bort samma rader från den aktiva tabellen. Arkivet lagras under \filepath{data/admin-log-archives/}, som är programdata snarare än ett offentligt galleriområde och har en ytterligare skyddsregel för webbservern.

Arkivunderhåll kan återupptas. Det skriver tillfälliga filer innan ett färdigt arkiv tas i bruk, registrerar en liten tillståndsfil och använder ett lås så att två underhållsbegäranden inte behandlar samma historik samtidigt. Om webbhotellsgränser, ett anslutningsavbrott eller ett annat problem stoppar arbetet kan administrationsskärmen visa det ofullständiga tillståndet, och nästa underhållsförsök kan verifiera, fortsätta eller återställa det. Ett arkiv betraktas aldrig som skäl att radera aktiva poster enbart för att en fil med liknande namn finns; kontroller av radantal och arkivmetadata ingår i säkerhetsgränsen.

Administratörer bör fortfarande ta med databasen och mappen \filepath{data/admin-log-archives/} i samordnade säkerhetskopior. Den aktiva databasen innehåller ny drifthistorik, medan arkivmappen innehåller äldre historik som valts enligt den uttryckliga lagringspolicyn. Att behålla båda ger ägaren en obruten historik och gör en återställd installation lättare att undersöka efter en uppdatering, migrering eller oväntad händelse.
\subsection{Fullständig referens för aktuella funktioner}
\label{sec:complete-feature-reference}
\index{f00090@funktionsreferens}
\index{f00089@funktionsförteckning}

Denna referens är manualens fullständighetschecklista för version~\version{}. Den listar aktuella produktfunktioner som exponeras genom offentliga rutter, administrationsnavigering, galleri- och bildredigerare, underhållsverktyg och valfria integrationer. Interna hjälpfunktioner listas inte som separata funktioner när de endast finns för att genomföra en rad nedan.

\begin{longtable}{>{\RaggedRight\arraybackslash}>{\hspace{0pt}}p{0.23\textwidth} >{\RaggedRight\arraybackslash}>{\hspace{0pt}}p{0.66\textwidth}}
\toprule
\textbf{Funktion} & \textbf{Implementerat beteende}\\
\midrule
\endfirsthead
\toprule
\textbf{Funktion} & \textbf{Implementerat beteende}\\
\midrule
\endhead
Filsystembaserade gallerier & Nästlade fysiska mappar som speglas av databasposter, praktiskt obegränsad hierarki, upptäckt och import, manuellt skapande, fysiska flyttar, rena offentliga sökvägar, sidofilsmetadata och stabil syskon- och bildordning.\\
Gallerimetadata & Titel, beskrivning, underhållna översättningar, manuellt datum eller datumperiod, EXIF-datumförslag, URL-namn, mappnamn, övergalleri, ordning, taggar, omslags- eller titelbild, uppladdad galleriminiatyrbild, profilering, bakgrund och presentationsåsidosättningar.\\
Galleriåtkomst & Offentligt, opublicerat, privat, lösenordslås, delningslänkar med eller utan utgångstid, ärvd åtkomstbedömning, administratörsåtkomst och bekräftelsegränser för NSFW / 18+.\\
Massadministration av gallerier & Grenmedveten radering, synlighetsuppdateringar, ändringar av EXIF/GPS-kartåsidosättningar, röstning, filnamnsvisning och Bildspelet, sökvägsomgenerering, miniatyrbilds- och skanningsåtgärder samt flöden för syskonomordning och datumsortering.\\
Smarta gallerier & Sparade nästlade booleska frågor över fysiska gallerier, taggar, datum, EXIF/GPS, text, AI-metadata, dubblettstatus, filegenskaper och privata redaktionella betyg; offentlig, olistad, rot- eller kopplad placering; ordning ovanför eller nedanför per övergalleri; cykelskydd; oberoende presentationsåsidosättningar; fullständig ljuslådebläddring över sidor; begränsad ZIP-hämtning.\\
Bildredigerare & Titel, beskrivning och översättningar, synlighet, NSFW, sorteringsordning, privat redaktionellt betyg, taggar, miniatyrbildsgränser per fotografi, EXIF/GPS-granskning, granskning av interna AI-metadata och OpenAI-utkastverktyg.\\
Massadministration av fotografier & Flytt till befintligt eller nytt galleri, radering, omslagsval, synlighet, NSFW, miniatyrbildsskapande och ordning genom dragning eller filnamn.\\
Taggar & Delade galleri- och bildtaggar, normaliserade återanvändbara taggnamn och URL-namn, offentliga beskrivningar, användningsantal, viktade redigerarförslag, namnbyte, redigering och radering i administrationen, offentliga tagglandningssidor, taggspecifika rutnäts- och kortinställningar samt ordning, stegvis visning och rullning för huvudpanelens taggar.\\
Offentliga start- och gallerisidor & Hierarkiska kort, brödsmulor, omslag och kollage, paginering, direkta bildrutnät, galleriantal, datum, beskrivningar, profilering och bakgrunder, favoritgenvägar i sidhuvudet, sammanhangsberoende administrationskontroller för redigering, borttagning och omordning för inloggade administratörer, flytande hjälp tillbaka till toppen för långa listor, responsiva offentliga sökvägar och frågesträngsreserv när omskrivning är avstängd.\\
Offentlig sökning & Valfri webbplats- eller grensökning i galleri- och bildtitlar, beskrivningar, filnamn, taggmetadata, underhållen språktext och tillgänglig intern sökbar AI-text; minsta frågelängd och begränsat resultatantal; åtkomst- och listningspolicy tillämpas innan resultat returneras.\\
Robotregler, webbplatskarta och SEO & Åtkomstmedveten \codeword{robots.txt}, bildmedveten \codeword{sitemap.xml}, kanoniska, sociala och strukturerade metadata samt det ruttspecifika frågesträngsskydd som beskrivits ovan.\\
Offentlig mediebehörighet & Särskilda rutter för original, miniatyrbilder, omslag, profilering, bakgrund och favicon utvärderar galleri- och bildsynlighet samt skyddad åtkomst på nytt i stället för att exponera godtyckliga filsystemssökvägar.\\
Miniatyrbildsgenerering & Responsiva storleksvarianter, JPEG/WebP-kompatibilitetspolicy, responsiva och progressiva renderare för valt galleri, lat inläsning och beteende nära visningsområdet, storleksgränser per galleri och fotografi, DNG-visningsmaster, metadataavstämning, ombyggnadsflöden i administration och webbläsare samt signerad åtkomstkontrollerad offentlig reservförvärmning med liten budget.\\
Ljuslåda & Behörighetskontrollerade asynkrona metadata, tangentbords- och peknavigering, bildspel och helskärm, bläddringslägen med enskild bild, bildremsa eller 3D-karusell, EXIF- och kartöverlägg, röstningstillstånd, zoom på 100--400~\% med panorering, nypning och markörförankring, omedelbar övergång till originalkvalitet vid zoomning och Skift+pil-hopp på 10 fotografier.\\
EXIF och GPS & Uttag av kamera, objektiv, exponering, datum och GPS, global standard med arv, tvingad aktivering eller tvingad avstängning per galleri, bildmarkörer, gallerikartdata, datumförslag, dubblettunderlag och integritetskänslig döljning offentligt.\\
Röstning & Återkallelig positiv röst per bild med poängsynkronisering, anonym besökarhash eller inloggad tillhörighet, frekvensbegränsning för nya gilla-markeringar och aktivering eller avstängning per galleri.\\
Bildspelet & Valfri parjämförelse på grennivå med lämpliga offentliga bilder, beständigt par- och vinnartillstånd, resultatlista och sammanställning samt schemasäker vägran när lagringen är okänd.\\
Bildhanteraren & Administrationsöverlägg i den offentliga vyn för inloggade administratörer med blandat urval av fotografier och fysiska undergallerier, återställningsbar radering, skapande av fysiskt galleri under valt övergalleri med kopiering av underträd och återställning vid fel, flyttning, kopiering, dragning och delning endast för fotografier, inbyggd Web Share för flera filer där det stöds och ZIP-reserv eller hämtning för valda fotografier.\\
Hämtningar & Behörighetskontrollerade galleri-ZIP-filer, smarta galleri-ZIP-filer, ZIP-filer för valda fotografier, administratörs-ZIP för alla gallerier, normaliserade arkivsökvägar, innehållssignaturer för cache, begränsade originalantal och byteantal för smarta gallerier samt säker cacherensning.\\
Webbläsaruppladdningar & Uppladdningsflöden till befintliga eller nya gallerier, uttryckligt val av webbläsar- eller serverbearbetning, fullständig validering av förberedda miniatyrbilder, begränsade arbetare och omgångar med enskilda överstora paket under PHP:s hårda gräns, formatpolicyanvisningar, lokal ZIP-import och automatiskt filnamnsbyte efter skanning.\\
Filsystemsupptäckt & Sessionsbaserad begränsad mappskanning; hoppar över symboliska länkar, dolda och genererade mappar; importerar bildinnehållande hierarki på plats; kan flytta ohanterade fotografier till ett befintligt galleri eller radera säkra ohanterade kandidatträd; rapporterar sidofiler med enbart metadata separat.\\
Metadata-sidofiler i filsystemet & Valfri \filepath{gallery.json} per galleri, skriven från redigerbart galleritillstånd och läst vid upptäckt eller övergallerireparation för stödda metadata; användbar för flyttbarhet i filsystemet men ingen ersättning för databassäkerhetskopiering.\\
Uppladdnings-API & Galleribegränsade API-nycklar som visas en gång, central API-hanterare, återkallning av nycklar, gemensam inläsningspipeline och autentiserad automatiseringsadress.\\
Kompletterande Windows-program & Uppladdningar från bevakad mapp, manuell massuppladdning, valfri lokal generering av responsiva miniatyrbilder, parallellt och överlappande pipelinearbete, återförsöks- och tillståndsfiler, meddelandefält, valfria SimConnect-kamerakoordinater, originalradering efter bekräftad framgång och valfri lokal AI-metadataarbetare.\\
Mobil WebDAV & Galleribegränsade lösenordsuppgifter som visas en gång, Basic-autentisering, WebDAV-upptäckts- och skrivmetoder som mobilöverföringsklienter behöver, JPG/PNG/GIF/WebP-inläsning, spårning av senaste användning och återkallning.\\
Mediefilnamnsbyte & Provkörningsplan för fysiska filnamn, standarden \codeword{\{gallery_path\}_\{seq4\}} och dokumenterade jokerteckenplatshållare, kollisions- och säkerhetshantering, galleriets ordningssammanhang, samordnade uppdateringar av databas, offentliga sökvägar, derivat och cache, begränsad webbplatsövergripande tillämpning och integrerat automatiskt filnamnsbyte efter uppladdning.\\
Metadataorganiserare & EXIF-datumgruppering enbart som förhandsvisning med filter för tidigaste och senaste datum, ISO-undermappar enligt \codeword{YYYY-MM-DD}, återanvändning av matchande direkta datumundergallerier, bekräftade fysiska flyttar med derivat och begränsade tillämpningsomgångar; platsgruppering är fortfarande planerad och inte aktiv.\\
Dubblettdetektor för fotografier & Exakta SHA-256-par och metadataunderbyggda möjliga par, sammanhangslänkar, raderingsflöde, register per administratör för granskade par och exakta gallerier samt beständigt granskningstillstånd.\\
OpenAI-texthjälp & Krypterade API-nyckel- och modellinställningar per konto, valfri uttrycklig samtyckesväxel för miniatyrbilder, galleri- och bildbeskrivningsutkast, textförbättring, sammanhangssammanfattning, beskrivning av synligt innehåll från små genererade miniatyrbilder, översättningsförslag och infogning endast för manuell granskning.\\
Interna AI-metadata & Separata sökbara maskinmetadata och kötabeller, autentiserade jobbtilldelningar och resultat för Windows-arbetare, skrivskyddad granskning i bildredigeraren, underlag för smarta gallerier och sökning, spårning av modellgeneration och tvingad ombearbetning av gren.\\
Flygkartor och SimBrief & Manuell rutttext, uttrycklig koordinatsyntax, hämtning av SimBriefs senaste OFP, redigerbar genererad Markdown, sparade OFP-data och koordinater samt offentlig rendering endast från lagrade punkter.\\
Navigeringsdata & Lokal databasimport från OurAirports, medföljande offline-CSV, cachelagrade leverantörspunkter, valfri Navigraph OAuth-, paket- och AIRAC-integration, tolkning med offline först och uppslagstest i administrationen.\\
Gallerimigrering & API-överföring med källstyrd sändning eller målstyrd hämtning mellan samma versioner, versionshanterade manifest, metadata- och resurskontrollsummor, begränsade begäranden för en resurs, beständigt måltillstånd, återanslutning och återupptagning samt separat slutförande.\\
Tema och layout & Direkta palettfärger, typsnittsläge, rundade hörn, sidbredd, galleri- och kortrutnät, paginering, kortlayout, antalmarkering, offentlig miniatyrbildsrenderare, ljuslådestandard, GPS-markörpresentation, tre favorit- eller sidhuvudsgenvägar och direktförhandsvisning. Aktuellt temagränssnitt har ingen separat växel för ljust eller mörkt läge.\\
Profilering och webbläsaridentitet & Webbplatsens sidhuvudsbanderoll, avdelarmått och sträckning, optimerad bakgrund och opacitet, åsidosättningar per galleri, offentliga omslagsresurser, kvadratiskt beskuren favicon med genererade varianter på 32/48/180~pixlar samt egna CSS-förval och redigerare.\\
Lokalisering & Engelska som kanonisk reserv samt underhållna valbara gränssnittspaket på tjeckiska, tyska och svenska, webbläsarval per besökare, utformningar med flaggor och språkens egna namn, gemensamma webbläsar- och serverkataloger, paketdiagnostik, redigerare, import och export samt underhållet språkinnehåll per galleri och fotografi.\\
Funktionsväxlar & Globala ruttmedvetna växlar för offentlig sökning, ljuslådelägen, Bildhanteraren, bildröstning, Bildspelet, hämtningar, besökarkonton och privata samlingar, gallerikartor, flygkartor, navigeringsdata, SimBrief, OpenAI, lokala AI-metadata, uppladdnings-API, mobil WebDAV, gallerimigrering, Mediefilnamnsbyte och telemetri.\\
Centrala Inställningar & Vägledd stegvis installationsguide med slutligt godkännande, direktförhandsvisning och säkra ansvarsadaptrar; sökbart kategoriserat register, tangentbordsstyrd Spotlight-navigering, kanonisk delegering av sparande för godkända globala värden, djuplänkar till specialsidor, arvsmedvetet ansvar, standard- och källetiketter samt döljning av hemligheter.\\
Administrationens galleriarbetsyta & Synlighetsfilter med avmarkering av dolda rader, beständigt hopfällningstillstånd per gren, Fäll ihop/Fäll ut alla, hierarkiomordning, offentliga omordningsverktyg för synlig sida, sidopanelsredigerare och sammanhangsberoende kontroller för redigering, borttagning och synlighet på offentliga sidor för inloggade administratörer.\\
Administratörsautentisering & Inloggning med användarnamn eller återställningsadress, CSRF- och sessionsskydd, hashade subjekt för frekvensbegränsning, valfri beständig inloggningstoken, lösenordsåterställning via PHP mail/SMTP, konto- och profiländringar skyddade med kontroll av aktuellt lösenord samt valfri Google-inloggning som först måste kopplas från en autentiserad lösenordssession och kräver aktuellt lösenord för bortkoppling.\\
Besökarkonton, livscykel, favoriter, privata samlingar och olistad delning & Direkt skapande och radering av besökare av administratör samt registrering endast via inbjudan, administrationskontroller för avstängning, återställning och utloggning överallt med central återkallning av besökarbehörighet, tvingat byte av administratörstilldelade tillfälliga lösenord före vanlig besökarbehörighet, skannersäker e-postverifiering och aktivering, separat besökarinloggning och utloggning, särskild roterande kom-ihåg-uppgift, generisk skannersäker lösenordsåterställning, nylig lösenordsåterautentisering för känsliga kontoåtgärder, lösenords- och e-postbyte, egen kontoradering, privata svar utan lagring för konto, favoriter, samlingar och livscykel, favoritkontroller på behöriga kort och ljuslådebilder, besökarägda ordnade samlingar med aktuell källbehörighet, en återkallelig 30-dagars olistad skrivskyddad samlingslänk med hemlighet som visas en gång och ren omdirigering, skärning med källans ACL i mottagarens sammanhang utan administratörsundantag och strikt åtskillnad från galleribehörighet; ingen öppen registrering, offentlig samlingsupptäckt eller offentliga profiler, mottagarsamarbete, kommentarer eller uppladdningar.\\
Databasberedskap och Systemstatus & Schemaundersökning med tre tillstånd som skiljer Tillgängligt, Saknas, Okänt samt Avstängt där det är tillämpligt; separata funktionspaneler för säkerhet och autentisering, ändringskänslighet samt valfri presentation och rapportering; skrivningar nekas av säkerhetsskäl vid okänt obligatoriskt schema.\\
Databasunderhåll & Uttrycklig inventering av schema, migreringar och kodreferenser, lagringsinformation, provkörning av rensningskandidater, begränsad logisk rensning, särskild beredskap för äldre reparation, vald ANALYZE, optimeringsförhandsvisning för valda tabeller och separat bekräftad OPTIMIZE TABLE.\\
Lagringsstatistik & Skanning av original och genererade medier på begäran, filtyps- och storleksfördelning, största gallerier och filer, redovisning av DNG-master och miniatyrbilder, hantering av inaktuell ögonblicksbild, separat vy för databas- och tabellanvändning, den kombinerade åtgärden \ui{Uppdatera allt} för filer, databasuppskattningar och skrivskyddad inspektion samt begränsat förlopp.\\
Fullständig rapport & Fristående HTML-export för lagring, databas, galleri-, bild-, EXIF- och GPS-statistik, funktions- och inställningstillstånd, telemetri, loggar och körningsdiagnostik, med begränsad behandling av bilder och genererade medier, återförsök vid tillfälliga begärandefel och begränsade sammanfattningar med många olika värden.\\
Telemetri & Valfria integritetskontrollerade anonyma mätningar av användning, prestanda, cache och databas, respekt för DNT, uteslutning av administratörer, begränsad visningstid, lagring av rådata, timdata och dygnsdata, sammanställning och rensning, administrationsöversikt och HTML-export.\\
Administrationsloggar & Strukturerade händelser, filter och status, gruppering av upprepade händelser med utökning till medlemmar, begränsade exporter och ZIP-export, konfigurerad aktiv lagringstid, skyddade dagliga JSON+HTML-arkiv, återupptagnings- och återställningstillstånd samt verktyg för hämtning och visning av arkiv.\\
Underhållscenter & Beständigt Analysera--Granska--Utför--Verifiera-flöde, grupperad uppgiftsgranskning med en infälld grupp för valfria uppgifter utan upptäckt arbete, fortsatt synlig djup medieverifiering och fullständig optimering, servergodkänd uppgiftsplan, begränsade kontrollpunkter, paus, återupptagning och avbrott, central ändringsreservation, fysiska databassteg med en tabell och verifierad före- och efterrapport.\\
Schemalagt webbplatsunderhåll & Dagligt UTC-fönster, valfri begärandeutlöst fortsättning, begränsad bildomgång och tidsbudget, CLI- och webb-cron-stöd, åtgärd för att köra en del, återställning av avbrutet tillstånd, roterbar dold webb-cron-token, uppgifter för miniatyrbilder, metadata, rensning, telemetri och loggarkivering samt avstående när Underhållscentret äger den centrala ändringsreservationen.\\
Körningsdiagnostik & PHP-, webbhotells- och databasberedskap, GD/Imagick/ImageMagick- och formatstöd, DNG-policy, kopierbar supportsammanfattning, utvecklingsdiagnostik endast för administratörer för visarens förladdning, cache, minne, nätverk och bildrutetider samt prestandatestingångar endast för utveckling.\\
Integritet & Verifiering av installerade kärnmanifesthashar för hanterade programfiler samt manifestgenerering och validering vid utgivning; åtskilt från säkerhetskopiering av användardata.\\
Galleriprestandatest & Autentiserat webbläsarprestandatest endast för utveckling som använder anonyma förhandsvisningsinläsningar, dolda ramar med samma ursprung, serverrenderingsräknare och navigeringstider och sedan exporterar JSON för före- och efterjämförelser.\\
Uppdaterare & Upptäckt av stabila och beta-utgåvor, cachelagrade patchanteckningar, återupptagningsbara beständiga jobb, återupptagen hämtning med Range, arkiv- och manifestverifiering, förberedelse, återställningsögonblicksbild, aktivering, migreringar, rensning, lås och återställning av gamla lås, återförsök och avbrott före aktivering, filåterställning efter aktivering, återställning av stabil utgåva, ren ominstallation eller beta-installation, CLI-fortsättning och nödfilen \filepath{reset.php}.\\
Driftsättningshjälpmedel & Windows- och Linux-driftsättningsskript, utgåvemetadata, kärnmanifestgenerering och paketvalidering, med driftsättning utformad så att en lokalt installerad PHP-körbar fil inte blir ett ovillkorligt körningskrav för själva galleriet.\\
\bottomrule
\end{longtable}
\subsection{Programuppdateringar}

Den kompakta versionsöversikten samlar installerad och tillgänglig version, kanal, status och huvudåtgärder; detaljer om kodförrådet och GitHub API finns kvar i ett hopfällt diagnostikavsnitt. När ett jobb slutförs uppdaterar webbläsaren översikten direkt från lokala metadata, håller panelen öppen och URL:en oförändrad och döljer en inaktuell åtgärd för stabil uppdatering. Detta gör ingen ny GitHub-kontroll. Om läsningen misslyckas kan knappen för statusuppdatering försöka läsa lokalt igen utan att upprepa installationen.

Uppdateraren kontrollerar konfigurerade kanaler, validerar paket, förbereder innehåll, verifierar nödvändiga körningsfiler och storlekar, bevarar överskrivna filer i återställningslagring och aktiverar uppdateringen. Hantering av frekvensbegränsning och återförsökstillstånd förhindrar upprepade fjärrbegäranden under leverantörens väntetider. Ett återupptagningsbart jobbkort visas både under \ui{Status} och \ui{Avancerade verktyg}, så att en beta-installation, återställning eller ominstallation fortsätter visa sitt förlopp där den startades. Sidan kan stängas och öppnas igen utan att färdiga kontrollpunkter kastas.

Administrationens uppdateringsområde visar också cachelagrad utgåveinformation och versionsbundna patchanteckningar för tillgänglig kanal, så att ägaren kan granska en utgåva innan jobbet startas. Stabila kanaler och betakanaler förblir separata val; etiketten Beta hör till utgåvevalet i administrationsgränssnittet och ändrar inte uppdaterarens terminologi för beständiga jobb eller manifest.

Tillfälliga hämtnings- eller webbhotellsfel kan försöks igen. Ett paket vars filer inte stämmer med dess eget integritetsmanifest är annorlunda: att hämta samma publicerade programbygge igen kan inte ändra dessa byte. PHP Gallery ber därför administratören att avbryta det förberedda jobbet och välja en nyare utgåva eller betakod i stället för att erbjuda en oändlig återförsöksslinga.

Upptäckt av stabila utgåvor och paketarbete är separata begränsade åtgärder. En upptäcktsbegäran kan skapa ett beständigt bakgrundsjobb, medan en senare säker begäran eller CLI-köraren driver en kort del genom hämtning, arkivvalidering, uppackning, manifestverifiering, förberedelse, säkerhetskopiering för återställning, aktivering, migreringar och rensning. Range-begäranden kan återuppta avbrutna hämtningar, arbetarlås förhindrar samtidiga framsteg och ett gammalt lås kan återställas när dess ägare löpt ut. På inaktiva delade webbhotell schemalägger du \codeword{php scripts/application_update.php} från cron, t.ex. var femte minut. Varje anrop gör begränsad upptäckt eller driver befintligt bakgrundsarbete framåt; korrektheten beror inte på \codeword{set_time_limit()} eller \codeword{ignore_user_abort()}.

Aktuella utgåvor stämmer också av den lilla uppsättningen programägda serverpolicyfiler under uppdateringsförberedelse och aktivering. Det skyddar installationer vars äldre uppdaterare hoppade över en policyfil: den validerade utgåvekopian jämförs och publiceras atomärt, återställningsmetadata utökas före ersättningen och en migreringsstödd uppstartsväg kan reparera ett redan aktiverat äldre träd. Avstämningen är begränsad till kända policysökvägar och ändrar inte gallerimedier eller användarkonfiguration.

Före aktivering kan administratören avbryta eller försöka igen med ett återställningsbart jobb. Efter aktivering kan administrationens återställningsåtgärd återställa hanterade programfiler från den sparade återställningsögonblicksbilden, men filåterställning gör inte databasmigreringar ogjorda. Återställning av stabil utgåva, ren ominstallation och beta-installation förblir återupptagningsbara jobb med egen paketvalidering. Den fristående nödvägen \filepath{reset.php} är ett separat återställningsverktyg för att återgå till den stabila grenen när den vanliga administrationsrutten inte kan användas; den kräver ändå samma omsorgsfulla säkerhetskopiering och integritetsgranskning.

Under den korta kritiska aktiveringsfasen förhindrar en tidig beroendefri körningsspärr att nya vanliga begäranden läser in ett delvis ersatt program. Sådana begäranden får ett privat svar \codeword{503 Service Unavailable} som inte kan cachelagras tills aktiveringens slutförande har registrerats beständigt. Den autentiserade vägen för uppdateringsåterställning och status förblir tillgänglig, så att en avbruten aktivering kan fortsätta. Skadat eller ofullständigt spärrtillstånd nekar åtkomst av säkerhetsskäl i stället för att framställa installationen som välfungerande.

Samma tidiga körningsgräns omvandlar ohanterade och fångstbara fatala PHP-fel till avgränsade \codeword{500}-svar när PHP fortfarande kan styra svaret. Produktionsmiljöer måste hålla \codeword{display\_errors} avstängt; annars kan PHP-miljön själv skriva ut varnings- eller fatalfeldetaljer innan programnivåns hantering kan formatera ett säkert svar.

\important{Gör en samordnad säkerhetskopia före en programuppdatering. Behåll säkerhetskopian tills körningskontroller, migreringar, offentliga gallerier, skyddad åtkomst, uppladdningar och administrationsåtgärder har verifierats.}

\subsection{Integritetsmanifest}
\index{k00144@kärnmanifest}

\filepath{app/core-manifest.json} registrerar programversionen och hashvärden för integritetshanterade kärnfiler. Den lokala generatorn under \filepath{scripts/generate_manifest.php} beräknar slutliga hashvärden efter körningsändringar. Inkludering av dokumentation och användardata följer generatorns etablerade policy.
\section{Valfri bakgrund: så fungerar PHP Gallery}
\label{sec:architecture}
\index{a00016@arkitektur}
\index{k00137@kontroller}
\index{t00249@tjänster}
\index{v00264@vyer}

Följande sidor beskriver den begärandeväg som programmet använder. De är främst relevanta för underhåll och utveckling; vanlig galleriadministration kräver dem inte.

\subsection{Uppstart, dirigering och fördelning}
\index{u00257@uppstart}
\index{d00061@dirigering}
\index{k00138@kontrollertilldelning}
\index{b00021@begäran!l00151@livscykel}

Vanliga begäranden kommer in genom \filepath{public/index.php}; \filepath{index.php} i kodförrådets rot är en stödd kompatibilitetsingång för andra webbhotellsupplägg. Den offentliga ingången läser in \filepath{app/early\_runtime.php} före \filepath{app/bootstrap.php}, och sedan når båda uppläggen samma vanliga begärandepipeline.

Uppstarten läser lokal konfiguration och etablerar begärans sammanhang: databasåtkomst där det krävs, språk, sessionstillstånd, säkerhetsförutsättningar och gemensamma körningstjänster. Dirigeringen översätter sedan inkommande URL eller frågesträngsreserv till en intern rutt med parametrar såsom gallerisökväg, sidnummer, mediestorlek eller delningstoken. Fördelningen anropar kontrollern som är registrerad för rutten.

Ordningen är avsiktlig. En begäran måste klassificeras innan koden avgör hur den ska besvaras, och åtkomstkontroller måste köras innan skyddad utmatning sammanställs. En bild- eller miniatyrbilds-URL som ser direkt ut är därför ändå en programbegäran; mediekontrollern behörighetskontrollerar den innan byte returneras.

\begin{lstlisting}
browser request
  -> index.php or public/index.php
  -> app/bootstrap.php
  -> cms_run()
  -> cms_route_from_request()
  -> controller function
  -> service functions
  -> HTML, JSON, media, archive, or redirect response
\end{lstlisting}

\subsection{Kontroller}
\index{k00137@kontroller}

Kontroller ansvarar för samordning på HTTP-nivå. De får en bestämd rutt och ett begärandesammanhang, validerar åtgärden, anropar tjänsterna som behövs och väljer svarstyp. En offentlig gallerirutt returnerar normalt HTML; en åtgärd i administrationens sidopanel kan returnera JSON eller ett HTML-fragment; en medierutt returnerar en behörighetskontrollerad fil; en hämtningsrutt kan strömma ett arkiv.

Offentlig gallerirendering börjar i \filepath{app/controllers/public_gallery.php}, med relaterade ansvarsområden uppdelade i avgränsade kontrollerfiler. Uppladdningar, underhåll, uppdateringar, hämtningar, medieåtkomst och administrationsfunktioner har egna ingångar, medan gemensamma domänregler finns kvar i tjänster.

Kontroller är ingen genväg förbi behörighet. Rutter för miniatyrbilder, kartor, ljuslåda, sökning, hämtning och direkta medier upprepar relevant åtkomstbeslut i stället för att förutsätta att besökaren tidigare lyckats öppna gallerisidan.

\subsection{Tjänster}
\index{t00249@tjänster}

Återanvändbara programregler finns i \filepath{app/services/}. Tjänster hittar gallerier och bilder, bedömer åtkomst, normaliserar inställningar, hanterar uppladdningar och ändringar, förbereder miniatyrbilds- och mediemanifest, söker, underhåller metadata och samordnar annat domänarbete. En kontroller ska kunna begära en sådan åtgärd utan att implementera reglerna igen.

Åtskillnaden är viktig eftersom samma fotografi kan nås från flera rutter. Åtkomstpolicy, miniatyrbildsval och raderingsbetydelse ska inte bero på om begäran kom från ett gallerikort, sökresultat, ljuslåda, administrationsverktyg eller direktlänk. Originalfotografier finns kvar i galleriträdet; tjänster samordnar databasposter och genererade derivat som beskriver eller stöder dem.

\subsection{Vyer och webbläsarmoduler}
\index{v00264@vyer}
\index{w00267@webbläsarmoduler}

Vyer renderar förberedda data till HTML. De ansvarar för struktur och presentation, inte för att avgöra om besökaren är behörig eller hur en miniatyrbild genereras. Serverutmatningen innehåller semantiska länkar, formulär, titlar, alternativtext och inledande mediekandidater innan JavaScript-utökningen börjar.

Webbläsarmoduler lägger till ljuslådebeteende, sökförslag, röstning, uppladdningsförlopp, dra-och-släpp-åtgärder, miniatyrbildsuppgraderingar, diagnostik och administrationens sidopanel. Koden förblir utan ramverk. Offentliga sidor behåller därför en användbar serverrenderad grund, medan autentiserade arbetsflöden kan ge rikare interaktioner på plats.

Administrationens sidopanel är dynamisk: dess fragment kan ersättas utan att lämna sidan. Moduler som ansvarar för panelkontroller använder därför delegerad eller livscykelmedveten händelsebindning så att en nyligen renderad kopia fortsätter fungera och ett gammalt fragment inte längre äger lyssnare.

\subsection{Offentlig renderingspipeline för valt galleri}
\index{r00202@renderingspipeline}

En begäran för valt galleri hittar först galleriet och dess effektiva åtkomstpolicy, inklusive ärvt skydd, lösenords- och delningstillstånd, publiceringstillstånd och åldersbegränsningar. Endast tillåtet innehåll går vidare till rendering.

Kontrollern läser sedan in direkta undergallerier och aktuell sida med fotografier tillsammans med metadata som behövs för titlar, beskrivningar, taggar, datum, röster, ordning och navigering. Miniatyrbildsmetadata förbereds i omgångar, och begärandelokala mediemanifest beskriver behöriga kandidater som är tillgängliga för varje synlig bild.

Vyn tar emot den förberedda modellen och renderar sidtitel, profilering, navigering, brödsmulor, kort, paginering, ljuslådekonfiguration och sidfotsresurser. Webbläsarmoduler utökar resultatet efter leveransen. Åtkomst förblir serverns ansvar; JavaScript kan göra en interaktion snabbare eller bekvämare men gör inte en otillåten kandidat tillåten.

Denna uppdelning med servern först bevarar också användbart beteende när skript blockeras eller kommer sent. Direktlänkar, semantiskt innehåll och de första bildvalen finns redan i HTML-koden, medan JavaScript hanterar funktioner som behöver en interaktiv klient.

\section{Valfri bakgrund: vad databasen kommer ihåg}
\label{sec:database}
\index{d00052@databasschema}
\index{m00163@migreringar}

Databasen lagrar programtillstånd som inte hör hemma i bildfilnamn: beskrivningar, integritet, ordning, taggar, röster, datum, offentliga adresser, konton och underhållshistorik. Översikten nedan räcker för planering av säkerhetskopiering och diskussioner med webbhotell; definitioner på kolumnnivå finns kvar i databasreferensen \cite{database}.

\subsection{Huvudsakliga domäner}

\begin{longtable}{>{\bfseries}>{\hspace{0pt}}p{0.25\linewidth}>{\hspace{0pt}}p{0.67\linewidth}}
\toprule
Domän & Beständigt ansvar\\
\midrule
\endhead
Användare och autentisering & Administratörsidentitet, lösenordshashar, återställningsadress, externa identitetskopplingar, återställningstoken och avgränsade inloggningsuppgifter.\\
Gallerier & Mappsökväg, hierarki, offentlig sökväg, metadata, synlighet, åtkomst, profilering, presentationsåsidosättningar och driftordning.\\
Bilder & Galleritillhörighet, relativ sökväg, filnamn, metadata, mått, synlighet, kontrollsumma, EXIF-data, ordning och härlett tillstånd.\\
Taggar & Återanvändbar taggidentitet samt galleri- och bildrelationer.\\
Röster & Galleri- och bildpoängposter med besökarkoppling.\\
Miniatyrbilder & Förteckning över giltiga genererade varianter, mått, format, status och derivatversion.\\
Inställningar & Skalära värden för program, tema, funktioner, uppdatering, underhåll och rendering.\\
Granskning och telemetri & Administratörshändelser, driftdiagnostik, integritetskontrollerade mätningar och sammanställningar.\\
Dubblettregister & Kanoniska parbeslut per administratör och regler för döljning av exakta gallerier.\\
Jobb och token & Begränsat tillstånd för webbläsare, detektor, migrering, WebDAV, återställning, delning och underhåll där beständig lagring krävs.\\
\bottomrule
\end{longtable}

\subsection{Relationer och kaskader}

Främmande nycklar uttrycker tillhörighet där schemautveckling och kompatibilitet tillåter det. Kaskader tar bort beroende arbetsflödesposter när deras överordnade identitet försvinner. Innehållsradering går fortfarande genom programtjänster så att filsystems- och databaseffekter förblir samordnade. Unikhetsvillkor uttrycker identiteter såsom kanoniska dubblettpar, exakta administratörs- och galleriregler, URL-namn, token och metadatavarianter.

\subsection{Programinställningar}

Tabellen \codeword{app_settings} lagrar nyckel-värde-konfiguration. \codeword{app_setting()} läser värden, \codeword{set_app_setting()} lagrar normaliserade värden och avgränsade tjänster tolkar domänspecifika inställningar. Offentlig språkkonfiguration använder \setting{public_language} för webbplatsstandarden, \setting{public_language_selector_enabled} för besökaråsidosättningsfunktionen som är~aktiverad som standard och \setting{public_language_selector_languages} för dess ordnade icke-tomma språklista. Den offentliga miniatyrbildsrenderaren använder \setting{public_thumbnail_rendering_mode} med värdena \setting{responsive} och \setting{progressive}.

\section{Valfri referens för utvecklare}
\label{sec:development}
\index{u00262@utveckling}
\index{k00134@kodstil}

Denna referens är för personer som ändrar själva PHP Gallery. Ägare som bara behöver kontroller för publicering och underhåll kan fortsätta till avsnitt~\ref{sec:testing}.

\subsection{Kodförrådets organisation}

\begin{longtable}{>{\ttfamily}>{\hspace{0pt}}p{0.27\linewidth}>{\hspace{0pt}}p{0.65\linewidth}}
\toprule
\textrm{Plats} & \textrm{Ansvar}\\
\midrule
\endhead
app/controllers/ & HTTP-kontroller och svarssamordning.\\
app/services/ & Återanvändbara domän- och infrastrukturtjänster.\\
app/views/ & Layout och återanvändbara serverrenderade vyhjälpmedel.\\
app/lang/ & Översättningskataloger för server och webbläsare.\\
database/migrations/ & Tidsstämplad schemautveckling.\\
public/assets/ & JavaScript utan ramverk, CSS och offentliga resurser.\\
scripts/ & Verktyg för migrering, integritet, driftsättning och underhåll.\\
tests/ & Fristående PHP-tester och riktade JavaScript-tester.\\
winapp/ & Windows-specifika verktyg och integrationer.\\
\bottomrule
\end{longtable}

\subsection{Kodkonventioner}

Nya PHP-filer använder strikta typer och indrag med fyra blanksteg. Kontroller samordnar begärandebeteende, medan tjänster ansvarar för återanvändbar logik. Ny JavaScript och CSS förblir utan ramverk och hör hemma bland offentliga resurser. Migreringar använder filnamn med tidsstämpelprefix. Befintliga förklarande kommentarer, dokumentationssträngar och PHPDoc bevaras och rättas när beteendet utvecklas.

Dokumentationssträngar krävs för namngivna funktioner, metoder och klasser eller liknande deklarationer, inklusive gränssnitt, traits och uppräkningstyper. Placera dokumentationen omedelbart ovanför deklarationen. Namngivna funktioner och metoder behöver en saklig beskrivning av syftet, en typad och beskriven post för varje uttrycklig parameter samt en typad returpost, inklusive void. Klasser behöver en saklig beskrivning av syftet. Anonyma funktioner, closures, pilfunktioner, callbacks, variabler, konstanter och klassegenskaper kräver inga dokumentationssträngar för deklarationen; namngivna deklarationer inuti callbacks kontrolleras fortfarande. Filhuvuden med författaruppgifter, inbyggda signaturtyper, kommentarer om driftpolicy och parsertäckning förblir separata krav.

Atomära moduler ska ansvara för ett sammanhängande område. Gemensamma kontrakt behöver normaliserade indata, uttryckliga utdata och riktade tester. Dynamiska administrationsfragment och offentliga fragment kräver livscykelsäker initiering och avveckling.

\subsection{Lägga till en funktion}

En typisk funktionsändring följer denna ordning:

\begin{enumerate}
  \item Spåra aktuell implementation och identifiera ansvarig kontroller, tjänster, vyer, webbläsarmoduler, lagring och tester.
  \item Definiera kontrakt för åtkomst, CSRF, omfattning, reservbeteende, migrering och användning utan JavaScript.
  \item Lägg till avgränsade tjänster och ruttintegration.
  \item Lägg till migreringsarbete enbart genom nya filer där beständiga schemaändringar krävs.
  \item Lägg till översatta gränssnittssträngar.
  \item Lägg till riktade automatiserade tester och dokumenterad manuell verifiering.
  \item Uppdatera arkitektur, kodkarta, databas, testning och användarvägledning enligt respektive ansvar.
  \item Generera om integritetsmetadata efter slutliga körningsändringar.
  \item Granska hela skillnaden och kör relevant lokal testsvit.
\end{enumerate}

\section{Kontrollera galleriet före publicering}
\label{sec:testing}
\index{t00246@testning}
\index{k00143@kvalitetssäkring}

För en galleriägare innebär testning att gå igenom webbplatsen som den avsedda målgruppen. En tekniskt välfungerande sida kan fortfarande innehålla en privat plats, en otydlig titel, ett felaktigt datum, en oavsiktlig hämtningsbehörighet eller en besvärlig mobillayout. En kort mänsklig granskning fångar dessa presentations- och integritetsproblem.

\subsection{Publiceringskontroll ur besökarens perspektiv}
\index{p00196@publiceringschecklista}

Öppna ett privat webbläsarfönster och följ samma länk som du tänker dela. Kontrollera:

\begin{enumerate}
  \item Titeln och första stycket förklarar samlingen.
  \item Valt omslag förblir tydligt i liten storlek.
  \item Undergallerier visas i en rimlig ordning.
  \item Bildtitlar och beskrivningar stämmer med synliga motiv.
  \item Privata fotografier och gallerier förblir dolda.
  \item Lösenords- eller delningsåtkomst fungerar från en ny session.
  \item GPS-kartor avslöjar bara avsedda platser.
  \item Första sidan läses in smidigt på en telefon.
  \item Den stora visaren, framåt- och bakåtkontrollerna och stängningen känns naturliga.
  \item Sökning, taggar, röstning och hämtningar fungerar enligt galleriets syfte.
\end{enumerate}

Be någon annan prova galleriet utan instruktioner. Personens första klick och första fråga avslöjar ofta navigering eller formuleringar som bara ägaren tyckte var självklara.

\subsection{Automatiserade kontroller för programunderhållare}

Utvecklare kan köra projektets körbara testskript efter ändringar i programkoden. Den samlade köraren omfattar projektsviten, medan riktade skript verifierar särskilda funktioner.

\begin{lstlisting}[language=bash]
php tests/run.php
php tests/public_thumbnail_rendering_model_test.php
php tests/public_thumbnail_markup_test.php
php tests/duplicate_photo_detector_test.php
php tests/duplicate_photo_ledger_test.php
php tests/migration_consistency_test.php
node tests/progressive_thumbnail_renderer_test.mjs
\end{lstlisting}

\subsection{Manuell webbläsarverifiering}

\subsubsection{Nätverks- och interaktionskontroller}

Webbläsarberoende programkontroller använder representativa data, anonyma och autentiserade sessioner, tom cache, mobila och stationära visningsområden, drift utan JavaScript där det är relevant, inställningar för minskad rörelse och nätverksbegränsning.

För offentlig miniatyrbildsrendering undersöker du panelerna Elements och Network. Responsivt läge ska exponera hela sin aktiva kandidatuppsättning. Progressivt läge ska exponera en liten aktiv kandidat och overksamma större kandidater före aktivering och sedan utföra högst det dokumenterade antalet samtidiga uppgraderingar för relevanta kort. Kontrollera layoutstabilitet, cacheåteranvändning, felreserv, ljuslådeinteraktion, kartor, röstning, skyddade medier och direktlänkar.

\subsection{Rutiner vid utgivning}

En utgåvegranskning bekräftar konsekvent körningsversion, ordnad migreringskompatibilitet, manifesthashar, dokumentation, översättningar, cacheuppdatering för webbläsare, tester, paketundantag och användardatasäkerhet. Underhållaren börjar från en ren utgåvegren, jämför den med exakt föregående utgåvetagg, granskar varje mellanliggande commit och sökväg och registrerar allt som avsiktligt utesluts. Körningsversionen i \filepath{app/bootstrap.php}, nyckeln och taggen \filepath{release-metadata.json}, senaste rubriken i \filepath{PATCH_NOTES.md}, dokumentationsmarkörer, avsett arkivnamn och slutlig Git-tagg måste alla stämma överens.

Nya migreringar granskas i tidsstämpelordning och provas genom tester för migreringskonsekvens och äldre körare. Underhållaren bygger om denna PDF med hela sekvensen \filepath{docs/LATEX_BUILD.md} och undersöker titelblad, utgåveöversikt, innehållsförteckning, bokmärken, index och ändrade funktionsavsnitt. Varje ändrad PHP-fil får \codeword{php -l}; ändrad JavaScript och Node-testdata får syntaxkontroller; riktade funktionstester, översättningskonsekvens, dokumentationstäckning, hela PHP-sviten och \codeword{git diff --check} måste godkännas.

Först efter den sista källkods- och dokumentationsändringen kör underhållaren \codeword{php scripts/generate_manifest.php} och dess läge \codeword{--check}. Driftsättningshjälpmedlet skapar sedan begärd lokal mapp eller ZIP-fil; det paketerar bara filer och kör inte PHP eller reparerar ett gammalt manifest. Undersök arkivförteckningen och bekräfta att den innehåller aktuell PDF, patchanteckningar, utgåvemetadata, migreringar och \filepath{app/core-manifest.json}, medan skyddad konfiguration, cache, loggar, tillfälliga filer, lokala verktyg, driftsättningsutdata och användarmedier följer paketeringspolicyn. Kontrollera manifestet igen efter paketeringen, skapa utgåvecommit och kommenterad \codeword{v_<version>}-tagg från det verifierade trädet och gör ett grundläggande funktionstest av en uppdateraruppgradering från föregående stabila utgåva. Saknade lokala beroenden måste rapporteras med exakta blockerade kommandon och miljödetaljer i stället för att kalla utgåvan fullt verifierad.

\section{Få galleriet att kännas snabbt}
\label{sec:performance}
\index{p00191@prestanda}
\index{h00120@hastighet}

Upplevd hastighet är inte ett enda tal. Besökare märker när sidan först visas, när de första fotografierna blir igenkännbara, om rullningen förblir responsiv, hur snabbt kontrollerna reagerar och hur lång tid en större bild tar efter att de begärt den. Anslutningskvalitet, serveravstånd, sidstorlek, miniatyrbildspolicy, profileringsresurser och vald renderare bidrar alla.

\subsection{Praktiska sätt att förbättra hastigheten}

\begin{itemize}
  \item Håll antalet kort per sida rimligt för målgruppen och webbhotellsavtalet.
  \item Generera konfigurerade miniatyrbildsstorlekar i förväg för stora offentliga samlingar.
  \item Undvik onödigt stora huvudpanels-, bakgrunds-, logotyp- och andra temaresurser.
  \item Föredra den responsiva miniatyrbildsrenderaren när förutsägbarhet och funktion utan JavaScript är viktigast; använd den progressiva renderaren när dess stegvisa skärpning passar samlingen och webbläsarfördelningen.
  \item Placera galleriet på ett webbhotell geografiskt nära huvudmålgruppen där det är praktiskt.
  \item Testa med tom cache. En varm utvecklarwebbläsare kan dölja kostsamma överföringar vid första besöket.
  \item Kontrollera PHP- och databasfördröjning separat från bildöverföringstid innan du antar att varje långsamt galleri har ett miniatyrbildsproblem.
\end{itemize}

\subsection{Test med långsam anslutning}

Öppna galleriet i ett privat webbläsarfönster, välj ett representativt mobilt visningsområde, aktivera nätverksbegränsning i utvecklarverktygen och ladda om med tom cache. Se i vilken ordning sidan, omslagen, synliga kort och senare bilduppgraderingar kommer. Upprepa med den renderare eller sidstorleksinställning du tänker jämföra; annars gör cachelagrade filer jämförelsen otillförlitlig.

\subsection{Tekniska mätningar för underhållare}

Offentlig renderingsprofilering och webbläsardiagnostik är användbara när en subjektiv rapport behöver siffror. Registrera minst:

\begin{itemize}
  \item tid till första byte och renderingstid på servern;
  \item HTML-överföringsstorlek;
  \item antal begäranden och överförda byte för inledande miniatyrbilder;
  \item tid tills det största synliga innehållet har stabiliserats;
  \item layoutförskjutningar medan kort och medier läses in;
  \item arbete i huvudtråden som orsakas av webbläsarmoduler;
  \item antal och storlek på progressiva kandidatuppgraderingar;
  \item totalt antal databasfrågor från utvecklingsprofileraren där den finns;
  \item bildavkodning och minnesbelastning vid test av mycket stora original eller djup ljuslådezoom.
\end{itemize}

Utvecklingsprofileraren och registrerade prestandatestposter är avsedda för före- och efterjämförelser mellan renderarlägen och gallerilayouter. Använd representativa skärmar med hög pixeltäthet och begränsade anslutningar i stället för att bara förlita dig på en snabb dator med varm cache.
\section{Felsökning}
\label{sec:troubleshooting}
\index{f00077@felsökning}

\subsection{Programmet startar inte}

Bekräfta PHP-version, nödvändiga tillägg, konfigurationsplats, databasanslutning, skrivbara sökvägar och webbrotsdirigering. Granska PHP- och webbserverloggar. Kör syntaxkontroller på nyligen ändrade PHP-filer och undersök väntande migreringar.

\subsection{Galleri eller bilder saknas}

Bekräfta filsystemssökvägar, galleriets upptäcktstillstånd, databasposter, synlighet, övergalleriets åtkomstregler, lösenordsupplåsningstillstånd, NSFW-policy, stödd filtyp och bildens relativa sökväg. Testa både som administratör och anonym besökare.

\subsection{Miniatyrbilder saknas eller är oskarpa}

Granska miniatyrbildsmetadata och underhållsrapporter. Bekräfta genererade storlekar, formatstöd, konfigurerade gränser, originalgeometri, derivatstatus och behörighetskontroll för offentlig adress. I responsivt läge undersöker du webbläsarens valda kandidat och \codeword{sizes}. I progressivt läge undersöker du den lilla kandidaten, kötillstånd, vald ersättning och avkodningsresultat.

\subsection{Administrationsåtgärder lämnar sidopanelen}

Bekräfta att matchande JavaScript-modul har lästs in, delegerade hanterare känner igen formuläret, AJAX-huvuden och svarskontrakt är korrekta, returnerat fragment innehåller förväntade livscykelmarkörer och generiska administrationsformulärhanterare utesluter funktionsägda formulär. Direkt POST-beteende kan fortfarande ge den dokumenterade reservrutten.

\subsection{Migrering misslyckas}

Behåll en verifierad databassäkerhetskopia och registrera migreringsfilnamn, Gallery-version, databasserverversion, avgränsat skärmmeddelande, begärandereferens och befintligt schema- och registertillstånd. Försök igen genom galleriadministrationssidan så att den vanliga idempotenta kompatibilitetsvägen kan slutföra en avbruten migrering. Version~\version{} behöver inte utlösare, lagrade rutiner, \codeword{SUPER}, \codeword{log\_bin\_trust\_function\_creators} eller en global databasändring. Om det avgränsade felet kvarstår undersöker du Körningsdiagnostik och ger begärandereferensen till underhållaren eller webbhotellet; klistra inte in databasuppgifter eller råa produktionsdata i en supportrapport.

\subsection{Kompilering av PDF-manualen misslyckas}

Kör kommandona i \filepath{docs/LATEX_BUILD.md}. Bekräfta att \codeword{pdflatex} och \codeword{makeindex} är installerade. En första MiKTeX-körning kan begära vanliga paket. Radera ignorerade mellanliggande byggfiler efter en misslyckad paketövergång och upprepa den dokumenterade sekvensen.

\section{Checklista för säkerhetskopiering och återställning}
\label{sec:backup}
\index{z00277@återställning}

\begin{enumerate}
  \item Exportera hela programdatabasen.
  \item Kopiera \filepath{galleries/} med originalfiler och galleriägda resurser.
  \item Kopiera lokal konfiguration och relevanta installationshemligheter säkert.
  \item Bevara egna temaresurser och egen CSS.
  \item Registrera programversion och väntande migreringstillstånd.
  \item Förvara säkerhetskopior utanför det aktiva webbträdet.
  \item Testa återställning regelbundet i en isolerad miljö.
  \item Verifiera offentliga rutter, skyddad åtkomst, administratörsinloggning, miniatyrbilder, uppladdningar och migreringar efter återställning.
\end{enumerate}

Säkerhetskopior skapade av uppdateraren bevarar överskrivna programfiler för uppdateringsåterställning. En fullständig driftsäkerhetskopia omfattar även databas, galleriträd, lokal konfiguration och installationsägda resurser.

\subsection{Validera en återställningsuppsättning}
\index{z00277@återställning!s00239@säkerställande}

Håll databasexport, originalfotografier, lokal konfiguration, relevant \filepath{data/}-lagring, papperskorgsinnehåll, egna resurser och exakt programversion tillsammans som en samordnad återställningsuppsättning. En leverantörsögonblicksbild eller ett begränsat underhållsfönster kan förhindra att skrivningar splittrar databas- och filtillstånd. Förvara en kopia utanför webbhotellet och kom överens om hur mycket nytt arbete som får gå förlorat och hur snabbt tjänsten måste återkomma.

CLI-ingången \filepath{scripts/recovery.php} kan inventera uppsättningen och validera återställda filprov och manifesthashar på en separat plats. Följ \filepath{docs/RECOVERY_ASSURANCE.md} för underlagsformat och kommandon. Verktyget kör inte en databasdump, återställer inte produktion och verifierar inte en säkerhetskopieringsleverantör. Dess syntetiska skadetest bevisar att kontrollverktyget upptäcker skador i testdata, inte att din verkliga säkerhetskopia kan återställa webbplatsen.

Innan en återställd installation öppnas isolerar du e-post, schemalagda jobb, integrationer och skrivåtkomst från produktionsmål. Verifiera återställd databas, relationer, administratörsinloggning, nekad åtkomst till privata medier och papperskorgsåterställning. Registrera säkerhetskopians ålder, förfluten återställningstid och operatörens observationer utan att kopiera hemligheter till rapporter. Använd \filepath{docs/RECOVERY_OFF_HOST.md} för hela förfarandet. Uppdaterarens filåterställning gör fortfarande inte databasmigreringar ogjorda.

\section{Underhållarverifiering och utgåveunderlag}
\label{sec:release-evidence}
\index{u00260@utgåvekvalificering}
\index{t00246@testning!t00247@tillfälliga arbetsflöden}

Automatiserad regression och driftsgodkännande besvarar olika frågor. Den centrala granskningen kontrollerar kod, kontrakt, syntax, tillgängliga webbläsartestfall och utgåvekonsekvens. Den isolerade arbetsflödestestmiljön lägger till riktiga databas-, HTTP- och Chromium-flöden med genererade fotografier och konton. Den kräver uttrycklig konfiguration enbart för en tillfällig miljö och får aldrig använda installationens databas eller medier. Följ \filepath{docs/GALLERY_WORKFLOWS.md} för förutsättningar; en saknad miljö är en lucka i täckningen, inte bevis för en fungerande funktion.

För utgåveförberedelse följer du \filepath{RELEASE.md}: slutför källkod och dokumentation, bygg om och undersök denna PDF, generera manifestet och initiera kvalificeringsposten. Kör en central utgåvegranskning, inte en sekvens med quick/full/release. Med de isolerade förutsättningarna konfigurerade accepterar \filepath{scripts/gallery_workflow_mysql.php} \codeword{--release} för att förbereda sin testmiljö kring samma utgåvegranskning.

Kvalificeringsverktyget \filepath{scripts/release_qualification.php} registrerar det exakta fingeravtrycket för källkod och PDF. Dess kommando \codeword{render} skapar begränsade sidförhandsvisningar; rendering ensam godkänner inte deras utseende. Kommandot \codeword{record-audit} accepterar endast den matchande centrala utgåverapporten, inte en äldre fullständig granskning. En ändrad källfil eller ombyggd PDF väljer nytt väntande underlag i stället för att tyst återanvända tidigare godkännande.

Granska titelblad, innehållsförteckning, index, ändrade sidor, interna länkar och representativa webbläsarflöden. Registrera vad som faktiskt kontrollerades, av vem och med vilket resultat. Låt olösta kontroller vara väntande och håll grundläggande funktionstest av uppdateraren efter publicering separat. En grön automatiserad rapport är inte tillstånd att publicera, och verktyget gör aldrig commit, taggar, laddar upp eller publicerar utgåvan. Se \filepath{docs/RELEASE_QUALIFICATION.md} för postkommandona.

\section{Ordlista}
\label{sec:glossary}
\index{o00181@ordlista}

\begin{description}
  \item[AJAX] Webbläsarbegäran som används för att uppdatera en del av sidan och samtidigt bevara det omgivande gränssnittet.
  \item[Grenomfattning] Ett galleri och dess undergallerier när vald funktion definierar rekursiv omfattning.
  \item[Kanoniskt par] Två bildidentifierare lagrade i stabil ordning från lägre till högre.
  \item[CSRF-token] Sessionsbundet värde som används för att validera avsiktliga formulärinlämningar som ändrar tillstånd.
  \item[Derivat] Genererad medierepresentation såsom en JPEG- eller WebP-miniatyrbild.
  \item[EXIF] Kamera- och tagningsmetadata inbäddade i stödda bildfiler.
  \item[Integritetsmanifest] Versionshanterad lista med kärnsökvägar och förväntade hashvärden.
  \item[Ljuslåda] Bildvisare i helskärm eller överlägg med navigering och relaterade kontroller.
  \item[Mediemanifest] Begärandelokala renderingsdata förberedda från miniatyrbildsmetadata för synliga bilder.
  \item[NSFW-skydd] Ålders- och galleri- eller bildpolicy som styr åtkomst till begränsade medier.
  \item[Progressiv skärpning] Miniatyrbildspipeline med små bilder först som senare aktiverar en avkodad större kandidat.
  \item[Responsivt val] Webbläsarens inbyggda kandidatval från en omedelbart aktiv källuppsättning.
  \item[Valt galleri] Galleri som representeras av aktuell offentlig rutt.
  \item[Sidopanel] Administrationens sammanhangsberoende gränssnitt till höger, som läses in och uppdateras genom fragmentbegäranden.
  \item[Miniatyrbildsgränser] Gallerimedvetna minsta och största genererade storlekar som kan väljas.
\end{description}

\clearpage
\section{Referenser}
\label{sec:references}

\begin{thebibliography}{99}

\bibitem{readme}
PHP Gallery-projektet, \emph{README.md}, lokal produktöversikt, funktionsguide, krav och installationsdokumentation, version \version.

\bibitem{architecture}
PHP Gallery-projektet, \emph{ARCHITECTURE.md}, lokal referens för programarkitektur och delsystems livscykler, version \version.

\bibitem{database}
PHP Gallery-projektet, \emph{DATABASE.md}, lokal referens för migreringar och beständigt schema, version \version.

\bibitem{codemap}
PHP Gallery-projektet, \emph{CODEMAP.md}, lokal karta över källkodsansvar och underhåll, version \version.

\bibitem{settingsinventory}
PHP Gallery-projektet, \emph{docs/ADMIN\_SETTINGS\_INVENTORY.md}, kanonisk förteckning över ansvar, reservbeteende, känslighet och migreringar för globala Inställningar, version \version.

\bibitem{testing}
PHP Gallery-projektet, \emph{TESTING.md}, lokal guide för automatiserad och manuell verifiering, version \version.

\bibitem{agents}
PHP Gallery-projektet, \emph{AGENTS.md}, lokala riktlinjer för bidrag till kodförrådet och säkerhet.

\bibitem{phpmanual}
The PHP Documentation Group, \emph{PHP-manualen}, \url{https://www.php.net/manual/en/}.

\bibitem{mysqlmanual}
Oracle Corporation, \emph{MySQL-referensmanualen}, \url{https://dev.mysql.com/doc/}.

\bibitem{mariadbmanual}
MariaDB Foundation, \emph{MariaDB Server-dokumentation}, \url{https://mariadb.com/docs/server/}.

\bibitem{htmlimages}
WHATWG, \emph{HTML:s levande standard: bilder}, \url{https://html.spec.whatwg.org/multipage/images.html}.

\bibitem{owasp}
OWASP Foundation, \emph{Säkerhetsvägledning för webbapplikationer}, \url{https://owasp.org/}.

\end{thebibliography}

\clearpage
\phantomsection
\addcontentsline{toc}{section}{Index}
\printindex

\end{document}
